\documentclass[12pt, ukrainian]{article}
\usepackage[a4paper]{geometry}
\setlength\parindent{0mm}
%\setlength\parskip{4mm}
\geometry{top=1.3cm,bottom=1.3cm,left=1.5cm,right=1.5cm}
\pagestyle{empty}
\usepackage[T2A]{fontenc}
\usepackage[utf8]{inputenc}
\usepackage[ukrainian]{babel}
\usepackage{verbatim}
\usepackage{fancyvrb}
\def\code#1{\Verb!#1!}
\begin{document}
%begin
{{\begin {scriptsize}\par
{ \center  Київський національний університет імені Тараса Шевченка \par}
{\parbox {5 cm}{Освітня програма \hfill}{\parbox {7 cm}{Інформатика\hfill}}\parbox {6 cm}{\hfill Семестр 1}}\par
{\parbox {5 cm} {Навчальний предмет \hfill}\parbox {7 cm}{Програмування \hfill}\parbox {6cm}{\hfill}\par}\vspace{-15pt} \end{scriptsize}}{\center Екзаменаційний білет \No \textbf { 1 }\par}}
{%beginitem
1. Які з наведених типів не є цілими (integral), якщо такі є серед перелічених?

int*, signed char, unsigned long, string?
%enditem
\par
\vspace{2mm}
%beginitem
2. Скільки змінних, що є вказівниками, тут оголошено?
char *s1, pc, **s2;?
%enditem
\par
\vspace{2mm}
%beginitem
3. За наявних оголошень \code{void f(int\&); int n;} які з виразів \code{f(13), f(++n), f(\&n)} є синтаксично коректними?
%enditem
\par
\vspace{2mm}
%beginitem
4. Що таке просторова складність алгоритму?
%enditem
\par
\vspace{2mm}
%beginitem
5. Яку часову оцінку складності має алгоритм швидкого сортування?
%enditem
\par
\vspace{2mm}
%beginitem
6. Написати програму, що запитує у користувача значення дійсних параметрів $x$ та $y$, обчислює для них значення виразу та повiдомляє введенi данi та результат обчислення.
$$\frac{\pi x + 70}{(\tg x + 0.3) (y - 93)}$$ 

%enditem
\par
%beginitem
7. Написати функцію, що знаходить найменше число з послідовності цілих (задається масивом), у десятковому запису якого нема цифри 5. Повертає індекс останнього входження цього числа або $-1$ (за відсутності). Використовувати додаткові структури даних (у тому числі рядки) заборонено.

По завданням 2 та 3 оцінюється правильність та якість коду, а також економність обчислень.

%enditem
}
{\smallskip\smallskip \begin {scriptsize} Затверджено на засіданні кафедри теоретичної кібернетики, протокол \No 3  від   08.11.2025 р.\par\smallskip\smallskip\smallskip\smallskip
{\parbox {6.5 cm} {Зав. кафедри \dotfill Ю.В.~Крак}}\hspace {0.7cm} {\hbox to 6.5 cm { Екзаменатор \dotfill Т.О.~Карнаух}} \end{scriptsize}}
%end
\vspace{0.9cm}
%begin
{{\begin {scriptsize}\par
{ \center  Київський національний університет імені Тараса Шевченка \par}
{\parbox {5 cm}{Освітня програма \hfill}{\parbox {7 cm}{Інформатика\hfill}}\parbox {6 cm}{\hfill Семестр 1}}\par
{\parbox {5 cm} {Навчальний предмет \hfill}\parbox {7 cm}{Програмування \hfill}\parbox {6cm}{\hfill}\par}\vspace{-15pt} \end{scriptsize}}{\center Екзаменаційний білет \No \textbf { 2 }\par}}
{%beginitem
1. Які з наведених типів не є цілими (integral), якщо такі є серед перелічених?

double, int, long*, unsigned char?
%enditem
\par
\vspace{2mm}
%beginitem
2. Скільки змінних, що є вказівниками, тут оголошено?
double x, *px1, px2, ***px3;?
%enditem
\par
\vspace{2mm}
%beginitem
3. За наявних оголошень \code{void f(int*); int n;} які з виразів \code{f(n++), f(++n), f(\&n)} є синтаксично коректними?
%enditem
\par
\vspace{2mm}
%beginitem
4. Що таке просторова складність задачі?
%enditem
\par
\vspace{2mm}
%beginitem
5. Яку часову оцінку складності має алгоритм сортування бульбашкою?
%enditem
\par
\vspace{2mm}
%beginitem
6. Написати програму, що запитує у користувача значення дійсних параметрів $x$ та $y$, обчислює для них значення виразу та повiдомляє введенi данi та результат обчислення.
$$\frac{\ x - 61}{(\cos x - 0.7) (y + 68)}$$ 

%enditem
\par
%beginitem
7. Написати функцію, що для цілого числа знаходить кількість його простих дільників, що входять у його розклад рівно у другому степеню. Наприклад, для числа $3^2{\cdot}7^9{\cdot}11{\cdot}23^{2}$ таких дільників два – $3$ та $23$, тому функція має повернути $2$.

По завданням 2 та 3 оцінюється правильність та якість коду, а також економність обчислень.

%enditem
}
{\smallskip\smallskip \begin {scriptsize} Затверджено на засіданні кафедри теоретичної кібернетики, протокол \No 3  від   08.11.2025 р.\par\smallskip\smallskip\smallskip\smallskip
{\parbox {6.5 cm} {Зав. кафедри \dotfill Ю.В.~Крак}}\hspace {0.7cm} {\hbox to 6.5 cm { Екзаменатор \dotfill Т.О.~Карнаух}} \end{scriptsize}}
%end
\newpage
%begin
{{\begin {scriptsize}\par
{ \center  Київський національний університет імені Тараса Шевченка \par}
{\parbox {5 cm}{Освітня програма \hfill}{\parbox {7 cm}{Інформатика\hfill}}\parbox {6 cm}{\hfill Семестр 1}}\par
{\parbox {5 cm} {Навчальний предмет \hfill}\parbox {7 cm}{Програмування \hfill}\parbox {6cm}{\hfill}\par}\vspace{-15pt} \end{scriptsize}}{\center Екзаменаційний білет \No \textbf { 3 }\par}}
{%beginitem
1. Які з наведених типів не є цілими (integral), якщо такі є серед перелічених?

double, int, long, unsigned char?
%enditem
\par
\vspace{2mm}
%beginitem
2. Скільки змінних, що є вказівниками, тут оголошено?
long pointer1, pointer2, pointer3, pointer4;?
%enditem
\par
\vspace{2mm}
%beginitem
3. За наявних оголошень \code{void f(char); char c;} які з виразів \code{f(c*c), f(c+=2), f(\&c)} є синтаксично коректними?
%enditem
\par
\vspace{2mm}
%beginitem
4. У чому відмінність літералів від констант?
%enditem
\par
\vspace{2mm}
%beginitem
5. Яку часову оцінку складності має алгоритм сортування вибором?
%enditem
\par
\vspace{2mm}
%beginitem
6. Написати програму, що запитує у користувача значення дійсних параметрів $x$ та $y$, обчислює для них значення виразу та повiдомляє введенi данi та результат обчислення.
$$\frac{\ln x - 63}{(\arcsin x - 0.4) (y - 21)}$$ 

%enditem
\par
%beginitem
7. Написати функцію, що в інтервалі $[a, b)$ знаходить ціле число, що ділиться на найбільший степінь числа 3. Якщо таких чисел декілька, повернути найменше.

По завданням 2 та 3 оцінюється правильність та якість коду, а також економність обчислень.

%enditem
}
{\smallskip\smallskip \begin {scriptsize} Затверджено на засіданні кафедри теоретичної кібернетики, протокол \No 3  від   08.11.2025 р.\par\smallskip\smallskip\smallskip\smallskip
{\parbox {6.5 cm} {Зав. кафедри \dotfill Ю.В.~Крак}}\hspace {0.7cm} {\hbox to 6.5 cm { Екзаменатор \dotfill Т.О.~Карнаух}} \end{scriptsize}}
%end
\vspace{0.9cm}
%begin
{{\begin {scriptsize}\par
{ \center  Київський національний університет імені Тараса Шевченка \par}
{\parbox {5 cm}{Освітня програма \hfill}{\parbox {7 cm}{Інформатика\hfill}}\parbox {6 cm}{\hfill Семестр 1}}\par
{\parbox {5 cm} {Навчальний предмет \hfill}\parbox {7 cm}{Програмування \hfill}\parbox {6cm}{\hfill}\par}\vspace{-15pt} \end{scriptsize}}{\center Екзаменаційний білет \No \textbf { 4 }\par}}
{%beginitem
1. Які з наведених типів не є цілими (integral), якщо такі є серед перелічених?

int, signed char, unsigned long, string?
%enditem
\par
\vspace{2mm}
%beginitem
2. Скільки змінних, що є вказівниками, тут оголошено?
int *pn1, pn2, pn3, pn4;?
%enditem
\par
\vspace{2mm}
%beginitem
3. За наявних оголошень \code{void f(char*); char c;} які з виразів \code{f(c*c), f(c+=2), f(\&c)} є синтаксично коректними?
%enditem
\par
\vspace{2mm}
%beginitem
4. Що таке часова складність алгоритму?
%enditem
\par
\vspace{2mm}
%beginitem
5. Яку часову оцінку складності має алгоритм сортування злиттям?
%enditem
\par
\vspace{2mm}
%beginitem
6. Написати програму, що запитує у користувача значення дійсних параметрів $x$ та $y$, обчислює для них значення виразу та повiдомляє введенi данi та результат обчислення.
$$\frac{\sqrt x + 13}{(\ctg x + 0.5) (y + 17)}$$ 

%enditem
\par
%beginitem
7. Написати функцію, що повертає найбільше число з послідовності цілих (задається масивом), що має рівно три різних простих дільники. За відсутності має повертати None. Використовувати додаткові структури даних (у тому числі рядки) заборонено.

По завданням 2 та 3 оцінюється правильність та якість коду, а також економність обчислень.

%enditem
}
{\smallskip\smallskip \begin {scriptsize} Затверджено на засіданні кафедри теоретичної кібернетики, протокол \No 3  від   08.11.2025 р.\par\smallskip\smallskip\smallskip\smallskip
{\parbox {6.5 cm} {Зав. кафедри \dotfill Ю.В.~Крак}}\hspace {0.7cm} {\hbox to 6.5 cm { Екзаменатор \dotfill Т.О.~Карнаух}} \end{scriptsize}}
%end
\newpage
%begin
{{\begin {scriptsize}\par
{ \center  Київський національний університет імені Тараса Шевченка \par}
{\parbox {5 cm}{Освітня програма \hfill}{\parbox {7 cm}{Інформатика\hfill}}\parbox {6 cm}{\hfill Семестр 1}}\par
{\parbox {5 cm} {Навчальний предмет \hfill}\parbox {7 cm}{Програмування \hfill}\parbox {6cm}{\hfill}\par}\vspace{-15pt} \end{scriptsize}}{\center Екзаменаційний білет \No \textbf { 5 }\par}}
{%beginitem
1. Які з наведених типів не є цілими (integral), якщо такі є серед перелічених?

char, double, int, float, unsigned?
%enditem
\par
\vspace{2mm}
%beginitem
2. Скільки змінних, що є вказівниками, тут оголошено?
int n, m, *p1, p2;?
%enditem
\par
\vspace{2mm}
%beginitem
3. За наявних оголошень \code{void f(int); int n;} які з виразів \code{f(3), f(++n), f(\&n)} є синтаксично коректними?
%enditem
\par
\vspace{2mm}
%beginitem
4. Що таке С-рядок?
%enditem
\par
\vspace{2mm}
%beginitem
5. Яку часову оцінку складності має алгоритм швидкого сортування?
%enditem
\par
\vspace{2mm}
%beginitem
6. Написати програму, що запитує у користувача значення дійсних параметрів $x$ та $y$, обчислює для них значення виразу та повiдомляє введенi данi та результат обчислення.
$$\frac{\log_{10} x + 26}{(\arctg x + 0.6) (y + 18)}$$ 

%enditem
\par
%beginitem
7. Написати функцію, що для послідовності цілих (задається масивом) знаходить найдовший відрізок, в якому всі сусідні числа різні. Якщо таких відрізків кілька, то повернути останній (як індекс початку та довжину). Використовувати додаткові структури даних (у тому числі рядки) заборонено.

По завданням 2 та 3 оцінюється правильність та якість коду, а також економність обчислень.

%enditem
}
{\smallskip\smallskip \begin {scriptsize} Затверджено на засіданні кафедри теоретичної кібернетики, протокол \No 3  від   08.11.2025 р.\par\smallskip\smallskip\smallskip\smallskip
{\parbox {6.5 cm} {Зав. кафедри \dotfill Ю.В.~Крак}}\hspace {0.7cm} {\hbox to 6.5 cm { Екзаменатор \dotfill Т.О.~Карнаух}} \end{scriptsize}}
%end
\vspace{0.9cm}
%begin
{{\begin {scriptsize}\par
{ \center  Київський національний університет імені Тараса Шевченка \par}
{\parbox {5 cm}{Освітня програма \hfill}{\parbox {7 cm}{Інформатика\hfill}}\parbox {6 cm}{\hfill Семестр 1}}\par
{\parbox {5 cm} {Навчальний предмет \hfill}\parbox {7 cm}{Програмування \hfill}\parbox {6cm}{\hfill}\par}\vspace{-15pt} \end{scriptsize}}{\center Екзаменаційний білет \No \textbf { 6 }\par}}
{%beginitem
1. Які з наведених типів не є цілими (integral), якщо такі є серед перелічених?

bool, int, long, long long, unsigned char?
%enditem
\par
\vspace{2mm}
%beginitem
2. Скільки змінних, що є вказівниками, тут оголошено?
long *n1, *n2, n3, pn4;?
%enditem
\par
\vspace{2mm}
%beginitem
3. За наявних оголошень \code{void f(char\&); char c;} які з виразів \code{f(c*c), f(c++), f(\&c)} є синтаксично коректними?
%enditem
\par
\vspace{2mm}
%beginitem
4. Що означає термін ``часова складність у середньому''?
%enditem
\par
\vspace{2mm}
%beginitem
5. Яку часову оцінку складності має алгоритм сортування купою (він же деревом, пірамідою)?
%enditem
\par
\vspace{2mm}
%beginitem
6. Написати програму, що запитує у користувача значення дійсних параметрів $x$ та $y$, обчислює для них значення виразу та повiдомляє введенi данi та результат обчислення.
$$\frac{\log_{2} x - 85}{(\sin x - 0.2) (y - 22)}$$ 

%enditem
\par
%beginitem
7. Написати функцію, що знаходить найменше число з послідовності цілих (задається масивом), у десятковому запису якого нема цифри 7. Повертає індекс першого входження цього числа або $-1$ (за відсутності). Використовувати додаткові структури даних (у тому числі рядки) заборонено.

По завданням 2 та 3 оцінюється правильність та якість коду, а також економність обчислень.

%enditem
}
{\smallskip\smallskip \begin {scriptsize} Затверджено на засіданні кафедри теоретичної кібернетики, протокол \No 3  від   08.11.2025 р.\par\smallskip\smallskip\smallskip\smallskip
{\parbox {6.5 cm} {Зав. кафедри \dotfill Ю.В.~Крак}}\hspace {0.7cm} {\hbox to 6.5 cm { Екзаменатор \dotfill Т.О.~Карнаух}} \end{scriptsize}}
%end
\newpage
%begin
{{\begin {scriptsize}\par
{ \center  Київський національний університет імені Тараса Шевченка \par}
{\parbox {5 cm}{Освітня програма \hfill}{\parbox {7 cm}{Інформатика\hfill}}\parbox {6 cm}{\hfill Семестр 1}}\par
{\parbox {5 cm} {Навчальний предмет \hfill}\parbox {7 cm}{Програмування \hfill}\parbox {6cm}{\hfill}\par}\vspace{-15pt} \end{scriptsize}}{\center Екзаменаційний білет \No \textbf { 7 }\par}}
{%beginitem
1. Які з наведених типів не є цілими (integral), якщо такі є серед перелічених?

char, double, int, float, unsigned*?
%enditem
\par
\vspace{2mm}
%beginitem
2. Скільки змінних, що є вказівниками, тут оголошено?
double *px, *py, **ppx;?
%enditem
\par
\vspace{2mm}
%beginitem
3. За наявних оголошень \code{void f(double*); double x;} які з виразів \code{f(x*x), f(x+=2), f(\&x)} є синтаксично коректними?
%enditem
\par
\vspace{2mm}
%beginitem
4. Що таке масив?
%enditem
\par
\vspace{2mm}
%beginitem
5. Яку часову оцінку складності має алгоритм сортування вставками?
%enditem
\par
\vspace{2mm}
%beginitem
6. Написати програму, що запитує у користувача значення дійсних параметрів $x$ та $y$, обчислює для них значення виразу та повiдомляє введенi данi та результат обчислення.
$$\frac{\ln x - 48}{(\arccos x - 0.9) (y - 38)}$$ 

%enditem
\par
%beginitem
7. Написати функцію, що в інтервалі $[a, b)$ знаходить ціле число, що ділиться на найбільший степінь числа 3. Якщо таких чисел декілька, повернути найбільше.

По завданням 2 та 3 оцінюється правильність та якість коду, а також економність обчислень.

%enditem
}
{\smallskip\smallskip \begin {scriptsize} Затверджено на засіданні кафедри теоретичної кібернетики, протокол \No 3  від   08.11.2025 р.\par\smallskip\smallskip\smallskip\smallskip
{\parbox {6.5 cm} {Зав. кафедри \dotfill Ю.В.~Крак}}\hspace {0.7cm} {\hbox to 6.5 cm { Екзаменатор \dotfill Т.О.~Карнаух}} \end{scriptsize}}
%end
\vspace{0.9cm}
%begin
{{\begin {scriptsize}\par
{ \center  Київський національний університет імені Тараса Шевченка \par}
{\parbox {5 cm}{Освітня програма \hfill}{\parbox {7 cm}{Інформатика\hfill}}\parbox {6 cm}{\hfill Семестр 1}}\par
{\parbox {5 cm} {Навчальний предмет \hfill}\parbox {7 cm}{Програмування \hfill}\parbox {6cm}{\hfill}\par}\vspace{-15pt} \end{scriptsize}}{\center Екзаменаційний білет \No \textbf { 8 }\par}}
{%beginitem
1. Які з наведених типів не є цілими (integral), якщо такі є серед перелічених?

bool*, char, int, unsigned?
%enditem
\par
\vspace{2mm}
%beginitem
2. Скільки змінних, що є вказівниками, тут оголошено?
char *s1, pc, **s2;?
%enditem
\par
\vspace{2mm}
%beginitem
3. За наявних оголошень \code{void f(double); double x;} які з виразів \code{f(3.1), f(x+=2), f(\&x)} є синтаксично коректними?
%enditem
\par
\vspace{2mm}
%beginitem
4. Чим префіксний оператор ++ відрізняється від постфіксного?
%enditem
\par
\vspace{2mm}
%beginitem
5. Яку часову оцінку складності має алгоритм сортування злиттям?
%enditem
\par
\vspace{2mm}
%beginitem
6. Написати програму, що запитує у користувача значення дійсних параметрів $x$ та $y$, обчислює для них значення виразу та повiдомляє введенi данi та результат обчислення.
$$\frac{\pi x + 55}{(\tg x + 0.1) (y + 86)}$$ 

%enditem
\par
%beginitem
7. Написати функцію, що для інтервалу $[a, b)$ знаходить найближчу пару чисел з цього інтервалу, що мають від трьох до п'яти (включно) простих дільників.

По завданням 2 та 3 оцінюється правильність та якість коду, а також економність обчислень.

%enditem
}
{\smallskip\smallskip \begin {scriptsize} Затверджено на засіданні кафедри теоретичної кібернетики, протокол \No 3  від   08.11.2025 р.\par\smallskip\smallskip\smallskip\smallskip
{\parbox {6.5 cm} {Зав. кафедри \dotfill Ю.В.~Крак}}\hspace {0.7cm} {\hbox to 6.5 cm { Екзаменатор \dotfill Т.О.~Карнаух}} \end{scriptsize}}
%end
\newpage
%begin
{{\begin {scriptsize}\par
{ \center  Київський національний університет імені Тараса Шевченка \par}
{\parbox {5 cm}{Освітня програма \hfill}{\parbox {7 cm}{Інформатика\hfill}}\parbox {6 cm}{\hfill Семестр 1}}\par
{\parbox {5 cm} {Навчальний предмет \hfill}\parbox {7 cm}{Програмування \hfill}\parbox {6cm}{\hfill}\par}\vspace{-15pt} \end{scriptsize}}{\center Екзаменаційний білет \No \textbf { 9 }\par}}
{%beginitem
1. Які з наведених типів не є цілими (integral), якщо такі є серед перелічених?

bool, int, long, long long, unsigned char?
%enditem
\par
\vspace{2mm}
%beginitem
2. Скільки змінних, що є вказівниками, тут оголошено?
long *n1, *n2, n3, pn4;?
%enditem
\par
\vspace{2mm}
%beginitem
3. За наявних оголошень \code{void f(double\&); double x;} які з виразів \code{f(x*x), f(x+=2), f(\&x)} є синтаксично коректними?
%enditem
\par
\vspace{2mm}
%beginitem
4. Що таке часова складність задачі?
%enditem
\par
\vspace{2mm}
%beginitem
5. Яку часову оцінку складності має алгоритм сортування купою (він же деревом, пірамідою)?
%enditem
\par
\vspace{2mm}
%beginitem
6. Написати програму, що запитує у користувача значення дійсних параметрів $x$ та $y$, обчислює для них значення виразу та повiдомляє введенi данi та результат обчислення.
$$\frac{\log_{10} x + 24}{(\cos x + 0.8) (y - 28)}$$ 

%enditem
\par
%beginitem
7. Написати функцію, що для цілого числа знаходить кількість простих дільників, які входять у його розклад у найбільшому степені. Наприклад, для числа $3^5{\cdot}7^9{\cdot}11{\cdot}23^{9}$ таких дільників два – $23$ та $7$, тому функція має повернути $2$.

По завданням 2 та 3 оцінюється правильність та якість коду, а також економність обчислень.

%enditem
}
{\smallskip\smallskip \begin {scriptsize} Затверджено на засіданні кафедри теоретичної кібернетики, протокол \No 3  від   08.11.2025 р.\par\smallskip\smallskip\smallskip\smallskip
{\parbox {6.5 cm} {Зав. кафедри \dotfill Ю.В.~Крак}}\hspace {0.7cm} {\hbox to 6.5 cm { Екзаменатор \dotfill Т.О.~Карнаух}} \end{scriptsize}}
%end
\vspace{0.9cm}
%begin
{{\begin {scriptsize}\par
{ \center  Київський національний університет імені Тараса Шевченка \par}
{\parbox {5 cm}{Освітня програма \hfill}{\parbox {7 cm}{Інформатика\hfill}}\parbox {6 cm}{\hfill Семестр 1}}\par
{\parbox {5 cm} {Навчальний предмет \hfill}\parbox {7 cm}{Програмування \hfill}\parbox {6cm}{\hfill}\par}\vspace{-15pt} \end{scriptsize}}{\center Екзаменаційний білет \No \textbf { 10 }\par}}
{%beginitem
1. Які з наведених типів не є цілими (integral), якщо такі є серед перелічених?

bool, char, int, unsigned?
%enditem
\par
\vspace{2mm}
%beginitem
2. Скільки змінних, що є вказівниками, тут оголошено?
int *pn1, pn2, pn3, pn4;?
%enditem
\par
\vspace{2mm}
%beginitem
3. За наявних оголошень \code{void f(int*); int n;} які з виразів \code{f(n++), f(++n), f(\&n)} є синтаксично коректними?
%enditem
\par
\vspace{2mm}
%beginitem
4. Що таке тип даних?
%enditem
\par
\vspace{2mm}
%beginitem
5. Яку часову оцінку складності має алгоритм сортування вибором?
%enditem
\par
\vspace{2mm}
%beginitem
6. Написати програму, що запитує у користувача значення дійсних параметрів $x$ та $y$, обчислює для них значення виразу та повiдомляє введенi данi та результат обчислення.
$$\frac{\log_{2} x - 60}{(\ctg x - 0.1) (y + 44)}$$ 

%enditem
\par
%beginitem
7. Написати функцію, яка для цілого числа знаходить кількість його різних простих дільників, що мають у десятковому запису цифру 1.

По завданням 2 та 3 оцінюється правильність та якість коду, а також економність обчислень.

%enditem
}
{\smallskip\smallskip \begin {scriptsize} Затверджено на засіданні кафедри теоретичної кібернетики, протокол \No 3  від   08.11.2025 р.\par\smallskip\smallskip\smallskip\smallskip
{\parbox {6.5 cm} {Зав. кафедри \dotfill Ю.В.~Крак}}\hspace {0.7cm} {\hbox to 6.5 cm { Екзаменатор \dotfill Т.О.~Карнаух}} \end{scriptsize}}
%end
\newpage
%begin
{{\begin {scriptsize}\par
{ \center  Київський національний університет імені Тараса Шевченка \par}
{\parbox {5 cm}{Освітня програма \hfill}{\parbox {7 cm}{Інформатика\hfill}}\parbox {6 cm}{\hfill Семестр 1}}\par
{\parbox {5 cm} {Навчальний предмет \hfill}\parbox {7 cm}{Програмування \hfill}\parbox {6cm}{\hfill}\par}\vspace{-15pt} \end{scriptsize}}{\center Екзаменаційний білет \No \textbf { 11 }\par}}
{%beginitem
1. Які з наведених типів не є цілими (integral), якщо такі є серед перелічених?

char, double, int, float, unsigned?
%enditem
\par
\vspace{2mm}
%beginitem
2. Скільки змінних, що є вказівниками, тут оголошено?
int n, m, *p1, p2;?
%enditem
\par
\vspace{2mm}
%beginitem
3. За наявних оголошень \code{void f(int\&); int n;} які з виразів \code{f(13), f(++n), f(\&n)} є синтаксично коректними?
%enditem
\par
\vspace{2mm}
%beginitem
4. Які параметри може приймати функція main? Що вони позначають та яких вони типів?
%enditem
\par
\vspace{2mm}
%beginitem
5. Яку часову оцінку складності має алгоритм сортування вставками?
%enditem
\par
\vspace{2mm}
%beginitem
6. Написати програму, що запитує у користувача значення дійсних параметрів $x$ та $y$, обчислює для них значення виразу та повiдомляє введенi данi та результат обчислення.
$$\frac{\ x + 37}{(\sin x + 0.9) (y + 54)}$$ 

%enditem
\par
%beginitem
7. Написати функцію, що повертає число з послідовності цілих (задається масивом), що має найбільшу кількість різних простих дільників. Якщо таких чисел кілька, то повернути найменше. Використовувати додаткові структури даних (у тому числі рядки) заборонено.

По завданням 2 та 3 оцінюється правильність та якість коду, а також економність обчислень.

%enditem
}
{\smallskip\smallskip \begin {scriptsize} Затверджено на засіданні кафедри теоретичної кібернетики, протокол \No 3  від   08.11.2025 р.\par\smallskip\smallskip\smallskip\smallskip
{\parbox {6.5 cm} {Зав. кафедри \dotfill Ю.В.~Крак}}\hspace {0.7cm} {\hbox to 6.5 cm { Екзаменатор \dotfill Т.О.~Карнаух}} \end{scriptsize}}
%end
\vspace{0.9cm}
%begin
{{\begin {scriptsize}\par
{ \center  Київський національний університет імені Тараса Шевченка \par}
{\parbox {5 cm}{Освітня програма \hfill}{\parbox {7 cm}{Інформатика\hfill}}\parbox {6 cm}{\hfill Семестр 1}}\par
{\parbox {5 cm} {Навчальний предмет \hfill}\parbox {7 cm}{Програмування \hfill}\parbox {6cm}{\hfill}\par}\vspace{-15pt} \end{scriptsize}}{\center Екзаменаційний білет \No \textbf { 12 }\par}}
{%beginitem
1. Які з наведених типів не є цілими (integral), якщо такі є серед перелічених?

bool, int, long, long long, unsigned char?
%enditem
\par
\vspace{2mm}
%beginitem
2. Скільки змінних, що є вказівниками, тут оголошено?
double x, *px1, px2, ***px3;?
%enditem
\par
\vspace{2mm}
%beginitem
3. За наявних оголошень \code{void f(char*); char c;} які з виразів \code{f(c*c), f(c+=2), f(\&c)} є синтаксично коректними?
%enditem
\par
\vspace{2mm}
%beginitem
4. Що таке тип даних?
%enditem
\par
\vspace{2mm}
%beginitem
5. Яку часову оцінку складності має алгоритм сортування бульбашкою?
%enditem
\par
\vspace{2mm}
%beginitem
6. Написати програму, що запитує у користувача значення дійсних параметрів $x$ та $y$, обчислює для них значення виразу та повiдомляє введенi данi та результат обчислення.
$$\frac{\sqrt x - 83}{(\arcsin x - 0.8) (y - 59)}$$ 

%enditem
\par
%beginitem
7. Написати функцію, що для послідовності цілих (задається масивом) знаходить довжину найдовшого відрізку, що складається з чисел, у сімковому запису якого нема цифри 2. За відсутності має повертати $0$. Використовувати додаткові структури даних (у тому числі рядки) заборонено.

По завданням 2 та 3 оцінюється правильність та якість коду, а також економність обчислень.

%enditem
}
{\smallskip\smallskip \begin {scriptsize} Затверджено на засіданні кафедри теоретичної кібернетики, протокол \No 3  від   08.11.2025 р.\par\smallskip\smallskip\smallskip\smallskip
{\parbox {6.5 cm} {Зав. кафедри \dotfill Ю.В.~Крак}}\hspace {0.7cm} {\hbox to 6.5 cm { Екзаменатор \dotfill Т.О.~Карнаух}} \end{scriptsize}}
%end
\newpage
%begin
{{\begin {scriptsize}\par
{ \center  Київський національний університет імені Тараса Шевченка \par}
{\parbox {5 cm}{Освітня програма \hfill}{\parbox {7 cm}{Інформатика\hfill}}\parbox {6 cm}{\hfill Семестр 1}}\par
{\parbox {5 cm} {Навчальний предмет \hfill}\parbox {7 cm}{Програмування \hfill}\parbox {6cm}{\hfill}\par}\vspace{-15pt} \end{scriptsize}}{\center Екзаменаційний білет \No \textbf { 13 }\par}}
{%beginitem
1. Які з наведених типів не є цілими (integral), якщо такі є серед перелічених?

int*, signed char, unsigned long, string?
%enditem
\par
\vspace{2mm}
%beginitem
2. Скільки змінних, що є вказівниками, тут оголошено?
long pointer1, pointer2, pointer3, pointer4;?
%enditem
\par
\vspace{2mm}
%beginitem
3. За наявних оголошень \code{void f(char\&); char c;} які з виразів \code{f(c*c), f(c++), f(\&c)} є синтаксично коректними?
%enditem
\par
\vspace{2mm}
%beginitem
4. Що означає термін ``часова складність у середньому''?
%enditem
\par
\vspace{2mm}
%beginitem
5. Яку часову оцінку складності має алгоритм швидкого сортування?
%enditem
\par
\vspace{2mm}
%beginitem
6. Написати програму, що запитує у користувача значення дійсних параметрів $x$ та $y$, обчислює для них значення виразу та повiдомляє введенi данi та результат обчислення.
$$\frac{\log_{10} x + 98}{(\arccos x + 0.7) (y + 36)}$$ 

%enditem
\par
%beginitem
7. Написати функцію, що для цілого числа знаходить простий дільник, який входить у його розклад ц найбільшому степені. Якщо таких кілька, то повернути найбільший. Наприклад, для числа $2^{9}{\cdot}3^5{\cdot}7^9{\cdot}11$ функція має повернути $7$.

По завданням 2 та 3 оцінюється правильність та якість коду, а також економність обчислень.

%enditem
}
{\smallskip\smallskip \begin {scriptsize} Затверджено на засіданні кафедри теоретичної кібернетики, протокол \No 3  від   08.11.2025 р.\par\smallskip\smallskip\smallskip\smallskip
{\parbox {6.5 cm} {Зав. кафедри \dotfill Ю.В.~Крак}}\hspace {0.7cm} {\hbox to 6.5 cm { Екзаменатор \dotfill Т.О.~Карнаух}} \end{scriptsize}}
%end
\vspace{0.9cm}
%begin
{{\begin {scriptsize}\par
{ \center  Київський національний університет імені Тараса Шевченка \par}
{\parbox {5 cm}{Освітня програма \hfill}{\parbox {7 cm}{Інформатика\hfill}}\parbox {6 cm}{\hfill Семестр 1}}\par
{\parbox {5 cm} {Навчальний предмет \hfill}\parbox {7 cm}{Програмування \hfill}\parbox {6cm}{\hfill}\par}\vspace{-15pt} \end{scriptsize}}{\center Екзаменаційний білет \No \textbf { 14 }\par}}
{%beginitem
1. Які з наведених типів не є цілими (integral), якщо такі є серед перелічених?

double, int, long*, unsigned char?
%enditem
\par
\vspace{2mm}
%beginitem
2. Скільки змінних, що є вказівниками, тут оголошено?
double *px, *py, **ppx;?
%enditem
\par
\vspace{2mm}
%beginitem
3. За наявних оголошень \code{void f(int); int n;} які з виразів \code{f(3), f(++n), f(\&n)} є синтаксично коректними?
%enditem
\par
\vspace{2mm}
%beginitem
4. Які параметри може приймати функція main? Що вони позначають та яких вони типів?
%enditem
\par
\vspace{2mm}
%beginitem
5. Яку часову оцінку складності має алгоритм швидкого сортування?
%enditem
\par
\vspace{2mm}
%beginitem
6. Написати програму, що запитує у користувача значення дійсних параметрів $x$ та $y$, обчислює для них значення виразу та повiдомляє введенi данi та результат обчислення.
$$\frac{\pi x - 30}{(\arctg x - 0.5) (y - 43)}$$ 

%enditem
\par
%beginitem
7. Написати функцію, що для цілого числа знаходить простий дільник, який входить у його розклад у найбільшому степені. Якщо таких кілька, то повернути найменший. Наприклад, для числа $2^{9}{\cdot}3^5{\cdot}7^9{\cdot}11$ функція має повернути $2$.

По завданням 2 та 3 оцінюється правильність та якість коду, а також економність обчислень.

%enditem
}
{\smallskip\smallskip \begin {scriptsize} Затверджено на засіданні кафедри теоретичної кібернетики, протокол \No 3  від   08.11.2025 р.\par\smallskip\smallskip\smallskip\smallskip
{\parbox {6.5 cm} {Зав. кафедри \dotfill Ю.В.~Крак}}\hspace {0.7cm} {\hbox to 6.5 cm { Екзаменатор \dotfill Т.О.~Карнаух}} \end{scriptsize}}
%end
\newpage
%begin
{{\begin {scriptsize}\par
{ \center  Київський національний університет імені Тараса Шевченка \par}
{\parbox {5 cm}{Освітня програма \hfill}{\parbox {7 cm}{Інформатика\hfill}}\parbox {6 cm}{\hfill Семестр 1}}\par
{\parbox {5 cm} {Навчальний предмет \hfill}\parbox {7 cm}{Програмування \hfill}\parbox {6cm}{\hfill}\par}\vspace{-15pt} \end{scriptsize}}{\center Екзаменаційний білет \No \textbf { 15 }\par}}
{%beginitem
1. Які з наведених типів не є цілими (integral), якщо такі є серед перелічених?

int, signed char, unsigned long, string?
%enditem
\par
\vspace{2mm}
%beginitem
2. Скільки змінних, що є вказівниками, тут оголошено?
char *s1, pc, **s2;?
%enditem
\par
\vspace{2mm}
%beginitem
3. За наявних оголошень \code{void f(double); double x;} які з виразів \code{f(3.1), f(x+=2), f(\&x)} є синтаксично коректними?
%enditem
\par
\vspace{2mm}
%beginitem
4. Що таке просторова складність задачі?
%enditem
\par
\vspace{2mm}
%beginitem
5. Яку часову оцінку складності має алгоритм сортування вибором?
%enditem
\par
\vspace{2mm}
%beginitem
6. Написати програму, що запитує у користувача значення дійсних параметрів $x$ та $y$, обчислює для них значення виразу та повiдомляє введенi данi та результат обчислення.
$$\frac{\log_{2} x - 92}{(\arcsin x + 0.6) (y - 79)}$$ 

%enditem
\par
%beginitem
7. Написати функцію, що для інтервалу $[a, b)$ знаходить найближчу пару чисел з цього інтервалу, що мають не менше трьох простих дільників та мають цифру 3 у своєму десятковому запису.

По завданням 2 та 3 оцінюється правильність та якість коду, а також економність обчислень.

%enditem
}
{\smallskip\smallskip \begin {scriptsize} Затверджено на засіданні кафедри теоретичної кібернетики, протокол \No 3  від   08.11.2025 р.\par\smallskip\smallskip\smallskip\smallskip
{\parbox {6.5 cm} {Зав. кафедри \dotfill Ю.В.~Крак}}\hspace {0.7cm} {\hbox to 6.5 cm { Екзаменатор \dotfill Т.О.~Карнаух}} \end{scriptsize}}
%end
\vspace{0.9cm}
%begin
{{\begin {scriptsize}\par
{ \center  Київський національний університет імені Тараса Шевченка \par}
{\parbox {5 cm}{Освітня програма \hfill}{\parbox {7 cm}{Інформатика\hfill}}\parbox {6 cm}{\hfill Семестр 1}}\par
{\parbox {5 cm} {Навчальний предмет \hfill}\parbox {7 cm}{Програмування \hfill}\parbox {6cm}{\hfill}\par}\vspace{-15pt} \end{scriptsize}}{\center Екзаменаційний білет \No \textbf { 16 }\par}}
{%beginitem
1. Які з наведених типів не є цілими (integral), якщо такі є серед перелічених?

bool, char, int, unsigned?
%enditem
\par
\vspace{2mm}
%beginitem
2. Скільки змінних, що є вказівниками, тут оголошено?
long *n1, *n2, n3, pn4;?
%enditem
\par
\vspace{2mm}
%beginitem
3. За наявних оголошень \code{void f(double*); double x;} які з виразів \code{f(x*x), f(x+=2), f(\&x)} є синтаксично коректними?
%enditem
\par
\vspace{2mm}
%beginitem
4. У чому відмінність літералів від констант?
%enditem
\par
\vspace{2mm}
%beginitem
5. Яку часову оцінку складності має алгоритм сортування бульбашкою?
%enditem
\par
\vspace{2mm}
%beginitem
6. Написати програму, що запитує у користувача значення дійсних параметрів $x$ та $y$, обчислює для них значення виразу та повiдомляє введенi данi та результат обчислення.
$$\frac{\ln x + 90}{(\sin x - 0.2) (y + 65)}$$ 

%enditem
\par
%beginitem
7. Написати функцію, що повертає найбільше число з послідовності цілих (задається масивом), що має не більше трьох різних простих дільників. За відсутності має повертати None. Використовувати додаткові структури даних (у тому числі рядки) заборонено.

По завданням 2 та 3 оцінюється правильність та якість коду, а також економність обчислень.

%enditem
}
{\smallskip\smallskip \begin {scriptsize} Затверджено на засіданні кафедри теоретичної кібернетики, протокол \No 3  від   08.11.2025 р.\par\smallskip\smallskip\smallskip\smallskip
{\parbox {6.5 cm} {Зав. кафедри \dotfill Ю.В.~Крак}}\hspace {0.7cm} {\hbox to 6.5 cm { Екзаменатор \dotfill Т.О.~Карнаух}} \end{scriptsize}}
%end
\newpage
%begin
{{\begin {scriptsize}\par
{ \center  Київський національний університет імені Тараса Шевченка \par}
{\parbox {5 cm}{Освітня програма \hfill}{\parbox {7 cm}{Інформатика\hfill}}\parbox {6 cm}{\hfill Семестр 1}}\par
{\parbox {5 cm} {Навчальний предмет \hfill}\parbox {7 cm}{Програмування \hfill}\parbox {6cm}{\hfill}\par}\vspace{-15pt} \end{scriptsize}}{\center Екзаменаційний білет \No \textbf { 17 }\par}}
{%beginitem
1. Які з наведених типів не є цілими (integral), якщо такі є серед перелічених?

char, double, int, float, unsigned*?
%enditem
\par
\vspace{2mm}
%beginitem
2. Скільки змінних, що є вказівниками, тут оголошено?
double x, *px1, px2, ***px3;?
%enditem
\par
\vspace{2mm}
%beginitem
3. За наявних оголошень \code{void f(double\&); double x;} які з виразів \code{f(x*x), f(x+=2), f(\&x)} є синтаксично коректними?
%enditem
\par
\vspace{2mm}
%beginitem
4. Що таке масив?
%enditem
\par
\vspace{2mm}
%beginitem
5. Яку часову оцінку складності має алгоритм сортування злиттям?
%enditem
\par
\vspace{2mm}
%beginitem
6. Написати програму, що запитує у користувача значення дійсних параметрів $x$ та $y$, обчислює для них значення виразу та повiдомляє введенi данi та результат обчислення.
$$\frac{\ x - 87}{(\cos x + 0.3) (y - 34)}$$ 

%enditem
\par
%beginitem
7. Написати функцію, що повертає число з послідовності цілих (задається масивом), що має парну кількість різних простих дільників. Якщо таких чисел кілька, то повернути найбільше. Використовувати додаткові структури даних (у тому числі рядки) заборонено.

По завданням 2 та 3 оцінюється правильність та якість коду, а також економність обчислень.

%enditem
}
{\smallskip\smallskip \begin {scriptsize} Затверджено на засіданні кафедри теоретичної кібернетики, протокол \No 3  від   08.11.2025 р.\par\smallskip\smallskip\smallskip\smallskip
{\parbox {6.5 cm} {Зав. кафедри \dotfill Ю.В.~Крак}}\hspace {0.7cm} {\hbox to 6.5 cm { Екзаменатор \dotfill Т.О.~Карнаух}} \end{scriptsize}}
%end
\vspace{0.9cm}
%begin
{{\begin {scriptsize}\par
{ \center  Київський національний університет імені Тараса Шевченка \par}
{\parbox {5 cm}{Освітня програма \hfill}{\parbox {7 cm}{Інформатика\hfill}}\parbox {6 cm}{\hfill Семестр 1}}\par
{\parbox {5 cm} {Навчальний предмет \hfill}\parbox {7 cm}{Програмування \hfill}\parbox {6cm}{\hfill}\par}\vspace{-15pt} \end{scriptsize}}{\center Екзаменаційний білет \No \textbf { 18 }\par}}
{%beginitem
1. Які з наведених типів не є цілими (integral), якщо такі є серед перелічених?

bool, int, long, long long, unsigned char?
%enditem
\par
\vspace{2mm}
%beginitem
2. Скільки змінних, що є вказівниками, тут оголошено?
long pointer1, pointer2, pointer3, pointer4;?
%enditem
\par
\vspace{2mm}
%beginitem
3. За наявних оголошень \code{void f(char); char c;} які з виразів \code{f(c*c), f(c+=2), f(\&c)} є синтаксично коректними?
%enditem
\par
\vspace{2mm}
%beginitem
4. Що таке С-рядок?
%enditem
\par
\vspace{2mm}
%beginitem
5. Яку часову оцінку складності має алгоритм швидкого сортування?
%enditem
\par
\vspace{2mm}
%beginitem
6. Написати програму, що запитує у користувача значення дійсних параметрів $x$ та $y$, обчислює для них значення виразу та повiдомляє введенi данi та результат обчислення.
$$\frac{\sqrt x + 76}{(\ctg x - 0.4) (y + 32)}$$ 

%enditem
\par
%beginitem
7. Написати функцію, що для послідовності цілих (задається масивом) знаходить найдовший відрізок, в якому всі сусідні числа різні. Якщо таких відрізків кілька, то повернути перший (як індекс початку та довжину). Використовувати додаткові структури даних (у тому числі рядки) заборонено.

По завданням 2 та 3 оцінюється правильність та якість коду, а також економність обчислень.

%enditem
}
{\smallskip\smallskip \begin {scriptsize} Затверджено на засіданні кафедри теоретичної кібернетики, протокол \No 3  від   08.11.2025 р.\par\smallskip\smallskip\smallskip\smallskip
{\parbox {6.5 cm} {Зав. кафедри \dotfill Ю.В.~Крак}}\hspace {0.7cm} {\hbox to 6.5 cm { Екзаменатор \dotfill Т.О.~Карнаух}} \end{scriptsize}}
%end
\newpage
%begin
{{\begin {scriptsize}\par
{ \center  Київський національний університет імені Тараса Шевченка \par}
{\parbox {5 cm}{Освітня програма \hfill}{\parbox {7 cm}{Інформатика\hfill}}\parbox {6 cm}{\hfill Семестр 1}}\par
{\parbox {5 cm} {Навчальний предмет \hfill}\parbox {7 cm}{Програмування \hfill}\parbox {6cm}{\hfill}\par}\vspace{-15pt} \end{scriptsize}}{\center Екзаменаційний білет \No \textbf { 19 }\par}}
{%beginitem
1. Які з наведених типів не є цілими (integral), якщо такі є серед перелічених?

double, int, long, unsigned char?
%enditem
\par
\vspace{2mm}
%beginitem
2. Скільки змінних, що є вказівниками, тут оголошено?
int n, m, *p1, p2;?
%enditem
\par
\vspace{2mm}
%beginitem
3. За наявних оголошень \code{void f(int); int n;} які з виразів \code{f(3), f(++n), f(\&n)} є синтаксично коректними?
%enditem
\par
\vspace{2mm}
%beginitem
4. Що таке часова складність алгоритму?
%enditem
\par
\vspace{2mm}
%beginitem
5. Яку часову оцінку складності має алгоритм сортування купою (він же деревом, пірамідою)?
%enditem
\par
\vspace{2mm}
%beginitem
6. Написати програму, що запитує у користувача значення дійсних параметрів $x$ та $y$, обчислює для них значення виразу та повiдомляє введенi данi та результат обчислення.
$$\frac{\log_{10} x - 73}{(\tg x + 0.9) (y + 41)}$$ 

%enditem
\par
%beginitem
7. Написати функцію, що повертає найбільше число з послідовності цілих (задається масивом), у сімковому запису якого нема цифри 2. За відсутності має повертати None. Використовувати додаткові структури даних (у тому числі рядки) заборонено.

По завданням 2 та 3 оцінюється правильність та якість коду, а також економність обчислень.

%enditem
}
{\smallskip\smallskip \begin {scriptsize} Затверджено на засіданні кафедри теоретичної кібернетики, протокол \No 3  від   08.11.2025 р.\par\smallskip\smallskip\smallskip\smallskip
{\parbox {6.5 cm} {Зав. кафедри \dotfill Ю.В.~Крак}}\hspace {0.7cm} {\hbox to 6.5 cm { Екзаменатор \dotfill Т.О.~Карнаух}} \end{scriptsize}}
%end
\vspace{0.9cm}
%begin
{{\begin {scriptsize}\par
{ \center  Київський національний університет імені Тараса Шевченка \par}
{\parbox {5 cm}{Освітня програма \hfill}{\parbox {7 cm}{Інформатика\hfill}}\parbox {6 cm}{\hfill Семестр 1}}\par
{\parbox {5 cm} {Навчальний предмет \hfill}\parbox {7 cm}{Програмування \hfill}\parbox {6cm}{\hfill}\par}\vspace{-15pt} \end{scriptsize}}{\center Екзаменаційний білет \No \textbf { 20 }\par}}
{%beginitem
1. Які з наведених типів не є цілими (integral), якщо такі є серед перелічених?

bool*, char, int, unsigned?
%enditem
\par
\vspace{2mm}
%beginitem
2. Скільки змінних, що є вказівниками, тут оголошено?
double *px, *py, **ppx;?
%enditem
\par
\vspace{2mm}
%beginitem
3. За наявних оголошень \code{void f(char); char c;} які з виразів \code{f(c*c), f(c+=2), f(\&c)} є синтаксично коректними?
%enditem
\par
\vspace{2mm}
%beginitem
4. Що таке просторова складність алгоритму?
%enditem
\par
\vspace{2mm}
%beginitem
5. Яку часову оцінку складності має алгоритм швидкого сортування?
%enditem
\par
\vspace{2mm}
%beginitem
6. Написати програму, що запитує у користувача значення дійсних параметрів $x$ та $y$, обчислює для них значення виразу та повiдомляє введенi данi та результат обчислення.
$$\frac{\ln x + 53}{(\arccos x - 0.3) (y - 75)}$$ 

%enditem
\par
%beginitem
7. Написати функцію, що знаходить найбільше число з послідовності цілих (задається масивом), у десятковому запису якого нема цифри 4. Повертає індекс першого входження цього числа або $-1$ (за відсутності). Використовувати додаткові структури даних (у тому числі рядки) заборонено.

По завданням 2 та 3 оцінюється правильність та якість коду, а також економність обчислень.

%enditem
}
{\smallskip\smallskip \begin {scriptsize} Затверджено на засіданні кафедри теоретичної кібернетики, протокол \No 3  від   08.11.2025 р.\par\smallskip\smallskip\smallskip\smallskip
{\parbox {6.5 cm} {Зав. кафедри \dotfill Ю.В.~Крак}}\hspace {0.7cm} {\hbox to 6.5 cm { Екзаменатор \dotfill Т.О.~Карнаух}} \end{scriptsize}}
%end
\newpage
%begin
{{\begin {scriptsize}\par
{ \center  Київський національний університет імені Тараса Шевченка \par}
{\parbox {5 cm}{Освітня програма \hfill}{\parbox {7 cm}{Інформатика\hfill}}\parbox {6 cm}{\hfill Семестр 1}}\par
{\parbox {5 cm} {Навчальний предмет \hfill}\parbox {7 cm}{Програмування \hfill}\parbox {6cm}{\hfill}\par}\vspace{-15pt} \end{scriptsize}}{\center Екзаменаційний білет \No \textbf { 21 }\par}}
{%beginitem
1. Які з наведених типів не є цілими (integral), якщо такі є серед перелічених?

bool, int, long, long long, unsigned char?
%enditem
\par
\vspace{2mm}
%beginitem
2. Скільки змінних, що є вказівниками, тут оголошено?
int *pn1, pn2, pn3, pn4;?
%enditem
\par
\vspace{2mm}
%beginitem
3. За наявних оголошень \code{void f(double*); double x;} які з виразів \code{f(x*x), f(x+=2), f(\&x)} є синтаксично коректними?
%enditem
\par
\vspace{2mm}
%beginitem
4. Чим префіксний оператор ++ відрізняється від постфіксного?
%enditem
\par
\vspace{2mm}
%beginitem
5. Яку часову оцінку складності має алгоритм сортування вставками?
%enditem
\par
\vspace{2mm}
%beginitem
6. Написати програму, що запитує у користувача значення дійсних параметрів $x$ та $y$, обчислює для них значення виразу та повiдомляє введенi данi та результат обчислення.
$$\frac{\sqrt x - 96}{(\arctg x + 0.4) (y + 33)}$$ 

%enditem
\par
%beginitem
7. Написати функцію, що для інтервалу $[a, b)$ знаходить найближчу пару чисел з цього інтервалу, що кожне з них має непарну кількість різних простих дільників. Якщо таких пар кілька, то цікавить пара з найбільшою сумою.

По завданням 2 та 3 оцінюється правильність та якість коду, а також економність обчислень.

%enditem
}
{\smallskip\smallskip \begin {scriptsize} Затверджено на засіданні кафедри теоретичної кібернетики, протокол \No 3  від   08.11.2025 р.\par\smallskip\smallskip\smallskip\smallskip
{\parbox {6.5 cm} {Зав. кафедри \dotfill Ю.В.~Крак}}\hspace {0.7cm} {\hbox to 6.5 cm { Екзаменатор \dotfill Т.О.~Карнаух}} \end{scriptsize}}
%end
\vspace{0.9cm}
%begin
{{\begin {scriptsize}\par
{ \center  Київський національний університет імені Тараса Шевченка \par}
{\parbox {5 cm}{Освітня програма \hfill}{\parbox {7 cm}{Інформатика\hfill}}\parbox {6 cm}{\hfill Семестр 1}}\par
{\parbox {5 cm} {Навчальний предмет \hfill}\parbox {7 cm}{Програмування \hfill}\parbox {6cm}{\hfill}\par}\vspace{-15pt} \end{scriptsize}}{\center Екзаменаційний білет \No \textbf { 22 }\par}}
{%beginitem
1. Які з наведених типів не є цілими (integral), якщо такі є серед перелічених?

int*, signed char, unsigned long, string?
%enditem
\par
\vspace{2mm}
%beginitem
2. Скільки змінних, що є вказівниками, тут оголошено?
int *pn1, pn2, pn3, pn4;?
%enditem
\par
\vspace{2mm}
%beginitem
3. За наявних оголошень \code{void f(double); double x;} які з виразів \code{f(3.1), f(x+=2), f(\&x)} є синтаксично коректними?
%enditem
\par
\vspace{2mm}
%beginitem
4. Що таке часова складність задачі?
%enditem
\par
\vspace{2mm}
%beginitem
5. Яку часову оцінку складності має алгоритм сортування вибором?
%enditem
\par
\vspace{2mm}
%beginitem
6. Написати програму, що запитує у користувача значення дійсних параметрів $x$ та $y$, обчислює для них значення виразу та повiдомляє введенi данi та результат обчислення.
$$\frac{\ x + 78}{(\cos x - 0.1) (y - 88)}$$ 

%enditem
\par
%beginitem
7. Написати функцію, що для інтервалу $[a, b)$ знаходить кількість чисел, квадрати яких діляться на суму квадратів їх цифр. (Розглядаємо десятковий запис.)

По завданням 2 та 3 оцінюється правильність та якість коду, а також економність обчислень.

%enditem
}
{\smallskip\smallskip \begin {scriptsize} Затверджено на засіданні кафедри теоретичної кібернетики, протокол \No 3  від   08.11.2025 р.\par\smallskip\smallskip\smallskip\smallskip
{\parbox {6.5 cm} {Зав. кафедри \dotfill Ю.В.~Крак}}\hspace {0.7cm} {\hbox to 6.5 cm { Екзаменатор \dotfill Т.О.~Карнаух}} \end{scriptsize}}
%end
\newpage
%begin
{{\begin {scriptsize}\par
{ \center  Київський національний університет імені Тараса Шевченка \par}
{\parbox {5 cm}{Освітня програма \hfill}{\parbox {7 cm}{Інформатика\hfill}}\parbox {6 cm}{\hfill Семестр 1}}\par
{\parbox {5 cm} {Навчальний предмет \hfill}\parbox {7 cm}{Програмування \hfill}\parbox {6cm}{\hfill}\par}\vspace{-15pt} \end{scriptsize}}{\center Екзаменаційний білет \No \textbf { 23 }\par}}
{%beginitem
1. Які з наведених типів не є цілими (integral), якщо такі є серед перелічених?

char, double, int, float, unsigned*?
%enditem
\par
\vspace{2mm}
%beginitem
2. Скільки змінних, що є вказівниками, тут оголошено?
long pointer1, pointer2, pointer3, pointer4;?
%enditem
\par
\vspace{2mm}
%beginitem
3. За наявних оголошень \code{void f(double\&); double x;} які з виразів \code{f(x*x), f(x+=2), f(\&x)} є синтаксично коректними?
%enditem
\par
\vspace{2mm}
%beginitem
4. Що таке просторова складність задачі?
%enditem
\par
\vspace{2mm}
%beginitem
5. Яку часову оцінку складності має алгоритм сортування злиттям?
%enditem
\par
\vspace{2mm}
%beginitem
6. Написати програму, що запитує у користувача значення дійсних параметрів $x$ та $y$, обчислює для них значення виразу та повiдомляє введенi данi та результат обчислення.
$$\frac{\pi x + 64}{(\ctg x - 0.8) (y - 99)}$$ 

%enditem
\par
%beginitem
7. Написати функцію, що для послідовності цілих (задається масивом) знаходить довжину найдовшого відрізку, в якому всі сусідні числа відрізняються якнайменш на 2. Використовувати додаткові структури даних (у тому числі рядки) заборонено.

По завданням 2 та 3 оцінюється правильність та якість коду, а також економність обчислень.

%enditem
}
{\smallskip\smallskip \begin {scriptsize} Затверджено на засіданні кафедри теоретичної кібернетики, протокол \No 3  від   08.11.2025 р.\par\smallskip\smallskip\smallskip\smallskip
{\parbox {6.5 cm} {Зав. кафедри \dotfill Ю.В.~Крак}}\hspace {0.7cm} {\hbox to 6.5 cm { Екзаменатор \dotfill Т.О.~Карнаух}} \end{scriptsize}}
%end
\vspace{0.9cm}
%begin
{{\begin {scriptsize}\par
{ \center  Київський національний університет імені Тараса Шевченка \par}
{\parbox {5 cm}{Освітня програма \hfill}{\parbox {7 cm}{Інформатика\hfill}}\parbox {6 cm}{\hfill Семестр 1}}\par
{\parbox {5 cm} {Навчальний предмет \hfill}\parbox {7 cm}{Програмування \hfill}\parbox {6cm}{\hfill}\par}\vspace{-15pt} \end{scriptsize}}{\center Екзаменаційний білет \No \textbf { 24 }\par}}
{%beginitem
1. Які з наведених типів не є цілими (integral), якщо такі є серед перелічених?

double, int, long*, unsigned char?
%enditem
\par
\vspace{2mm}
%beginitem
2. Скільки змінних, що є вказівниками, тут оголошено?
int n, m, *p1, p2;?
%enditem
\par
\vspace{2mm}
%beginitem
3. За наявних оголошень \code{void f(int\&); int n;} які з виразів \code{f(13), f(++n), f(\&n)} є синтаксично коректними?
%enditem
\par
\vspace{2mm}
%beginitem
4. Що таке часова складність задачі?
%enditem
\par
\vspace{2mm}
%beginitem
5. Яку часову оцінку складності має алгоритм швидкого сортування?
%enditem
\par
\vspace{2mm}
%beginitem
6. Написати програму, що запитує у користувача значення дійсних параметрів $x$ та $y$, обчислює для них значення виразу та повiдомляє введенi данi та результат обчислення.
$$\frac{\log_{2} x - 62}{(\arctg x + 0.5) (y + 35)}$$ 

%enditem
\par
%beginitem
7. Написати функцію, що повертає число з послідовності цілих (задається масивом), що має найменшу кількість різних простих дільників. Якщо таких чисел кілька, то повернути найбільше. Використовувати додаткові структури даних (у тому числі рядки) заборонено.

По завданням 2 та 3 оцінюється правильність та якість коду, а також економність обчислень.

%enditem
}
{\smallskip\smallskip \begin {scriptsize} Затверджено на засіданні кафедри теоретичної кібернетики, протокол \No 3  від   08.11.2025 р.\par\smallskip\smallskip\smallskip\smallskip
{\parbox {6.5 cm} {Зав. кафедри \dotfill Ю.В.~Крак}}\hspace {0.7cm} {\hbox to 6.5 cm { Екзаменатор \dotfill Т.О.~Карнаух}} \end{scriptsize}}
%end
\newpage
%begin
{{\begin {scriptsize}\par
{ \center  Київський національний університет імені Тараса Шевченка \par}
{\parbox {5 cm}{Освітня програма \hfill}{\parbox {7 cm}{Інформатика\hfill}}\parbox {6 cm}{\hfill Семестр 1}}\par
{\parbox {5 cm} {Навчальний предмет \hfill}\parbox {7 cm}{Програмування \hfill}\parbox {6cm}{\hfill}\par}\vspace{-15pt} \end{scriptsize}}{\center Екзаменаційний білет \No \textbf { 25 }\par}}
{%beginitem
1. Які з наведених типів не є цілими (integral), якщо такі є серед перелічених?

bool*, char, int, unsigned?
%enditem
\par
\vspace{2mm}
%beginitem
2. Скільки змінних, що є вказівниками, тут оголошено?
double x, *px1, px2, ***px3;?
%enditem
\par
\vspace{2mm}
%beginitem
3. За наявних оголошень \code{void f(char\&); char c;} які з виразів \code{f(c*c), f(c++), f(\&c)} є синтаксично коректними?
%enditem
\par
\vspace{2mm}
%beginitem
4. Що таке тип даних?
%enditem
\par
\vspace{2mm}
%beginitem
5. Яку часову оцінку складності має алгоритм сортування бульбашкою?
%enditem
\par
\vspace{2mm}
%beginitem
6. Написати програму, що запитує у користувача значення дійсних параметрів $x$ та $y$, обчислює для них значення виразу та повiдомляє введенi данi та результат обчислення.
$$\frac{\log_{2} x - 31}{(\arcsin x - 0.6) (y + 81)}$$ 

%enditem
\par
%beginitem
7. Написати функцію, що для інтервалу $[a, b)$ знаходить знаходить найбільший степінь числа 3, на який ділиться хоча б одне ціле число з цього інтервалу.

По завданням 2 та 3 оцінюється правильність та якість коду, а також економність обчислень.

%enditem
}
{\smallskip\smallskip \begin {scriptsize} Затверджено на засіданні кафедри теоретичної кібернетики, протокол \No 3  від   08.11.2025 р.\par\smallskip\smallskip\smallskip\smallskip
{\parbox {6.5 cm} {Зав. кафедри \dotfill Ю.В.~Крак}}\hspace {0.7cm} {\hbox to 6.5 cm { Екзаменатор \dotfill Т.О.~Карнаух}} \end{scriptsize}}
%end
\vspace{0.9cm}
%begin
{{\begin {scriptsize}\par
{ \center  Київський національний університет імені Тараса Шевченка \par}
{\parbox {5 cm}{Освітня програма \hfill}{\parbox {7 cm}{Інформатика\hfill}}\parbox {6 cm}{\hfill Семестр 1}}\par
{\parbox {5 cm} {Навчальний предмет \hfill}\parbox {7 cm}{Програмування \hfill}\parbox {6cm}{\hfill}\par}\vspace{-15pt} \end{scriptsize}}{\center Екзаменаційний білет \No \textbf { 26 }\par}}
{%beginitem
1. Які з наведених типів не є цілими (integral), якщо такі є серед перелічених?

int, signed char, unsigned long, string?
%enditem
\par
\vspace{2mm}
%beginitem
2. Скільки змінних, що є вказівниками, тут оголошено?
double *px, *py, **ppx;?
%enditem
\par
\vspace{2mm}
%beginitem
3. За наявних оголошень \code{void f(int*); int n;} які з виразів \code{f(n++), f(++n), f(\&n)} є синтаксично коректними?
%enditem
\par
\vspace{2mm}
%beginitem
4. Що таке С-рядок?
%enditem
\par
\vspace{2mm}
%beginitem
5. Яку часову оцінку складності має алгоритм сортування купою (він же деревом, пірамідою)?
%enditem
\par
\vspace{2mm}
%beginitem
6. Написати програму, що запитує у користувача значення дійсних параметрів $x$ та $y$, обчислює для них значення виразу та повiдомляє введенi данi та результат обчислення.
$$\frac{\pi x + 14}{(\sin x + 0.7) (y - 58)}$$ 

%enditem
\par
%beginitem
7. Написати функцію, що для послідовності цілих (задається масивом) знаходить довжину найдовшого відрізку, в якому всі сусідні числа різні. Використовувати додаткові структури даних (у тому числі рядки) заборонено.

По завданням 2 та 3 оцінюється правильність та якість коду, а також економність обчислень.

%enditem
}
{\smallskip\smallskip \begin {scriptsize} Затверджено на засіданні кафедри теоретичної кібернетики, протокол \No 3  від   08.11.2025 р.\par\smallskip\smallskip\smallskip\smallskip
{\parbox {6.5 cm} {Зав. кафедри \dotfill Ю.В.~Крак}}\hspace {0.7cm} {\hbox to 6.5 cm { Екзаменатор \dotfill Т.О.~Карнаух}} \end{scriptsize}}
%end
\newpage
%begin
{{\begin {scriptsize}\par
{ \center  Київський національний університет імені Тараса Шевченка \par}
{\parbox {5 cm}{Освітня програма \hfill}{\parbox {7 cm}{Інформатика\hfill}}\parbox {6 cm}{\hfill Семестр 1}}\par
{\parbox {5 cm} {Навчальний предмет \hfill}\parbox {7 cm}{Програмування \hfill}\parbox {6cm}{\hfill}\par}\vspace{-15pt} \end{scriptsize}}{\center Екзаменаційний білет \No \textbf { 27 }\par}}
{%beginitem
1. Які з наведених типів не є цілими (integral), якщо такі є серед перелічених?

bool, int, long, long long, unsigned char?
%enditem
\par
\vspace{2mm}
%beginitem
2. Скільки змінних, що є вказівниками, тут оголошено?
long *n1, *n2, n3, pn4;?
%enditem
\par
\vspace{2mm}
%beginitem
3. За наявних оголошень \code{void f(char*); char c;} які з виразів \code{f(c*c), f(c+=2), f(\&c)} є синтаксично коректними?
%enditem
\par
\vspace{2mm}
%beginitem
4. Чим префіксний оператор ++ відрізняється від постфіксного?
%enditem
\par
\vspace{2mm}
%beginitem
5. Яку часову оцінку складності має алгоритм сортування вставками?
%enditem
\par
\vspace{2mm}
%beginitem
6. Написати програму, що запитує у користувача значення дійсних параметрів $x$ та $y$, обчислює для них значення виразу та повiдомляє введенi данi та результат обчислення.
$$\frac{\sqrt x - 56}{(\arccos x - 0.2) (y + 12)}$$ 

%enditem
\par
%beginitem
7. Написати функцію, що для інтервалу $[a, b)$ знаходить найближчу пару чисел з цього інтервалу, що кожне з них має парну кількість різних простих дільників. Якщо таких пар кілька, то цікавить пара з найменшим добутком.

По завданням 2 та 3 оцінюється правильність та якість коду, а також економність обчислень.

%enditem
}
{\smallskip\smallskip \begin {scriptsize} Затверджено на засіданні кафедри теоретичної кібернетики, протокол \No 3  від   08.11.2025 р.\par\smallskip\smallskip\smallskip\smallskip
{\parbox {6.5 cm} {Зав. кафедри \dotfill Ю.В.~Крак}}\hspace {0.7cm} {\hbox to 6.5 cm { Екзаменатор \dotfill Т.О.~Карнаух}} \end{scriptsize}}
%end
\vspace{0.9cm}
%begin
{{\begin {scriptsize}\par
{ \center  Київський національний університет імені Тараса Шевченка \par}
{\parbox {5 cm}{Освітня програма \hfill}{\parbox {7 cm}{Інформатика\hfill}}\parbox {6 cm}{\hfill Семестр 1}}\par
{\parbox {5 cm} {Навчальний предмет \hfill}\parbox {7 cm}{Програмування \hfill}\parbox {6cm}{\hfill}\par}\vspace{-15pt} \end{scriptsize}}{\center Екзаменаційний білет \No \textbf { 28 }\par}}
{%beginitem
1. Які з наведених типів не є цілими (integral), якщо такі є серед перелічених?

double, int, long, unsigned char?
%enditem
\par
\vspace{2mm}
%beginitem
2. Скільки змінних, що є вказівниками, тут оголошено?
char *s1, pc, **s2;?
%enditem
\par
\vspace{2mm}
%beginitem
3. За наявних оголошень \code{void f(double); double x;} які з виразів \code{f(3.1), f(x+=2), f(\&x)} є синтаксично коректними?
%enditem
\par
\vspace{2mm}
%beginitem
4. Що таке часова складність алгоритму?
%enditem
\par
\vspace{2mm}
%beginitem
5. Яку часову оцінку складності має алгоритм швидкого сортування?
%enditem
\par
\vspace{2mm}
%beginitem
6. Написати програму, що запитує у користувача значення дійсних параметрів $x$ та $y$, обчислює для них значення виразу та повiдомляє введенi данi та результат обчислення.
$$\frac{\ln x + 89}{(\tg x + 0.7) (y - 94)}$$ 

%enditem
\par
%beginitem
7. Написати функцію, що для цілого числа знаходить суму його простих дільників, що входять у його розклад не більш ніж у 4му степеню. Наприклад, для числа $3^5{\cdot}7^9{\cdot}11{\cdot}23^{2}$ таких дільників два – $11$ та $23$, тому функція має повернути $34$.

По завданням 2 та 3 оцінюється правильність та якість коду, а також економність обчислень.

%enditem
}
{\smallskip\smallskip \begin {scriptsize} Затверджено на засіданні кафедри теоретичної кібернетики, протокол \No 3  від   08.11.2025 р.\par\smallskip\smallskip\smallskip\smallskip
{\parbox {6.5 cm} {Зав. кафедри \dotfill Ю.В.~Крак}}\hspace {0.7cm} {\hbox to 6.5 cm { Екзаменатор \dotfill Т.О.~Карнаух}} \end{scriptsize}}
%end
\newpage
%begin
{{\begin {scriptsize}\par
{ \center  Київський національний університет імені Тараса Шевченка \par}
{\parbox {5 cm}{Освітня програма \hfill}{\parbox {7 cm}{Інформатика\hfill}}\parbox {6 cm}{\hfill Семестр 1}}\par
{\parbox {5 cm} {Навчальний предмет \hfill}\parbox {7 cm}{Програмування \hfill}\parbox {6cm}{\hfill}\par}\vspace{-15pt} \end{scriptsize}}{\center Екзаменаційний білет \No \textbf { 29 }\par}}
{%beginitem
1. Які з наведених типів не є цілими (integral), якщо такі є серед перелічених?

char, double, int, float, unsigned?
%enditem
\par
\vspace{2mm}
%beginitem
2. Скільки змінних, що є вказівниками, тут оголошено?
long *n1, *n2, n3, pn4;?
%enditem
\par
\vspace{2mm}
%beginitem
3. За наявних оголошень \code{void f(char*); char c;} які з виразів \code{f(c*c), f(c+=2), f(\&c)} є синтаксично коректними?
%enditem
\par
\vspace{2mm}
%beginitem
4. Що означає термін ``часова складність у середньому''?
%enditem
\par
\vspace{2mm}
%beginitem
5. Яку часову оцінку складності має алгоритм швидкого сортування?
%enditem
\par
\vspace{2mm}
%beginitem
6. Написати програму, що запитує у користувача значення дійсних параметрів $x$ та $y$, обчислює для них значення виразу та повiдомляє введенi данi та результат обчислення.
$$\frac{\ x - 47}{(\arctg x + 0.4) (y + 19)}$$ 

%enditem
\par
%beginitem
7. Написати функцію, що для інтервалу $[a, b)$ знаходить найближчу пару чисел з цього інтервалу, що кожне з них має парну кількість різних простих дільників. Якщо таких пар кілька, то цікавить пара з найбільшою сумою.

По завданням 2 та 3 оцінюється правильність та якість коду, а також економність обчислень.

%enditem
}
{\smallskip\smallskip \begin {scriptsize} Затверджено на засіданні кафедри теоретичної кібернетики, протокол \No 3  від   08.11.2025 р.\par\smallskip\smallskip\smallskip\smallskip
{\parbox {6.5 cm} {Зав. кафедри \dotfill Ю.В.~Крак}}\hspace {0.7cm} {\hbox to 6.5 cm { Екзаменатор \dotfill Т.О.~Карнаух}} \end{scriptsize}}
%end
\vspace{0.9cm}
%begin
{{\begin {scriptsize}\par
{ \center  Київський національний університет імені Тараса Шевченка \par}
{\parbox {5 cm}{Освітня програма \hfill}{\parbox {7 cm}{Інформатика\hfill}}\parbox {6 cm}{\hfill Семестр 1}}\par
{\parbox {5 cm} {Навчальний предмет \hfill}\parbox {7 cm}{Програмування \hfill}\parbox {6cm}{\hfill}\par}\vspace{-15pt} \end{scriptsize}}{\center Екзаменаційний білет \No \textbf { 30 }\par}}
{%beginitem
1. Які з наведених типів не є цілими (integral), якщо такі є серед перелічених?

bool, char, int, unsigned?
%enditem
\par
\vspace{2mm}
%beginitem
2. Скільки змінних, що є вказівниками, тут оголошено?
int *pn1, pn2, pn3, pn4;?
%enditem
\par
\vspace{2mm}
%beginitem
3. За наявних оголошень \code{void f(double*); double x;} які з виразів \code{f(x*x), f(x+=2), f(\&x)} є синтаксично коректними?
%enditem
\par
\vspace{2mm}
%beginitem
4. Що таке просторова складність алгоритму?
%enditem
\par
\vspace{2mm}
%beginitem
5. Яку часову оцінку складності має алгоритм швидкого сортування?
%enditem
\par
\vspace{2mm}
%beginitem
6. Написати програму, що запитує у користувача значення дійсних параметрів $x$ та $y$, обчислює для них значення виразу та повiдомляє введенi данi та результат обчислення.
$$\frac{\log_{10} x + 69}{(\sin x - 0.3) (y - 67)}$$ 

%enditem
\par
%beginitem
7. Написати функцію, що повертає найменше число з послідовності цілих (задається масивом), що має не більше трьох різних простих дільників. За відсутності має повертати None. Використовувати додаткові структури даних (у тому числі рядки) заборонено.

По завданням 2 та 3 оцінюється правильність та якість коду, а також економність обчислень.

%enditem
}
{\smallskip\smallskip \begin {scriptsize} Затверджено на засіданні кафедри теоретичної кібернетики, протокол \No 3  від   08.11.2025 р.\par\smallskip\smallskip\smallskip\smallskip
{\parbox {6.5 cm} {Зав. кафедри \dotfill Ю.В.~Крак}}\hspace {0.7cm} {\hbox to 6.5 cm { Екзаменатор \dotfill Т.О.~Карнаух}} \end{scriptsize}}
%end
\newpage
%begin
{{\begin {scriptsize}\par
{ \center  Київський національний університет імені Тараса Шевченка \par}
{\parbox {5 cm}{Освітня програма \hfill}{\parbox {7 cm}{Інформатика\hfill}}\parbox {6 cm}{\hfill Семестр 1}}\par
{\parbox {5 cm} {Навчальний предмет \hfill}\parbox {7 cm}{Програмування \hfill}\parbox {6cm}{\hfill}\par}\vspace{-15pt} \end{scriptsize}}{\center Екзаменаційний білет \No \textbf { 31 }\par}}
{%beginitem
1. Які з наведених типів не є цілими (integral), якщо такі є серед перелічених?

double, int, long*, unsigned char?
%enditem
\par
\vspace{2mm}
%beginitem
2. Скільки змінних, що є вказівниками, тут оголошено?
int n, m, *p1, p2;?
%enditem
\par
\vspace{2mm}
%beginitem
3. За наявних оголошень \code{void f(char\&); char c;} які з виразів \code{f(c*c), f(c++), f(\&c)} є синтаксично коректними?
%enditem
\par
\vspace{2mm}
%beginitem
4. Що таке масив?
%enditem
\par
\vspace{2mm}
%beginitem
5. Яку часову оцінку складності має алгоритм сортування купою (він же деревом, пірамідою)?
%enditem
\par
\vspace{2mm}
%beginitem
6. Написати програму, що запитує у користувача значення дійсних параметрів $x$ та $y$, обчислює для них значення виразу та повiдомляє введенi данi та результат обчислення.
$$\frac{\log_{10} x - 40}{(\cos x - 0.1) (y - 25)}$$ 

%enditem
\par
%beginitem
7. Написати функцію, що для інтервалу $[a, b)$ знаходить найближчу пару чисел з цього інтервалу, що кожне з них має непарну кількість різних простих дільників. Якщо таких пар кілька, то цікавить пара з найменшим добутком.

По завданням 2 та 3 оцінюється правильність та якість коду, а також економність обчислень.

%enditem
}
{\smallskip\smallskip \begin {scriptsize} Затверджено на засіданні кафедри теоретичної кібернетики, протокол \No 3  від   08.11.2025 р.\par\smallskip\smallskip\smallskip\smallskip
{\parbox {6.5 cm} {Зав. кафедри \dotfill Ю.В.~Крак}}\hspace {0.7cm} {\hbox to 6.5 cm { Екзаменатор \dotfill Т.О.~Карнаух}} \end{scriptsize}}
%end
\vspace{0.9cm}
%begin
{{\begin {scriptsize}\par
{ \center  Київський національний університет імені Тараса Шевченка \par}
{\parbox {5 cm}{Освітня програма \hfill}{\parbox {7 cm}{Інформатика\hfill}}\parbox {6 cm}{\hfill Семестр 1}}\par
{\parbox {5 cm} {Навчальний предмет \hfill}\parbox {7 cm}{Програмування \hfill}\parbox {6cm}{\hfill}\par}\vspace{-15pt} \end{scriptsize}}{\center Екзаменаційний білет \No \textbf { 32 }\par}}
{%beginitem
1. Які з наведених типів не є цілими (integral), якщо такі є серед перелічених?

double, int, long, unsigned char?
%enditem
\par
\vspace{2mm}
%beginitem
2. Скільки змінних, що є вказівниками, тут оголошено?
char *s1, pc, **s2;?
%enditem
\par
\vspace{2mm}
%beginitem
3. За наявних оголошень \code{void f(char); char c;} які з виразів \code{f(c*c), f(c+=2), f(\&c)} є синтаксично коректними?
%enditem
\par
\vspace{2mm}
%beginitem
4. Які параметри може приймати функція main? Що вони позначають та яких вони типів?
%enditem
\par
\vspace{2mm}
%beginitem
5. Яку часову оцінку складності має алгоритм сортування злиттям?
%enditem
\par
\vspace{2mm}
%beginitem
6. Написати програму, що запитує у користувача значення дійсних параметрів $x$ та $y$, обчислює для них значення виразу та повiдомляє введенi данi та результат обчислення.
$$\frac{\ x + 80}{(\ctg x + 0.2) (y + 39)}$$ 

%enditem
\par
%beginitem
7. Написати функцію, що для послідовності цілих (задається масивом) знаходить найдовший відрізок, в якому всі сусідні числа відрізняються якнайменш на 3. Якщо таких відрізків кілька, то повернути перший (як індекс початку та довжину). Використовувати додаткові структури даних (у тому числі рядки) заборонено.

По завданням 2 та 3 оцінюється правильність та якість коду, а також економність обчислень.

%enditem
}
{\smallskip\smallskip \begin {scriptsize} Затверджено на засіданні кафедри теоретичної кібернетики, протокол \No 3  від   08.11.2025 р.\par\smallskip\smallskip\smallskip\smallskip
{\parbox {6.5 cm} {Зав. кафедри \dotfill Ю.В.~Крак}}\hspace {0.7cm} {\hbox to 6.5 cm { Екзаменатор \dotfill Т.О.~Карнаух}} \end{scriptsize}}
%end
\newpage
%begin
{{\begin {scriptsize}\par
{ \center  Київський національний університет імені Тараса Шевченка \par}
{\parbox {5 cm}{Освітня програма \hfill}{\parbox {7 cm}{Інформатика\hfill}}\parbox {6 cm}{\hfill Семестр 1}}\par
{\parbox {5 cm} {Навчальний предмет \hfill}\parbox {7 cm}{Програмування \hfill}\parbox {6cm}{\hfill}\par}\vspace{-15pt} \end{scriptsize}}{\center Екзаменаційний білет \No \textbf { 33 }\par}}
{%beginitem
1. Які з наведених типів не є цілими (integral), якщо такі є серед перелічених?

bool, int, long, long long, unsigned char?
%enditem
\par
\vspace{2mm}
%beginitem
2. Скільки змінних, що є вказівниками, тут оголошено?
long pointer1, pointer2, pointer3, pointer4;?
%enditem
\par
\vspace{2mm}
%beginitem
3. За наявних оголошень \code{void f(int); int n;} які з виразів \code{f(3), f(++n), f(\&n)} є синтаксично коректними?
%enditem
\par
\vspace{2mm}
%beginitem
4. У чому відмінність літералів від констант?
%enditem
\par
\vspace{2mm}
%beginitem
5. Яку часову оцінку складності має алгоритм сортування бульбашкою?
%enditem
\par
\vspace{2mm}
%beginitem
6. Написати програму, що запитує у користувача значення дійсних параметрів $x$ та $y$, обчислює для них значення виразу та повiдомляє введенi данi та результат обчислення.
$$\frac{\ln x - 27}{(\tg x - 0.8) (y - 20)}$$ 

%enditem
\par
%beginitem
7. Написати функцію, що знаходить найбільше число з послідовності цілих (задається масивом), у десятковому запису якого нема цифри 3. Повертає індекс останнього входження цього числа або $-1$ (за відсутності). Використовувати додаткові структури даних (у тому числі рядки) заборонено.

По завданням 2 та 3 оцінюється правильність та якість коду, а також економність обчислень.

%enditem
}
{\smallskip\smallskip \begin {scriptsize} Затверджено на засіданні кафедри теоретичної кібернетики, протокол \No 3  від   08.11.2025 р.\par\smallskip\smallskip\smallskip\smallskip
{\parbox {6.5 cm} {Зав. кафедри \dotfill Ю.В.~Крак}}\hspace {0.7cm} {\hbox to 6.5 cm { Екзаменатор \dotfill Т.О.~Карнаух}} \end{scriptsize}}
%end
\vspace{0.9cm}
%begin
{{\begin {scriptsize}\par
{ \center  Київський національний університет імені Тараса Шевченка \par}
{\parbox {5 cm}{Освітня програма \hfill}{\parbox {7 cm}{Інформатика\hfill}}\parbox {6 cm}{\hfill Семестр 1}}\par
{\parbox {5 cm} {Навчальний предмет \hfill}\parbox {7 cm}{Програмування \hfill}\parbox {6cm}{\hfill}\par}\vspace{-15pt} \end{scriptsize}}{\center Екзаменаційний білет \No \textbf { 34 }\par}}
{%beginitem
1. Які з наведених типів не є цілими (integral), якщо такі є серед перелічених?

char, double, int, float, unsigned?
%enditem
\par
\vspace{2mm}
%beginitem
2. Скільки змінних, що є вказівниками, тут оголошено?
double x, *px1, px2, ***px3;?
%enditem
\par
\vspace{2mm}
%beginitem
3. За наявних оголошень \code{void f(int\&); int n;} які з виразів \code{f(13), f(++n), f(\&n)} є синтаксично коректними?
%enditem
\par
\vspace{2mm}
%beginitem
4. Що таке просторова складність алгоритму?
%enditem
\par
\vspace{2mm}
%beginitem
5. Яку часову оцінку складності має алгоритм сортування вставками?
%enditem
\par
\vspace{2mm}
%beginitem
6. Написати програму, що запитує у користувача значення дійсних параметрів $x$ та $y$, обчислює для них значення виразу та повiдомляє введенi данi та результат обчислення.
$$\frac{\pi x + 77}{(\arcsin x + 0.9) (y + 46)}$$ 

%enditem
\par
%beginitem
7. Написати функцію, що для цілого числа знаходить кількість його простих дільників, що входять у його розклад не більш ніж у 5му степеню. Наприклад, для числа $3^8{\cdot}7^9{\cdot}11{\cdot}23^{2}$ таких дільників два – $11$ та $23$, тому функція має повернути $2$.

По завданням 2 та 3 оцінюється правильність та якість коду, а також економність обчислень.

%enditem
}
{\smallskip\smallskip \begin {scriptsize} Затверджено на засіданні кафедри теоретичної кібернетики, протокол \No 3  від   08.11.2025 р.\par\smallskip\smallskip\smallskip\smallskip
{\parbox {6.5 cm} {Зав. кафедри \dotfill Ю.В.~Крак}}\hspace {0.7cm} {\hbox to 6.5 cm { Екзаменатор \dotfill Т.О.~Карнаух}} \end{scriptsize}}
%end
\newpage
%begin
{{\begin {scriptsize}\par
{ \center  Київський національний університет імені Тараса Шевченка \par}
{\parbox {5 cm}{Освітня програма \hfill}{\parbox {7 cm}{Інформатика\hfill}}\parbox {6 cm}{\hfill Семестр 1}}\par
{\parbox {5 cm} {Навчальний предмет \hfill}\parbox {7 cm}{Програмування \hfill}\parbox {6cm}{\hfill}\par}\vspace{-15pt} \end{scriptsize}}{\center Екзаменаційний білет \No \textbf { 35 }\par}}
{%beginitem
1. Які з наведених типів не є цілими (integral), якщо такі є серед перелічених?

int, signed char, unsigned long, string?
%enditem
\par
\vspace{2mm}
%beginitem
2. Скільки змінних, що є вказівниками, тут оголошено?
double *px, *py, **ppx;?
%enditem
\par
\vspace{2mm}
%beginitem
3. За наявних оголошень \code{void f(double\&); double x;} які з виразів \code{f(x*x), f(x+=2), f(\&x)} є синтаксично коректними?
%enditem
\par
\vspace{2mm}
%beginitem
4. Що таке часова складність алгоритму?
%enditem
\par
\vspace{2mm}
%beginitem
5. Яку часову оцінку складності має алгоритм сортування вибором?
%enditem
\par
\vspace{2mm}
%beginitem
6. Написати програму, що запитує у користувача значення дійсних параметрів $x$ та $y$, обчислює для них значення виразу та повiдомляє введенi данi та результат обчислення.
$$\frac{\log_{2} x - 84}{(\arccos x + 0.6) (y - 49)}$$ 

%enditem
\par
%beginitem
7. Написати функцію, що для цілого числа знаходить кількість його простих дільників, що входять у його розклад рівно у другому степеню. Наприклад, для числа $3^2{\cdot}7^9{\cdot}11{\cdot}23^{2}$ таких дільників два – $3$ та $23$, тому функція має повернути $2$.

По завданням 2 та 3 оцінюється правильність та якість коду, а також економність обчислень.

%enditem
}
{\smallskip\smallskip \begin {scriptsize} Затверджено на засіданні кафедри теоретичної кібернетики, протокол \No 3  від   08.11.2025 р.\par\smallskip\smallskip\smallskip\smallskip
{\parbox {6.5 cm} {Зав. кафедри \dotfill Ю.В.~Крак}}\hspace {0.7cm} {\hbox to 6.5 cm { Екзаменатор \dotfill Т.О.~Карнаух}} \end{scriptsize}}
%end
\vspace{0.9cm}
%begin
{{\begin {scriptsize}\par
{ \center  Київський національний університет імені Тараса Шевченка \par}
{\parbox {5 cm}{Освітня програма \hfill}{\parbox {7 cm}{Інформатика\hfill}}\parbox {6 cm}{\hfill Семестр 1}}\par
{\parbox {5 cm} {Навчальний предмет \hfill}\parbox {7 cm}{Програмування \hfill}\parbox {6cm}{\hfill}\par}\vspace{-15pt} \end{scriptsize}}{\center Екзаменаційний білет \No \textbf { 36 }\par}}
{%beginitem
1. Які з наведених типів не є цілими (integral), якщо такі є серед перелічених?

bool, int, long, long long, unsigned char?
%enditem
\par
\vspace{2mm}
%beginitem
2. Скільки змінних, що є вказівниками, тут оголошено?
int *pn1, pn2, pn3, pn4;?
%enditem
\par
\vspace{2mm}
%beginitem
3. За наявних оголошень \code{void f(int*); int n;} які з виразів \code{f(n++), f(++n), f(\&n)} є синтаксично коректними?
%enditem
\par
\vspace{2mm}
%beginitem
4. Чим префіксний оператор ++ відрізняється від постфіксного?
%enditem
\par
\vspace{2mm}
%beginitem
5. Яку часову оцінку складності має алгоритм швидкого сортування?
%enditem
\par
\vspace{2mm}
%beginitem
6. Написати програму, що запитує у користувача значення дійсних параметрів $x$ та $y$, обчислює для них значення виразу та повiдомляє введенi данi та результат обчислення.
$$\frac{\sqrt x + 11}{(\arccos x - 0.5) (y + 51)}$$ 

%enditem
\par
%beginitem
7. Написати функцію, що для інтервалу $[a, b)$ знаходить знаходить найбільший степінь числа 3, на який ділиться хоча б одне ціле число з цього інтервалу.

По завданням 2 та 3 оцінюється правильність та якість коду, а також економність обчислень.

%enditem
}
{\smallskip\smallskip \begin {scriptsize} Затверджено на засіданні кафедри теоретичної кібернетики, протокол \No 3  від   08.11.2025 р.\par\smallskip\smallskip\smallskip\smallskip
{\parbox {6.5 cm} {Зав. кафедри \dotfill Ю.В.~Крак}}\hspace {0.7cm} {\hbox to 6.5 cm { Екзаменатор \dotfill Т.О.~Карнаух}} \end{scriptsize}}
%end
\newpage
%begin
{{\begin {scriptsize}\par
{ \center  Київський національний університет імені Тараса Шевченка \par}
{\parbox {5 cm}{Освітня програма \hfill}{\parbox {7 cm}{Інформатика\hfill}}\parbox {6 cm}{\hfill Семестр 1}}\par
{\parbox {5 cm} {Навчальний предмет \hfill}\parbox {7 cm}{Програмування \hfill}\parbox {6cm}{\hfill}\par}\vspace{-15pt} \end{scriptsize}}{\center Екзаменаційний білет \No \textbf { 37 }\par}}
{%beginitem
1. Які з наведених типів не є цілими (integral), якщо такі є серед перелічених?

bool*, char, int, unsigned?
%enditem
\par
\vspace{2mm}
%beginitem
2. Скільки змінних, що є вказівниками, тут оголошено?
int n, m, *p1, p2;?
%enditem
\par
\vspace{2mm}
%beginitem
3. За наявних оголошень \code{void f(int*); int n;} які з виразів \code{f(n++), f(++n), f(\&n)} є синтаксично коректними?
%enditem
\par
\vspace{2mm}
%beginitem
4. Які параметри може приймати функція main? Що вони позначають та яких вони типів?
%enditem
\par
\vspace{2mm}
%beginitem
5. Яку часову оцінку складності має алгоритм швидкого сортування?
%enditem
\par
\vspace{2mm}
%beginitem
6. Написати програму, що запитує у користувача значення дійсних параметрів $x$ та $y$, обчислює для них значення виразу та повiдомляє введенi данi та результат обчислення.
$$\frac{\log_{2} x - 57}{(\arcsin x - 0.6) (y - 50)}$$ 

%enditem
\par
%beginitem
7. Написати функцію, що знаходить найбільше число з послідовності цілих (задається масивом), у десятковому запису якого нема цифри 3. Повертає індекс останнього входження цього числа або $-1$ (за відсутності). Використовувати додаткові структури даних (у тому числі рядки) заборонено.

По завданням 2 та 3 оцінюється правильність та якість коду, а також економність обчислень.

%enditem
}
{\smallskip\smallskip \begin {scriptsize} Затверджено на засіданні кафедри теоретичної кібернетики, протокол \No 3  від   08.11.2025 р.\par\smallskip\smallskip\smallskip\smallskip
{\parbox {6.5 cm} {Зав. кафедри \dotfill Ю.В.~Крак}}\hspace {0.7cm} {\hbox to 6.5 cm { Екзаменатор \dotfill Т.О.~Карнаух}} \end{scriptsize}}
%end
\vspace{0.9cm}
%begin
{{\begin {scriptsize}\par
{ \center  Київський національний університет імені Тараса Шевченка \par}
{\parbox {5 cm}{Освітня програма \hfill}{\parbox {7 cm}{Інформатика\hfill}}\parbox {6 cm}{\hfill Семестр 1}}\par
{\parbox {5 cm} {Навчальний предмет \hfill}\parbox {7 cm}{Програмування \hfill}\parbox {6cm}{\hfill}\par}\vspace{-15pt} \end{scriptsize}}{\center Екзаменаційний білет \No \textbf { 38 }\par}}
{%beginitem
1. Які з наведених типів не є цілими (integral), якщо такі є серед перелічених?

char, double, int, float, unsigned*?
%enditem
\par
\vspace{2mm}
%beginitem
2. Скільки змінних, що є вказівниками, тут оголошено?
long pointer1, pointer2, pointer3, pointer4;?
%enditem
\par
\vspace{2mm}
%beginitem
3. За наявних оголошень \code{void f(char\&); char c;} які з виразів \code{f(c*c), f(c++), f(\&c)} є синтаксично коректними?
%enditem
\par
\vspace{2mm}
%beginitem
4. Що таке масив?
%enditem
\par
\vspace{2mm}
%beginitem
5. Яку часову оцінку складності має алгоритм сортування купою (він же деревом, пірамідою)?
%enditem
\par
\vspace{2mm}
%beginitem
6. Написати програму, що запитує у користувача значення дійсних параметрів $x$ та $y$, обчислює для них значення виразу та повiдомляє введенi данi та результат обчислення.
$$\frac{\pi x + 97}{(\arctg x + 0.4) (y + 45)}$$ 

%enditem
\par
%beginitem
7. Написати функцію, що для послідовності цілих (задається масивом) знаходить довжину найдовшого відрізку, в якому всі сусідні числа різні. Використовувати додаткові структури даних (у тому числі рядки) заборонено.

По завданням 2 та 3 оцінюється правильність та якість коду, а також економність обчислень.

%enditem
}
{\smallskip\smallskip \begin {scriptsize} Затверджено на засіданні кафедри теоретичної кібернетики, протокол \No 3  від   08.11.2025 р.\par\smallskip\smallskip\smallskip\smallskip
{\parbox {6.5 cm} {Зав. кафедри \dotfill Ю.В.~Крак}}\hspace {0.7cm} {\hbox to 6.5 cm { Екзаменатор \dotfill Т.О.~Карнаух}} \end{scriptsize}}
%end
\newpage
%begin
{{\begin {scriptsize}\par
{ \center  Київський національний університет імені Тараса Шевченка \par}
{\parbox {5 cm}{Освітня програма \hfill}{\parbox {7 cm}{Інформатика\hfill}}\parbox {6 cm}{\hfill Семестр 1}}\par
{\parbox {5 cm} {Навчальний предмет \hfill}\parbox {7 cm}{Програмування \hfill}\parbox {6cm}{\hfill}\par}\vspace{-15pt} \end{scriptsize}}{\center Екзаменаційний білет \No \textbf { 39 }\par}}
{%beginitem
1. Які з наведених типів не є цілими (integral), якщо такі є серед перелічених?

int*, signed char, unsigned long, string?
%enditem
\par
\vspace{2mm}
%beginitem
2. Скільки змінних, що є вказівниками, тут оголошено?
long *n1, *n2, n3, pn4;?
%enditem
\par
\vspace{2mm}
%beginitem
3. За наявних оголошень \code{void f(char*); char c;} які з виразів \code{f(c*c), f(c+=2), f(\&c)} є синтаксично коректними?
%enditem
\par
\vspace{2mm}
%beginitem
4. Що означає термін ``часова складність у середньому''?
%enditem
\par
\vspace{2mm}
%beginitem
5. Яку часову оцінку складності має алгоритм сортування вставками?
%enditem
\par
\vspace{2mm}
%beginitem
6. Написати програму, що запитує у користувача значення дійсних параметрів $x$ та $y$, обчислює для них значення виразу та повiдомляє введенi данi та результат обчислення.
$$\frac{\ln x - 23}{(\tg x + 0.7) (y - 82)}$$ 

%enditem
\par
%beginitem
7. Написати функцію, що повертає число з послідовності цілих (задається масивом), що має найбільшу кількість різних простих дільників. Якщо таких чисел кілька, то повернути найменше. Використовувати додаткові структури даних (у тому числі рядки) заборонено.

По завданням 2 та 3 оцінюється правильність та якість коду, а також економність обчислень.

%enditem
}
{\smallskip\smallskip \begin {scriptsize} Затверджено на засіданні кафедри теоретичної кібернетики, протокол \No 3  від   08.11.2025 р.\par\smallskip\smallskip\smallskip\smallskip
{\parbox {6.5 cm} {Зав. кафедри \dotfill Ю.В.~Крак}}\hspace {0.7cm} {\hbox to 6.5 cm { Екзаменатор \dotfill Т.О.~Карнаух}} \end{scriptsize}}
%end
\vspace{0.9cm}
%begin
{{\begin {scriptsize}\par
{ \center  Київський національний університет імені Тараса Шевченка \par}
{\parbox {5 cm}{Освітня програма \hfill}{\parbox {7 cm}{Інформатика\hfill}}\parbox {6 cm}{\hfill Семестр 1}}\par
{\parbox {5 cm} {Навчальний предмет \hfill}\parbox {7 cm}{Програмування \hfill}\parbox {6cm}{\hfill}\par}\vspace{-15pt} \end{scriptsize}}{\center Екзаменаційний білет \No \textbf { 40 }\par}}
{%beginitem
1. Які з наведених типів не є цілими (integral), якщо такі є серед перелічених?

bool, char, int, unsigned?
%enditem
\par
\vspace{2mm}
%beginitem
2. Скільки змінних, що є вказівниками, тут оголошено?
char *s1, pc, **s2;?
%enditem
\par
\vspace{2mm}
%beginitem
3. За наявних оголошень \code{void f(int\&); int n;} які з виразів \code{f(13), f(++n), f(\&n)} є синтаксично коректними?
%enditem
\par
\vspace{2mm}
%beginitem
4. Що таке С-рядок?
%enditem
\par
\vspace{2mm}
%beginitem
5. Яку часову оцінку складності має алгоритм сортування бульбашкою?
%enditem
\par
\vspace{2mm}
%beginitem
6. Написати програму, що запитує у користувача значення дійсних параметрів $x$ та $y$, обчислює для них значення виразу та повiдомляє введенi данi та результат обчислення.
$$\frac{\ x + 95}{(\ctg x - 0.5) (y + 71)}$$ 

%enditem
\par
%beginitem
7. Написати функцію, що для послідовності цілих (задається масивом) знаходить найдовший відрізок, в якому всі сусідні числа різні. Якщо таких відрізків кілька, то повернути останній (як індекс початку та довжину). Використовувати додаткові структури даних (у тому числі рядки) заборонено.

По завданням 2 та 3 оцінюється правильність та якість коду, а також економність обчислень.

%enditem
}
{\smallskip\smallskip \begin {scriptsize} Затверджено на засіданні кафедри теоретичної кібернетики, протокол \No 3  від   08.11.2025 р.\par\smallskip\smallskip\smallskip\smallskip
{\parbox {6.5 cm} {Зав. кафедри \dotfill Ю.В.~Крак}}\hspace {0.7cm} {\hbox to 6.5 cm { Екзаменатор \dotfill Т.О.~Карнаух}} \end{scriptsize}}
%end
\newpage
%begin
{{\begin {scriptsize}\par
{ \center  Київський національний університет імені Тараса Шевченка \par}
{\parbox {5 cm}{Освітня програма \hfill}{\parbox {7 cm}{Інформатика\hfill}}\parbox {6 cm}{\hfill Семестр 1}}\par
{\parbox {5 cm} {Навчальний предмет \hfill}\parbox {7 cm}{Програмування \hfill}\parbox {6cm}{\hfill}\par}\vspace{-15pt} \end{scriptsize}}{\center Екзаменаційний білет \No \textbf { 41 }\par}}
{%beginitem
1. Які з наведених типів не є цілими (integral), якщо такі є серед перелічених?

int, signed char, unsigned long, string?
%enditem
\par
\vspace{2mm}
%beginitem
2. Скільки змінних, що є вказівниками, тут оголошено?
double x, *px1, px2, ***px3;?
%enditem
\par
\vspace{2mm}
%beginitem
3. За наявних оголошень \code{void f(double\&); double x;} які з виразів \code{f(x*x), f(x+=2), f(\&x)} є синтаксично коректними?
%enditem
\par
\vspace{2mm}
%beginitem
4. Що таке просторова складність задачі?
%enditem
\par
\vspace{2mm}
%beginitem
5. Яку часову оцінку складності має алгоритм сортування злиттям?
%enditem
\par
\vspace{2mm}
%beginitem
6. Написати програму, що запитує у користувача значення дійсних параметрів $x$ та $y$, обчислює для них значення виразу та повiдомляє введенi данi та результат обчислення.
$$\frac{\sqrt x - 66}{(\sin x - 0.3) (y + 52)}$$ 

%enditem
\par
%beginitem
7. Написати функцію, що повертає найбільше число з послідовності цілих (задається масивом), у сімковому запису якого нема цифри 2. За відсутності має повертати None. Використовувати додаткові структури даних (у тому числі рядки) заборонено.

По завданням 2 та 3 оцінюється правильність та якість коду, а також економність обчислень.

%enditem
}
{\smallskip\smallskip \begin {scriptsize} Затверджено на засіданні кафедри теоретичної кібернетики, протокол \No 3  від   08.11.2025 р.\par\smallskip\smallskip\smallskip\smallskip
{\parbox {6.5 cm} {Зав. кафедри \dotfill Ю.В.~Крак}}\hspace {0.7cm} {\hbox to 6.5 cm { Екзаменатор \dotfill Т.О.~Карнаух}} \end{scriptsize}}
%end
\vspace{0.9cm}
%begin
{{\begin {scriptsize}\par
{ \center  Київський національний університет імені Тараса Шевченка \par}
{\parbox {5 cm}{Освітня програма \hfill}{\parbox {7 cm}{Інформатика\hfill}}\parbox {6 cm}{\hfill Семестр 1}}\par
{\parbox {5 cm} {Навчальний предмет \hfill}\parbox {7 cm}{Програмування \hfill}\parbox {6cm}{\hfill}\par}\vspace{-15pt} \end{scriptsize}}{\center Екзаменаційний білет \No \textbf { 42 }\par}}
{%beginitem
1. Які з наведених типів не є цілими (integral), якщо такі є серед перелічених?

char, double, int, float, unsigned?
%enditem
\par
\vspace{2mm}
%beginitem
2. Скільки змінних, що є вказівниками, тут оголошено?
double *px, *py, **ppx;?
%enditem
\par
\vspace{2mm}
%beginitem
3. За наявних оголошень \code{void f(double*); double x;} які з виразів \code{f(x*x), f(x+=2), f(\&x)} є синтаксично коректними?
%enditem
\par
\vspace{2mm}
%beginitem
4. Що таке часова складність задачі?
%enditem
\par
\vspace{2mm}
%beginitem
5. Яку часову оцінку складності має алгоритм сортування вибором?
%enditem
\par
\vspace{2mm}
%beginitem
6. Написати програму, що запитує у користувача значення дійсних параметрів $x$ та $y$, обчислює для них значення виразу та повiдомляє введенi данi та результат обчислення.
$$\frac{\log_{10} x + 15}{(\cos x + 0.8) (y - 16)}$$ 

%enditem
\par
%beginitem
7. Написати функцію, що для інтервалу $[a, b)$ знаходить найближчу пару чисел з цього інтервалу, що мають від трьох до п'яти (включно) простих дільників.

По завданням 2 та 3 оцінюється правильність та якість коду, а також економність обчислень.

%enditem
}
{\smallskip\smallskip \begin {scriptsize} Затверджено на засіданні кафедри теоретичної кібернетики, протокол \No 3  від   08.11.2025 р.\par\smallskip\smallskip\smallskip\smallskip
{\parbox {6.5 cm} {Зав. кафедри \dotfill Ю.В.~Крак}}\hspace {0.7cm} {\hbox to 6.5 cm { Екзаменатор \dotfill Т.О.~Карнаух}} \end{scriptsize}}
%end
\newpage
%begin
{{\begin {scriptsize}\par
{ \center  Київський національний університет імені Тараса Шевченка \par}
{\parbox {5 cm}{Освітня програма \hfill}{\parbox {7 cm}{Інформатика\hfill}}\parbox {6 cm}{\hfill Семестр 1}}\par
{\parbox {5 cm} {Навчальний предмет \hfill}\parbox {7 cm}{Програмування \hfill}\parbox {6cm}{\hfill}\par}\vspace{-15pt} \end{scriptsize}}{\center Екзаменаційний білет \No \textbf { 43 }\par}}
{%beginitem
1. Які з наведених типів не є цілими (integral), якщо такі є серед перелічених?

bool*, char, int, unsigned?
%enditem
\par
\vspace{2mm}
%beginitem
2. Скільки змінних, що є вказівниками, тут оголошено?
double x, *px1, px2, ***px3;?
%enditem
\par
\vspace{2mm}
%beginitem
3. За наявних оголошень \code{void f(double); double x;} які з виразів \code{f(3.1), f(x+=2), f(\&x)} є синтаксично коректними?
%enditem
\par
\vspace{2mm}
%beginitem
4. Що таке тип даних?
%enditem
\par
\vspace{2mm}
%beginitem
5. Яку часову оцінку складності має алгоритм сортування вставками?
%enditem
\par
\vspace{2mm}
%beginitem
6. Написати програму, що запитує у користувача значення дійсних параметрів $x$ та $y$, обчислює для них значення виразу та повiдомляє введенi данi та результат обчислення.
$$\frac{\log_{10} x - 72}{(\tg x + 0.9) (y - 42)}$$ 

%enditem
\par
%beginitem
7. Написати функцію, що для цілого числа знаходить кількість простих дільників, які входять у його розклад у найбільшому степені. Наприклад, для числа $3^5{\cdot}7^9{\cdot}11{\cdot}23^{9}$ таких дільників два – $23$ та $7$, тому функція має повернути $2$.

По завданням 2 та 3 оцінюється правильність та якість коду, а також економність обчислень.

%enditem
}
{\smallskip\smallskip \begin {scriptsize} Затверджено на засіданні кафедри теоретичної кібернетики, протокол \No 3  від   08.11.2025 р.\par\smallskip\smallskip\smallskip\smallskip
{\parbox {6.5 cm} {Зав. кафедри \dotfill Ю.В.~Крак}}\hspace {0.7cm} {\hbox to 6.5 cm { Екзаменатор \dotfill Т.О.~Карнаух}} \end{scriptsize}}
%end
\vspace{0.9cm}
%begin
{{\begin {scriptsize}\par
{ \center  Київський національний університет імені Тараса Шевченка \par}
{\parbox {5 cm}{Освітня програма \hfill}{\parbox {7 cm}{Інформатика\hfill}}\parbox {6 cm}{\hfill Семестр 1}}\par
{\parbox {5 cm} {Навчальний предмет \hfill}\parbox {7 cm}{Програмування \hfill}\parbox {6cm}{\hfill}\par}\vspace{-15pt} \end{scriptsize}}{\center Екзаменаційний білет \No \textbf { 44 }\par}}
{%beginitem
1. Які з наведених типів не є цілими (integral), якщо такі є серед перелічених?

bool, int, long, long long, unsigned char?
%enditem
\par
\vspace{2mm}
%beginitem
2. Скільки змінних, що є вказівниками, тут оголошено?
double *px, *py, **ppx;?
%enditem
\par
\vspace{2mm}
%beginitem
3. За наявних оголошень \code{void f(int); int n;} які з виразів \code{f(3), f(++n), f(\&n)} є синтаксично коректними?
%enditem
\par
\vspace{2mm}
%beginitem
4. У чому відмінність літералів від констант?
%enditem
\par
\vspace{2mm}
%beginitem
5. Яку часову оцінку складності має алгоритм сортування бульбашкою?
%enditem
\par
\vspace{2mm}
%beginitem
6. Написати програму, що запитує у користувача значення дійсних параметрів $x$ та $y$, обчислює для них значення виразу та повiдомляє введенi данi та результат обчислення.
$$\frac{\sqrt x + 91}{(\sin x - 0.1) (y + 74)}$$ 

%enditem
\par
%beginitem
7. Написати функцію, що для інтервалу $[a, b)$ знаходить кількість чисел, квадрати яких діляться на суму квадратів їх цифр. (Розглядаємо десятковий запис.)

По завданням 2 та 3 оцінюється правильність та якість коду, а також економність обчислень.

%enditem
}
{\smallskip\smallskip \begin {scriptsize} Затверджено на засіданні кафедри теоретичної кібернетики, протокол \No 3  від   08.11.2025 р.\par\smallskip\smallskip\smallskip\smallskip
{\parbox {6.5 cm} {Зав. кафедри \dotfill Ю.В.~Крак}}\hspace {0.7cm} {\hbox to 6.5 cm { Екзаменатор \dotfill Т.О.~Карнаух}} \end{scriptsize}}
%end
\newpage
%begin
{{\begin {scriptsize}\par
{ \center  Київський національний університет імені Тараса Шевченка \par}
{\parbox {5 cm}{Освітня програма \hfill}{\parbox {7 cm}{Інформатика\hfill}}\parbox {6 cm}{\hfill Семестр 1}}\par
{\parbox {5 cm} {Навчальний предмет \hfill}\parbox {7 cm}{Програмування \hfill}\parbox {6cm}{\hfill}\par}\vspace{-15pt} \end{scriptsize}}{\center Екзаменаційний білет \No \textbf { 45 }\par}}
{%beginitem
1. Які з наведених типів не є цілими (integral), якщо такі є серед перелічених?

char, double, int, float, unsigned*?
%enditem
\par
\vspace{2mm}
%beginitem
2. Скільки змінних, що є вказівниками, тут оголошено?
int *pn1, pn2, pn3, pn4;?
%enditem
\par
\vspace{2mm}
%beginitem
3. За наявних оголошень \code{void f(char); char c;} які з виразів \code{f(c*c), f(c+=2), f(\&c)} є синтаксично коректними?
%enditem
\par
\vspace{2mm}
%beginitem
4. Які параметри може приймати функція main? Що вони позначають та яких вони типів?
%enditem
\par
\vspace{2mm}
%beginitem
5. Яку часову оцінку складності має алгоритм швидкого сортування?
%enditem
\par
\vspace{2mm}
%beginitem
6. Написати програму, що запитує у користувача значення дійсних параметрів $x$ та $y$, обчислює для них значення виразу та повiдомляє введенi данi та результат обчислення.
$$\frac{\log_{2} x + 56}{(\arcsin x - 0.2) (y - 29)}$$ 

%enditem
\par
%beginitem
7. Написати функцію, що для цілого числа знаходить кількість його простих дільників, що входять у його розклад не більш ніж у 5му степеню. Наприклад, для числа $3^8{\cdot}7^9{\cdot}11{\cdot}23^{2}$ таких дільників два – $11$ та $23$, тому функція має повернути $2$.

По завданням 2 та 3 оцінюється правильність та якість коду, а також економність обчислень.

%enditem
}
{\smallskip\smallskip \begin {scriptsize} Затверджено на засіданні кафедри теоретичної кібернетики, протокол \No 3  від   08.11.2025 р.\par\smallskip\smallskip\smallskip\smallskip
{\parbox {6.5 cm} {Зав. кафедри \dotfill Ю.В.~Крак}}\hspace {0.7cm} {\hbox to 6.5 cm { Екзаменатор \dotfill Т.О.~Карнаух}} \end{scriptsize}}
%end
\vspace{0.9cm}
%begin
{{\begin {scriptsize}\par
{ \center  Київський національний університет імені Тараса Шевченка \par}
{\parbox {5 cm}{Освітня програма \hfill}{\parbox {7 cm}{Інформатика\hfill}}\parbox {6 cm}{\hfill Семестр 1}}\par
{\parbox {5 cm} {Навчальний предмет \hfill}\parbox {7 cm}{Програмування \hfill}\parbox {6cm}{\hfill}\par}\vspace{-15pt} \end{scriptsize}}{\center Екзаменаційний білет \No \textbf { 46 }\par}}
{%beginitem
1. Які з наведених типів не є цілими (integral), якщо такі є серед перелічених?

double, int, long, unsigned char?
%enditem
\par
\vspace{2mm}
%beginitem
2. Скільки змінних, що є вказівниками, тут оголошено?
long *n1, *n2, n3, pn4;?
%enditem
\par
\vspace{2mm}
%beginitem
3. За наявних оголошень \code{void f(char*); char c;} які з виразів \code{f(c*c), f(c+=2), f(\&c)} є синтаксично коректними?
%enditem
\par
\vspace{2mm}
%beginitem
4. Що означає термін ``часова складність у середньому''?
%enditem
\par
\vspace{2mm}
%beginitem
5. Яку часову оцінку складності має алгоритм сортування вибором?
%enditem
\par
\vspace{2mm}
%beginitem
6. Написати програму, що запитує у користувача значення дійсних параметрів $x$ та $y$, обчислює для них значення виразу та повiдомляє введенi данi та результат обчислення.
$$\frac{\ x - 39}{(\ctg x + 0.2) (y + 58)}$$ 

%enditem
\par
%beginitem
7. Написати функцію, що для інтервалу $[a, b)$ знаходить найближчу пару чисел з цього інтервалу, що кожне з них має парну кількість різних простих дільників. Якщо таких пар кілька, то цікавить пара з найменшим добутком.

По завданням 2 та 3 оцінюється правильність та якість коду, а також економність обчислень.

%enditem
}
{\smallskip\smallskip \begin {scriptsize} Затверджено на засіданні кафедри теоретичної кібернетики, протокол \No 3  від   08.11.2025 р.\par\smallskip\smallskip\smallskip\smallskip
{\parbox {6.5 cm} {Зав. кафедри \dotfill Ю.В.~Крак}}\hspace {0.7cm} {\hbox to 6.5 cm { Екзаменатор \dotfill Т.О.~Карнаух}} \end{scriptsize}}
%end
\newpage
%begin
{{\begin {scriptsize}\par
{ \center  Київський національний університет імені Тараса Шевченка \par}
{\parbox {5 cm}{Освітня програма \hfill}{\parbox {7 cm}{Інформатика\hfill}}\parbox {6 cm}{\hfill Семестр 1}}\par
{\parbox {5 cm} {Навчальний предмет \hfill}\parbox {7 cm}{Програмування \hfill}\parbox {6cm}{\hfill}\par}\vspace{-15pt} \end{scriptsize}}{\center Екзаменаційний білет \No \textbf { 47 }\par}}
{%beginitem
1. Які з наведених типів не є цілими (integral), якщо такі є серед перелічених?

bool, int, long, long long, unsigned char?
%enditem
\par
\vspace{2mm}
%beginitem
2. Скільки змінних, що є вказівниками, тут оголошено?
long pointer1, pointer2, pointer3, pointer4;?
%enditem
\par
\vspace{2mm}
%beginitem
3. За наявних оголошень \code{void f(double\&); double x;} які з виразів \code{f(x*x), f(x+=2), f(\&x)} є синтаксично коректними?
%enditem
\par
\vspace{2mm}
%beginitem
4. Що таке просторова складність алгоритму?
%enditem
\par
\vspace{2mm}
%beginitem
5. Яку часову оцінку складності має алгоритм сортування злиттям?
%enditem
\par
\vspace{2mm}
%beginitem
6. Написати програму, що запитує у користувача значення дійсних параметрів $x$ та $y$, обчислює для них значення виразу та повiдомляє введенi данi та результат обчислення.
$$\frac{\ln x - 48}{(\arctg x - 0.3) (y - 74)}$$ 

%enditem
\par
%beginitem
7. Написати функцію, що для послідовності цілих (задається масивом) знаходить довжину найдовшого відрізку, що складається з чисел, у сімковому запису якого нема цифри 2. За відсутності має повертати $0$. Використовувати додаткові структури даних (у тому числі рядки) заборонено.

По завданням 2 та 3 оцінюється правильність та якість коду, а також економність обчислень.

%enditem
}
{\smallskip\smallskip \begin {scriptsize} Затверджено на засіданні кафедри теоретичної кібернетики, протокол \No 3  від   08.11.2025 р.\par\smallskip\smallskip\smallskip\smallskip
{\parbox {6.5 cm} {Зав. кафедри \dotfill Ю.В.~Крак}}\hspace {0.7cm} {\hbox to 6.5 cm { Екзаменатор \dotfill Т.О.~Карнаух}} \end{scriptsize}}
%end
\vspace{0.9cm}
%begin
{{\begin {scriptsize}\par
{ \center  Київський національний університет імені Тараса Шевченка \par}
{\parbox {5 cm}{Освітня програма \hfill}{\parbox {7 cm}{Інформатика\hfill}}\parbox {6 cm}{\hfill Семестр 1}}\par
{\parbox {5 cm} {Навчальний предмет \hfill}\parbox {7 cm}{Програмування \hfill}\parbox {6cm}{\hfill}\par}\vspace{-15pt} \end{scriptsize}}{\center Екзаменаційний білет \No \textbf { 48 }\par}}
{%beginitem
1. Які з наведених типів не є цілими (integral), якщо такі є серед перелічених?

double, int, long*, unsigned char?
%enditem
\par
\vspace{2mm}
%beginitem
2. Скільки змінних, що є вказівниками, тут оголошено?
char *s1, pc, **s2;?
%enditem
\par
\vspace{2mm}
%beginitem
3. За наявних оголошень \code{void f(int\&); int n;} які з виразів \code{f(13), f(++n), f(\&n)} є синтаксично коректними?
%enditem
\par
\vspace{2mm}
%beginitem
4. Що таке С-рядок?
%enditem
\par
\vspace{2mm}
%beginitem
5. Яку часову оцінку складності має алгоритм швидкого сортування?
%enditem
\par
\vspace{2mm}
%beginitem
6. Написати програму, що запитує у користувача значення дійсних параметрів $x$ та $y$, обчислює для них значення виразу та повiдомляє введенi данi та результат обчислення.
$$\frac{\pi x + 27}{(\cos x + 0.4) (y + 84)}$$ 

%enditem
\par
%beginitem
7. Написати функцію, що повертає найбільше число з послідовності цілих (задається масивом), що має не більше трьох різних простих дільників. За відсутності має повертати None. Використовувати додаткові структури даних (у тому числі рядки) заборонено.

По завданням 2 та 3 оцінюється правильність та якість коду, а також економність обчислень.

%enditem
}
{\smallskip\smallskip \begin {scriptsize} Затверджено на засіданні кафедри теоретичної кібернетики, протокол \No 3  від   08.11.2025 р.\par\smallskip\smallskip\smallskip\smallskip
{\parbox {6.5 cm} {Зав. кафедри \dotfill Ю.В.~Крак}}\hspace {0.7cm} {\hbox to 6.5 cm { Екзаменатор \dotfill Т.О.~Карнаух}} \end{scriptsize}}
%end
\newpage
%begin
{{\begin {scriptsize}\par
{ \center  Київський національний університет імені Тараса Шевченка \par}
{\parbox {5 cm}{Освітня програма \hfill}{\parbox {7 cm}{Інформатика\hfill}}\parbox {6 cm}{\hfill Семестр 1}}\par
{\parbox {5 cm} {Навчальний предмет \hfill}\parbox {7 cm}{Програмування \hfill}\parbox {6cm}{\hfill}\par}\vspace{-15pt} \end{scriptsize}}{\center Екзаменаційний білет \No \textbf { 49 }\par}}
{%beginitem
1. Які з наведених типів не є цілими (integral), якщо такі є серед перелічених?

bool, char, int, unsigned?
%enditem
\par
\vspace{2mm}
%beginitem
2. Скільки змінних, що є вказівниками, тут оголошено?
int n, m, *p1, p2;?
%enditem
\par
\vspace{2mm}
%beginitem
3. За наявних оголошень \code{void f(char); char c;} які з виразів \code{f(c*c), f(c+=2), f(\&c)} є синтаксично коректними?
%enditem
\par
\vspace{2mm}
%beginitem
4. Що таке масив?
%enditem
\par
\vspace{2mm}
%beginitem
5. Яку часову оцінку складності має алгоритм сортування купою (він же деревом, пірамідою)?
%enditem
\par
\vspace{2mm}
%beginitem
6. Написати програму, що запитує у користувача значення дійсних параметрів $x$ та $y$, обчислює для них значення виразу та повiдомляє введенi данi та результат обчислення.
$$\frac{\log_{10} x + 24}{(\arccos x - 0.6) (y - 82)}$$ 

%enditem
\par
%beginitem
7. Написати функцію, яка для цілого числа знаходить кількість його різних простих дільників, що мають у десятковому запису цифру 1.

По завданням 2 та 3 оцінюється правильність та якість коду, а також економність обчислень.

%enditem
}
{\smallskip\smallskip \begin {scriptsize} Затверджено на засіданні кафедри теоретичної кібернетики, протокол \No 3  від   08.11.2025 р.\par\smallskip\smallskip\smallskip\smallskip
{\parbox {6.5 cm} {Зав. кафедри \dotfill Ю.В.~Крак}}\hspace {0.7cm} {\hbox to 6.5 cm { Екзаменатор \dotfill Т.О.~Карнаух}} \end{scriptsize}}
%end
\vspace{0.9cm}
%begin
{{\begin {scriptsize}\par
{ \center  Київський національний університет імені Тараса Шевченка \par}
{\parbox {5 cm}{Освітня програма \hfill}{\parbox {7 cm}{Інформатика\hfill}}\parbox {6 cm}{\hfill Семестр 1}}\par
{\parbox {5 cm} {Навчальний предмет \hfill}\parbox {7 cm}{Програмування \hfill}\parbox {6cm}{\hfill}\par}\vspace{-15pt} \end{scriptsize}}{\center Екзаменаційний білет \No \textbf { 50 }\par}}
{%beginitem
1. Які з наведених типів не є цілими (integral), якщо такі є серед перелічених?

int*, signed char, unsigned long, string?
%enditem
\par
\vspace{2mm}
%beginitem
2. Скільки змінних, що є вказівниками, тут оголошено?
long *n1, *n2, n3, pn4;?
%enditem
\par
\vspace{2mm}
%beginitem
3. За наявних оголошень \code{void f(int); int n;} які з виразів \code{f(3), f(++n), f(\&n)} є синтаксично коректними?
%enditem
\par
\vspace{2mm}
%beginitem
4. Що таке часова складність алгоритму?
%enditem
\par
\vspace{2mm}
%beginitem
5. Яку часову оцінку складності має алгоритм сортування злиттям?
%enditem
\par
\vspace{2mm}
%beginitem
6. Написати програму, що запитує у користувача значення дійсних параметрів $x$ та $y$, обчислює для них значення виразу та повiдомляє введенi данi та результат обчислення.
$$\frac{\ln x - 68}{(\cos x + 0.8) (y + 15)}$$ 

%enditem
\par
%beginitem
7. Написати функцію, що повертає число з послідовності цілих (задається масивом), що має парну кількість різних простих дільників. Якщо таких чисел кілька, то повернути найбільше. Використовувати додаткові структури даних (у тому числі рядки) заборонено.

По завданням 2 та 3 оцінюється правильність та якість коду, а також економність обчислень.

%enditem
}
{\smallskip\smallskip \begin {scriptsize} Затверджено на засіданні кафедри теоретичної кібернетики, протокол \No 3  від   08.11.2025 р.\par\smallskip\smallskip\smallskip\smallskip
{\parbox {6.5 cm} {Зав. кафедри \dotfill Ю.В.~Крак}}\hspace {0.7cm} {\hbox to 6.5 cm { Екзаменатор \dotfill Т.О.~Карнаух}} \end{scriptsize}}
%end
\newpage
%begin
{{\begin {scriptsize}\par
{ \center  Київський національний університет імені Тараса Шевченка \par}
{\parbox {5 cm}{Освітня програма \hfill}{\parbox {7 cm}{Інформатика\hfill}}\parbox {6 cm}{\hfill Семестр 1}}\par
{\parbox {5 cm} {Навчальний предмет \hfill}\parbox {7 cm}{Програмування \hfill}\parbox {6cm}{\hfill}\par}\vspace{-15pt} \end{scriptsize}}{\center Екзаменаційний білет \No \textbf { 51 }\par}}
{%beginitem
1. Які з наведених типів не є цілими (integral), якщо такі є серед перелічених?

double, int, long*, unsigned char?
%enditem
\par
\vspace{2mm}
%beginitem
2. Скільки змінних, що є вказівниками, тут оголошено?
int *pn1, pn2, pn3, pn4;?
%enditem
\par
\vspace{2mm}
%beginitem
3. За наявних оголошень \code{void f(double*); double x;} які з виразів \code{f(x*x), f(x+=2), f(\&x)} є синтаксично коректними?
%enditem
\par
\vspace{2mm}
%beginitem
4. Чим префіксний оператор ++ відрізняється від постфіксного?
%enditem
\par
\vspace{2mm}
%beginitem
5. Яку часову оцінку складності має алгоритм швидкого сортування?
%enditem
\par
\vspace{2mm}
%beginitem
6. Написати програму, що запитує у користувача значення дійсних параметрів $x$ та $y$, обчислює для них значення виразу та повiдомляє введенi данi та результат обчислення.
$$\frac{\log_{2} x - 70}{(\ctg x - 0.5) (y + 61)}$$ 

%enditem
\par
%beginitem
7. Написати функцію, що для інтервалу $[a, b)$ знаходить найближчу пару чисел з цього інтервалу, що кожне з них має непарну кількість різних простих дільників. Якщо таких пар кілька, то цікавить пара з найменшим добутком.

По завданням 2 та 3 оцінюється правильність та якість коду, а також економність обчислень.

%enditem
}
{\smallskip\smallskip \begin {scriptsize} Затверджено на засіданні кафедри теоретичної кібернетики, протокол \No 3  від   08.11.2025 р.\par\smallskip\smallskip\smallskip\smallskip
{\parbox {6.5 cm} {Зав. кафедри \dotfill Ю.В.~Крак}}\hspace {0.7cm} {\hbox to 6.5 cm { Екзаменатор \dotfill Т.О.~Карнаух}} \end{scriptsize}}
%end
\vspace{0.9cm}
%begin
{{\begin {scriptsize}\par
{ \center  Київський національний університет імені Тараса Шевченка \par}
{\parbox {5 cm}{Освітня програма \hfill}{\parbox {7 cm}{Інформатика\hfill}}\parbox {6 cm}{\hfill Семестр 1}}\par
{\parbox {5 cm} {Навчальний предмет \hfill}\parbox {7 cm}{Програмування \hfill}\parbox {6cm}{\hfill}\par}\vspace{-15pt} \end{scriptsize}}{\center Екзаменаційний білет \No \textbf { 52 }\par}}
{%beginitem
1. Які з наведених типів не є цілими (integral), якщо такі є серед перелічених?

int*, signed char, unsigned long, string?
%enditem
\par
\vspace{2mm}
%beginitem
2. Скільки змінних, що є вказівниками, тут оголошено?
long pointer1, pointer2, pointer3, pointer4;?
%enditem
\par
\vspace{2mm}
%beginitem
3. За наявних оголошень \code{void f(double); double x;} які з виразів \code{f(3.1), f(x+=2), f(\&x)} є синтаксично коректними?
%enditem
\par
\vspace{2mm}
%beginitem
4. Що таке часова складність задачі?
%enditem
\par
\vspace{2mm}
%beginitem
5. Яку часову оцінку складності має алгоритм сортування вставками?
%enditem
\par
\vspace{2mm}
%beginitem
6. Написати програму, що запитує у користувача значення дійсних параметрів $x$ та $y$, обчислює для них значення виразу та повiдомляє введенi данi та результат обчислення.
$$\frac{\ x + 36}{(\sin x + 0.7) (y - 44)}$$ 

%enditem
\par
%beginitem
7. Написати функцію, що повертає найменше число з послідовності цілих (задається масивом), що має не більше трьох різних простих дільників. За відсутності має повертати None. Використовувати додаткові структури даних (у тому числі рядки) заборонено.

По завданням 2 та 3 оцінюється правильність та якість коду, а також економність обчислень.

%enditem
}
{\smallskip\smallskip \begin {scriptsize} Затверджено на засіданні кафедри теоретичної кібернетики, протокол \No 3  від   08.11.2025 р.\par\smallskip\smallskip\smallskip\smallskip
{\parbox {6.5 cm} {Зав. кафедри \dotfill Ю.В.~Крак}}\hspace {0.7cm} {\hbox to 6.5 cm { Екзаменатор \dotfill Т.О.~Карнаух}} \end{scriptsize}}
%end
\newpage
%begin
{{\begin {scriptsize}\par
{ \center  Київський національний університет імені Тараса Шевченка \par}
{\parbox {5 cm}{Освітня програма \hfill}{\parbox {7 cm}{Інформатика\hfill}}\parbox {6 cm}{\hfill Семестр 1}}\par
{\parbox {5 cm} {Навчальний предмет \hfill}\parbox {7 cm}{Програмування \hfill}\parbox {6cm}{\hfill}\par}\vspace{-15pt} \end{scriptsize}}{\center Екзаменаційний білет \No \textbf { 53 }\par}}
{%beginitem
1. Які з наведених типів не є цілими (integral), якщо такі є серед перелічених?

bool, int, long, long long, unsigned char?
%enditem
\par
\vspace{2mm}
%beginitem
2. Скільки змінних, що є вказівниками, тут оголошено?
double x, *px1, px2, ***px3;?
%enditem
\par
\vspace{2mm}
%beginitem
3. За наявних оголошень \code{void f(int*); int n;} які з виразів \code{f(n++), f(++n), f(\&n)} є синтаксично коректними?
%enditem
\par
\vspace{2mm}
%beginitem
4. Що таке тип даних?
%enditem
\par
\vspace{2mm}
%beginitem
5. Яку часову оцінку складності має алгоритм швидкого сортування?
%enditem
\par
\vspace{2mm}
%beginitem
6. Написати програму, що запитує у користувача значення дійсних параметрів $x$ та $y$, обчислює для них значення виразу та повiдомляє введенi данi та результат обчислення.
$$\frac{\sqrt x + 38}{(\arccos x - 0.1) (y - 89)}$$ 

%enditem
\par
%beginitem
7. Написати функцію, що для цілого числа знаходить суму його простих дільників, що входять у його розклад не більш ніж у 4му степеню. Наприклад, для числа $3^5{\cdot}7^9{\cdot}11{\cdot}23^{2}$ таких дільників два – $11$ та $23$, тому функція має повернути $34$.

По завданням 2 та 3 оцінюється правильність та якість коду, а також економність обчислень.

%enditem
}
{\smallskip\smallskip \begin {scriptsize} Затверджено на засіданні кафедри теоретичної кібернетики, протокол \No 3  від   08.11.2025 р.\par\smallskip\smallskip\smallskip\smallskip
{\parbox {6.5 cm} {Зав. кафедри \dotfill Ю.В.~Крак}}\hspace {0.7cm} {\hbox to 6.5 cm { Екзаменатор \dotfill Т.О.~Карнаух}} \end{scriptsize}}
%end
\vspace{0.9cm}
%begin
{{\begin {scriptsize}\par
{ \center  Київський національний університет імені Тараса Шевченка \par}
{\parbox {5 cm}{Освітня програма \hfill}{\parbox {7 cm}{Інформатика\hfill}}\parbox {6 cm}{\hfill Семестр 1}}\par
{\parbox {5 cm} {Навчальний предмет \hfill}\parbox {7 cm}{Програмування \hfill}\parbox {6cm}{\hfill}\par}\vspace{-15pt} \end{scriptsize}}{\center Екзаменаційний білет \No \textbf { 54 }\par}}
{%beginitem
1. Які з наведених типів не є цілими (integral), якщо такі є серед перелічених?

bool, int, long, long long, unsigned char?
%enditem
\par
\vspace{2mm}
%beginitem
2. Скільки змінних, що є вказівниками, тут оголошено?
double *px, *py, **ppx;?
%enditem
\par
\vspace{2mm}
%beginitem
3. За наявних оголошень \code{void f(char\&); char c;} які з виразів \code{f(c*c), f(c++), f(\&c)} є синтаксично коректними?
%enditem
\par
\vspace{2mm}
%beginitem
4. У чому відмінність літералів від констант?
%enditem
\par
\vspace{2mm}
%beginitem
5. Яку часову оцінку складності має алгоритм сортування купою (він же деревом, пірамідою)?
%enditem
\par
\vspace{2mm}
%beginitem
6. Написати програму, що запитує у користувача значення дійсних параметрів $x$ та $y$, обчислює для них значення виразу та повiдомляє введенi данi та результат обчислення.
$$\frac{\pi x - 96}{(\tg x + 0.9) (y + 78)}$$ 

%enditem
\par
%beginitem
7. Написати функцію, що знаходить найбільше число з послідовності цілих (задається масивом), у десятковому запису якого нема цифри 4. Повертає індекс першого входження цього числа або $-1$ (за відсутності). Використовувати додаткові структури даних (у тому числі рядки) заборонено.

По завданням 2 та 3 оцінюється правильність та якість коду, а також економність обчислень.

%enditem
}
{\smallskip\smallskip \begin {scriptsize} Затверджено на засіданні кафедри теоретичної кібернетики, протокол \No 3  від   08.11.2025 р.\par\smallskip\smallskip\smallskip\smallskip
{\parbox {6.5 cm} {Зав. кафедри \dotfill Ю.В.~Крак}}\hspace {0.7cm} {\hbox to 6.5 cm { Екзаменатор \dotfill Т.О.~Карнаух}} \end{scriptsize}}
%end
\newpage
%begin
{{\begin {scriptsize}\par
{ \center  Київський національний університет імені Тараса Шевченка \par}
{\parbox {5 cm}{Освітня програма \hfill}{\parbox {7 cm}{Інформатика\hfill}}\parbox {6 cm}{\hfill Семестр 1}}\par
{\parbox {5 cm} {Навчальний предмет \hfill}\parbox {7 cm}{Програмування \hfill}\parbox {6cm}{\hfill}\par}\vspace{-15pt} \end{scriptsize}}{\center Екзаменаційний білет \No \textbf { 55 }\par}}
{%beginitem
1. Які з наведених типів не є цілими (integral), якщо такі є серед перелічених?

char, double, int, float, unsigned*?
%enditem
\par
\vspace{2mm}
%beginitem
2. Скільки змінних, що є вказівниками, тут оголошено?
char *s1, pc, **s2;?
%enditem
\par
\vspace{2mm}
%beginitem
3. За наявних оголошень \code{void f(double*); double x;} які з виразів \code{f(x*x), f(x+=2), f(\&x)} є синтаксично коректними?
%enditem
\par
\vspace{2mm}
%beginitem
4. Що таке просторова складність задачі?
%enditem
\par
\vspace{2mm}
%beginitem
5. Яку часову оцінку складності має алгоритм сортування вибором?
%enditem
\par
\vspace{2mm}
%beginitem
6. Написати програму, що запитує у користувача значення дійсних параметрів $x$ та $y$, обчислює для них значення виразу та повiдомляє введенi данi та результат обчислення.
$$\frac{\ x + 86}{(\arctg x + 0.4) (y + 42)}$$ 

%enditem
\par
%beginitem
7. Написати функцію, що для інтервалу $[a, b)$ знаходить найближчу пару чисел з цього інтервалу, що мають не менше трьох простих дільників та мають цифру 3 у своєму десятковому запису.

По завданням 2 та 3 оцінюється правильність та якість коду, а також економність обчислень.

%enditem
}
{\smallskip\smallskip \begin {scriptsize} Затверджено на засіданні кафедри теоретичної кібернетики, протокол \No 3  від   08.11.2025 р.\par\smallskip\smallskip\smallskip\smallskip
{\parbox {6.5 cm} {Зав. кафедри \dotfill Ю.В.~Крак}}\hspace {0.7cm} {\hbox to 6.5 cm { Екзаменатор \dotfill Т.О.~Карнаух}} \end{scriptsize}}
%end
\vspace{0.9cm}
%begin
{{\begin {scriptsize}\par
{ \center  Київський національний університет імені Тараса Шевченка \par}
{\parbox {5 cm}{Освітня програма \hfill}{\parbox {7 cm}{Інформатика\hfill}}\parbox {6 cm}{\hfill Семестр 1}}\par
{\parbox {5 cm} {Навчальний предмет \hfill}\parbox {7 cm}{Програмування \hfill}\parbox {6cm}{\hfill}\par}\vspace{-15pt} \end{scriptsize}}{\center Екзаменаційний білет \No \textbf { 56 }\par}}
{%beginitem
1. Які з наведених типів не є цілими (integral), якщо такі є серед перелічених?

bool, char, int, unsigned?
%enditem
\par
\vspace{2mm}
%beginitem
2. Скільки змінних, що є вказівниками, тут оголошено?
int n, m, *p1, p2;?
%enditem
\par
\vspace{2mm}
%beginitem
3. За наявних оголошень \code{void f(double); double x;} які з виразів \code{f(3.1), f(x+=2), f(\&x)} є синтаксично коректними?
%enditem
\par
\vspace{2mm}
%beginitem
4. Що таке С-рядок?
%enditem
\par
\vspace{2mm}
%beginitem
5. Яку часову оцінку складності має алгоритм сортування бульбашкою?
%enditem
\par
\vspace{2mm}
%beginitem
6. Написати програму, що запитує у користувача значення дійсних параметрів $x$ та $y$, обчислює для них значення виразу та повiдомляє введенi данi та результат обчислення.
$$\frac{\sqrt x - 41}{(\arcsin x - 0.3) (y - 52)}$$ 

%enditem
\par
%beginitem
7. Написати функцію, що для цілого числа знаходить простий дільник, який входить у його розклад ц найбільшому степені. Якщо таких кілька, то повернути найбільший. Наприклад, для числа $2^{9}{\cdot}3^5{\cdot}7^9{\cdot}11$ функція має повернути $7$.

По завданням 2 та 3 оцінюється правильність та якість коду, а також економність обчислень.

%enditem
}
{\smallskip\smallskip \begin {scriptsize} Затверджено на засіданні кафедри теоретичної кібернетики, протокол \No 3  від   08.11.2025 р.\par\smallskip\smallskip\smallskip\smallskip
{\parbox {6.5 cm} {Зав. кафедри \dotfill Ю.В.~Крак}}\hspace {0.7cm} {\hbox to 6.5 cm { Екзаменатор \dotfill Т.О.~Карнаух}} \end{scriptsize}}
%end
\newpage
%begin
{{\begin {scriptsize}\par
{ \center  Київський національний університет імені Тараса Шевченка \par}
{\parbox {5 cm}{Освітня програма \hfill}{\parbox {7 cm}{Інформатика\hfill}}\parbox {6 cm}{\hfill Семестр 1}}\par
{\parbox {5 cm} {Навчальний предмет \hfill}\parbox {7 cm}{Програмування \hfill}\parbox {6cm}{\hfill}\par}\vspace{-15pt} \end{scriptsize}}{\center Екзаменаційний білет \No \textbf { 57 }\par}}
{%beginitem
1. Які з наведених типів не є цілими (integral), якщо такі є серед перелічених?

bool*, char, int, unsigned?
%enditem
\par
\vspace{2mm}
%beginitem
2. Скільки змінних, що є вказівниками, тут оголошено?
long pointer1, pointer2, pointer3, pointer4;?
%enditem
\par
\vspace{2mm}
%beginitem
3. За наявних оголошень \code{void f(int*); int n;} які з виразів \code{f(n++), f(++n), f(\&n)} є синтаксично коректними?
%enditem
\par
\vspace{2mm}
%beginitem
4. Що таке тип даних?
%enditem
\par
\vspace{2mm}
%beginitem
5. Яку часову оцінку складності має алгоритм сортування вибором?
%enditem
\par
\vspace{2mm}
%beginitem
6. Написати програму, що запитує у користувача значення дійсних параметрів $x$ та $y$, обчислює для них значення виразу та повiдомляє введенi данi та результат обчислення.
$$\frac{\pi x - 87}{(\arccos x - 0.5) (y + 93)}$$ 

%enditem
\par
%beginitem
7. Написати функцію, що для послідовності цілих (задається масивом) знаходить найдовший відрізок, в якому всі сусідні числа відрізняються якнайменш на 3. Якщо таких відрізків кілька, то повернути перший (як індекс початку та довжину). Використовувати додаткові структури даних (у тому числі рядки) заборонено.

По завданням 2 та 3 оцінюється правильність та якість коду, а також економність обчислень.

%enditem
}
{\smallskip\smallskip \begin {scriptsize} Затверджено на засіданні кафедри теоретичної кібернетики, протокол \No 3  від   08.11.2025 р.\par\smallskip\smallskip\smallskip\smallskip
{\parbox {6.5 cm} {Зав. кафедри \dotfill Ю.В.~Крак}}\hspace {0.7cm} {\hbox to 6.5 cm { Екзаменатор \dotfill Т.О.~Карнаух}} \end{scriptsize}}
%end
\vspace{0.9cm}
%begin
{{\begin {scriptsize}\par
{ \center  Київський національний університет імені Тараса Шевченка \par}
{\parbox {5 cm}{Освітня програма \hfill}{\parbox {7 cm}{Інформатика\hfill}}\parbox {6 cm}{\hfill Семестр 1}}\par
{\parbox {5 cm} {Навчальний предмет \hfill}\parbox {7 cm}{Програмування \hfill}\parbox {6cm}{\hfill}\par}\vspace{-15pt} \end{scriptsize}}{\center Екзаменаційний білет \No \textbf { 58 }\par}}
{%beginitem
1. Які з наведених типів не є цілими (integral), якщо такі є серед перелічених?

char, double, int, float, unsigned?
%enditem
\par
\vspace{2mm}
%beginitem
2. Скільки змінних, що є вказівниками, тут оголошено?
int n, m, *p1, p2;?
%enditem
\par
\vspace{2mm}
%beginitem
3. За наявних оголошень \code{void f(char\&); char c;} які з виразів \code{f(c*c), f(c++), f(\&c)} є синтаксично коректними?
%enditem
\par
\vspace{2mm}
%beginitem
4. Що таке часова складність задачі?
%enditem
\par
\vspace{2mm}
%beginitem
5. Яку часову оцінку складності має алгоритм сортування купою (він же деревом, пірамідою)?
%enditem
\par
\vspace{2mm}
%beginitem
6. Написати програму, що запитує у користувача значення дійсних параметрів $x$ та $y$, обчислює для них значення виразу та повiдомляє введенi данi та результат обчислення.
$$\frac{\ln x + 76}{(\ctg x + 0.1) (y - 49)}$$ 

%enditem
\par
%beginitem
7. Написати функцію, що для інтервалу $[a, b)$ знаходить найближчу пару чисел з цього інтервалу, що кожне з них має непарну кількість різних простих дільників. Якщо таких пар кілька, то цікавить пара з найбільшою сумою.

По завданням 2 та 3 оцінюється правильність та якість коду, а також економність обчислень.

%enditem
}
{\smallskip\smallskip \begin {scriptsize} Затверджено на засіданні кафедри теоретичної кібернетики, протокол \No 3  від   08.11.2025 р.\par\smallskip\smallskip\smallskip\smallskip
{\parbox {6.5 cm} {Зав. кафедри \dotfill Ю.В.~Крак}}\hspace {0.7cm} {\hbox to 6.5 cm { Екзаменатор \dotfill Т.О.~Карнаух}} \end{scriptsize}}
%end
\newpage
%begin
{{\begin {scriptsize}\par
{ \center  Київський національний університет імені Тараса Шевченка \par}
{\parbox {5 cm}{Освітня програма \hfill}{\parbox {7 cm}{Інформатика\hfill}}\parbox {6 cm}{\hfill Семестр 1}}\par
{\parbox {5 cm} {Навчальний предмет \hfill}\parbox {7 cm}{Програмування \hfill}\parbox {6cm}{\hfill}\par}\vspace{-15pt} \end{scriptsize}}{\center Екзаменаційний білет \No \textbf { 59 }\par}}
{%beginitem
1. Які з наведених типів не є цілими (integral), якщо такі є серед перелічених?

double, int, long, unsigned char?
%enditem
\par
\vspace{2mm}
%beginitem
2. Скільки змінних, що є вказівниками, тут оголошено?
long *n1, *n2, n3, pn4;?
%enditem
\par
\vspace{2mm}
%beginitem
3. За наявних оголошень \code{void f(int\&); int n;} які з виразів \code{f(13), f(++n), f(\&n)} є синтаксично коректними?
%enditem
\par
\vspace{2mm}
%beginitem
4. Що таке просторова складність задачі?
%enditem
\par
\vspace{2mm}
%beginitem
5. Яку часову оцінку складності має алгоритм швидкого сортування?
%enditem
\par
\vspace{2mm}
%beginitem
6. Написати програму, що запитує у користувача значення дійсних параметрів $x$ та $y$, обчислює для них значення виразу та повiдомляє введенi данi та результат обчислення.
$$\frac{\log_{10} x + 60}{(\arcsin x - 0.2) (y - 50)}$$ 

%enditem
\par
%beginitem
7. Написати функцію, що для послідовності цілих (задається масивом) знаходить найдовший відрізок, в якому всі сусідні числа різні. Якщо таких відрізків кілька, то повернути перший (як індекс початку та довжину). Використовувати додаткові структури даних (у тому числі рядки) заборонено.

По завданням 2 та 3 оцінюється правильність та якість коду, а також економність обчислень.

%enditem
}
{\smallskip\smallskip \begin {scriptsize} Затверджено на засіданні кафедри теоретичної кібернетики, протокол \No 3  від   08.11.2025 р.\par\smallskip\smallskip\smallskip\smallskip
{\parbox {6.5 cm} {Зав. кафедри \dotfill Ю.В.~Крак}}\hspace {0.7cm} {\hbox to 6.5 cm { Екзаменатор \dotfill Т.О.~Карнаух}} \end{scriptsize}}
%end
\vspace{0.9cm}
%begin
{{\begin {scriptsize}\par
{ \center  Київський національний університет імені Тараса Шевченка \par}
{\parbox {5 cm}{Освітня програма \hfill}{\parbox {7 cm}{Інформатика\hfill}}\parbox {6 cm}{\hfill Семестр 1}}\par
{\parbox {5 cm} {Навчальний предмет \hfill}\parbox {7 cm}{Програмування \hfill}\parbox {6cm}{\hfill}\par}\vspace{-15pt} \end{scriptsize}}{\center Екзаменаційний білет \No \textbf { 60 }\par}}
{%beginitem
1. Які з наведених типів не є цілими (integral), якщо такі є серед перелічених?

int, signed char, unsigned long, string?
%enditem
\par
\vspace{2mm}
%beginitem
2. Скільки змінних, що є вказівниками, тут оголошено?
char *s1, pc, **s2;?
%enditem
\par
\vspace{2mm}
%beginitem
3. За наявних оголошень \code{void f(char); char c;} які з виразів \code{f(c*c), f(c+=2), f(\&c)} є синтаксично коректними?
%enditem
\par
\vspace{2mm}
%beginitem
4. Що означає термін ``часова складність у середньому''?
%enditem
\par
\vspace{2mm}
%beginitem
5. Яку часову оцінку складності має алгоритм сортування вставками?
%enditem
\par
\vspace{2mm}
%beginitem
6. Написати програму, що запитує у користувача значення дійсних параметрів $x$ та $y$, обчислює для них значення виразу та повiдомляє введенi данi та результат обчислення.
$$\frac{\log_{2} x - 47}{(\tg x + 0.8) (y + 99)}$$ 

%enditem
\par
%beginitem
7. Написати функцію, що для послідовності цілих (задається масивом) знаходить довжину найдовшого відрізку, в якому всі сусідні числа відрізняються якнайменш на 2. Використовувати додаткові структури даних (у тому числі рядки) заборонено.

По завданням 2 та 3 оцінюється правильність та якість коду, а також економність обчислень.

%enditem
}
{\smallskip\smallskip \begin {scriptsize} Затверджено на засіданні кафедри теоретичної кібернетики, протокол \No 3  від   08.11.2025 р.\par\smallskip\smallskip\smallskip\smallskip
{\parbox {6.5 cm} {Зав. кафедри \dotfill Ю.В.~Крак}}\hspace {0.7cm} {\hbox to 6.5 cm { Екзаменатор \dotfill Т.О.~Карнаух}} \end{scriptsize}}
%end
\newpage
%begin
{{\begin {scriptsize}\par
{ \center  Київський національний університет імені Тараса Шевченка \par}
{\parbox {5 cm}{Освітня програма \hfill}{\parbox {7 cm}{Інформатика\hfill}}\parbox {6 cm}{\hfill Семестр 1}}\par
{\parbox {5 cm} {Навчальний предмет \hfill}\parbox {7 cm}{Програмування \hfill}\parbox {6cm}{\hfill}\par}\vspace{-15pt} \end{scriptsize}}{\center Екзаменаційний білет \No \textbf { 61 }\par}}
{%beginitem
1. Які з наведених типів не є цілими (integral), якщо такі є серед перелічених?

bool, int, long, long long, unsigned char?
%enditem
\par
\vspace{2mm}
%beginitem
2. Скільки змінних, що є вказівниками, тут оголошено?
double *px, *py, **ppx;?
%enditem
\par
\vspace{2mm}
%beginitem
3. За наявних оголошень \code{void f(char*); char c;} які з виразів \code{f(c*c), f(c+=2), f(\&c)} є синтаксично коректними?
%enditem
\par
\vspace{2mm}
%beginitem
4. Що таке часова складність алгоритму?
%enditem
\par
\vspace{2mm}
%beginitem
5. Яку часову оцінку складності має алгоритм сортування злиттям?
%enditem
\par
\vspace{2mm}
%beginitem
6. Написати програму, що запитує у користувача значення дійсних параметрів $x$ та $y$, обчислює для них значення виразу та повiдомляє введенi данi та результат обчислення.
$$\frac{\ln x - 26}{(\cos x + 0.6) (y + 19)}$$ 

%enditem
\par
%beginitem
7. Написати функцію, що в інтервалі $[a, b)$ знаходить ціле число, що ділиться на найбільший степінь числа 3. Якщо таких чисел декілька, повернути найбільше.

По завданням 2 та 3 оцінюється правильність та якість коду, а також економність обчислень.

%enditem
}
{\smallskip\smallskip \begin {scriptsize} Затверджено на засіданні кафедри теоретичної кібернетики, протокол \No 3  від   08.11.2025 р.\par\smallskip\smallskip\smallskip\smallskip
{\parbox {6.5 cm} {Зав. кафедри \dotfill Ю.В.~Крак}}\hspace {0.7cm} {\hbox to 6.5 cm { Екзаменатор \dotfill Т.О.~Карнаух}} \end{scriptsize}}
%end
\vspace{0.9cm}
%begin
{{\begin {scriptsize}\par
{ \center  Київський національний університет імені Тараса Шевченка \par}
{\parbox {5 cm}{Освітня програма \hfill}{\parbox {7 cm}{Інформатика\hfill}}\parbox {6 cm}{\hfill Семестр 1}}\par
{\parbox {5 cm} {Навчальний предмет \hfill}\parbox {7 cm}{Програмування \hfill}\parbox {6cm}{\hfill}\par}\vspace{-15pt} \end{scriptsize}}{\center Екзаменаційний білет \No \textbf { 62 }\par}}
{%beginitem
1. Які з наведених типів не є цілими (integral), якщо такі є серед перелічених?

bool*, char, int, unsigned?
%enditem
\par
\vspace{2mm}
%beginitem
2. Скільки змінних, що є вказівниками, тут оголошено?
double x, *px1, px2, ***px3;?
%enditem
\par
\vspace{2mm}
%beginitem
3. За наявних оголошень \code{void f(double\&); double x;} які з виразів \code{f(x*x), f(x+=2), f(\&x)} є синтаксично коректними?
%enditem
\par
\vspace{2mm}
%beginitem
4. У чому відмінність літералів від констант?
%enditem
\par
\vspace{2mm}
%beginitem
5. Яку часову оцінку складності має алгоритм сортування бульбашкою?
%enditem
\par
\vspace{2mm}
%beginitem
6. Написати програму, що запитує у користувача значення дійсних параметрів $x$ та $y$, обчислює для них значення виразу та повiдомляє введенi данi та результат обчислення.
$$\frac{\ x + 32}{(\arctg x - 0.9) (y - 95)}$$ 

%enditem
\par
%beginitem
7. Написати функцію, що повертає найбільше число з послідовності цілих (задається масивом), що має рівно три різних простих дільники. За відсутності має повертати None. Використовувати додаткові структури даних (у тому числі рядки) заборонено.

По завданням 2 та 3 оцінюється правильність та якість коду, а також економність обчислень.

%enditem
}
{\smallskip\smallskip \begin {scriptsize} Затверджено на засіданні кафедри теоретичної кібернетики, протокол \No 3  від   08.11.2025 р.\par\smallskip\smallskip\smallskip\smallskip
{\parbox {6.5 cm} {Зав. кафедри \dotfill Ю.В.~Крак}}\hspace {0.7cm} {\hbox to 6.5 cm { Екзаменатор \dotfill Т.О.~Карнаух}} \end{scriptsize}}
%end
\newpage
%begin
{{\begin {scriptsize}\par
{ \center  Київський національний університет імені Тараса Шевченка \par}
{\parbox {5 cm}{Освітня програма \hfill}{\parbox {7 cm}{Інформатика\hfill}}\parbox {6 cm}{\hfill Семестр 1}}\par
{\parbox {5 cm} {Навчальний предмет \hfill}\parbox {7 cm}{Програмування \hfill}\parbox {6cm}{\hfill}\par}\vspace{-15pt} \end{scriptsize}}{\center Екзаменаційний білет \No \textbf { 63 }\par}}
{%beginitem
1. Які з наведених типів не є цілими (integral), якщо такі є серед перелічених?

int, signed char, unsigned long, string?
%enditem
\par
\vspace{2mm}
%beginitem
2. Скільки змінних, що є вказівниками, тут оголошено?
int *pn1, pn2, pn3, pn4;?
%enditem
\par
\vspace{2mm}
%beginitem
3. За наявних оголошень \code{void f(int); int n;} які з виразів \code{f(3), f(++n), f(\&n)} є синтаксично коректними?
%enditem
\par
\vspace{2mm}
%beginitem
4. Чим префіксний оператор ++ відрізняється від постфіксного?
%enditem
\par
\vspace{2mm}
%beginitem
5. Яку часову оцінку складності має алгоритм швидкого сортування?
%enditem
\par
\vspace{2mm}
%beginitem
6. Написати програму, що запитує у користувача значення дійсних параметрів $x$ та $y$, обчислює для них значення виразу та повiдомляє введенi данi та результат обчислення.
$$\frac{\log_{10} x + 12}{(\sin x + 0.7) (y - 88)}$$ 

%enditem
\par
%beginitem
7. Написати функцію, що повертає число з послідовності цілих (задається масивом), що має найменшу кількість різних простих дільників. Якщо таких чисел кілька, то повернути найбільше. Використовувати додаткові структури даних (у тому числі рядки) заборонено.

По завданням 2 та 3 оцінюється правильність та якість коду, а також економність обчислень.

%enditem
}
{\smallskip\smallskip \begin {scriptsize} Затверджено на засіданні кафедри теоретичної кібернетики, протокол \No 3  від   08.11.2025 р.\par\smallskip\smallskip\smallskip\smallskip
{\parbox {6.5 cm} {Зав. кафедри \dotfill Ю.В.~Крак}}\hspace {0.7cm} {\hbox to 6.5 cm { Екзаменатор \dotfill Т.О.~Карнаух}} \end{scriptsize}}
%end
\vspace{0.9cm}
%begin
{{\begin {scriptsize}\par
{ \center  Київський національний університет імені Тараса Шевченка \par}
{\parbox {5 cm}{Освітня програма \hfill}{\parbox {7 cm}{Інформатика\hfill}}\parbox {6 cm}{\hfill Семестр 1}}\par
{\parbox {5 cm} {Навчальний предмет \hfill}\parbox {7 cm}{Програмування \hfill}\parbox {6cm}{\hfill}\par}\vspace{-15pt} \end{scriptsize}}{\center Екзаменаційний білет \No \textbf { 64 }\par}}
{%beginitem
1. Які з наведених типів не є цілими (integral), якщо такі є серед перелічених?

bool, int, long, long long, unsigned char?
%enditem
\par
\vspace{2mm}
%beginitem
2. Скільки змінних, що є вказівниками, тут оголошено?
double *px, *py, **ppx;?
%enditem
\par
\vspace{2mm}
%beginitem
3. За наявних оголошень \code{void f(char\&); char c;} які з виразів \code{f(c*c), f(c++), f(\&c)} є синтаксично коректними?
%enditem
\par
\vspace{2mm}
%beginitem
4. Які параметри може приймати функція main? Що вони позначають та яких вони типів?
%enditem
\par
\vspace{2mm}
%beginitem
5. Яку часову оцінку складності має алгоритм сортування вставками?
%enditem
\par
\vspace{2mm}
%beginitem
6. Написати програму, що запитує у користувача значення дійсних параметрів $x$ та $y$, обчислює для них значення виразу та повiдомляє введенi данi та результат обчислення.
$$\frac{\log_{2} x - 64}{(\arcsin x - 0.2) (y + 62)}$$ 

%enditem
\par
%beginitem
7. Написати функцію, що знаходить найменше число з послідовності цілих (задається масивом), у десятковому запису якого нема цифри 7. Повертає індекс першого входження цього числа або $-1$ (за відсутності). Використовувати додаткові структури даних (у тому числі рядки) заборонено.

По завданням 2 та 3 оцінюється правильність та якість коду, а також економність обчислень.

%enditem
}
{\smallskip\smallskip \begin {scriptsize} Затверджено на засіданні кафедри теоретичної кібернетики, протокол \No 3  від   08.11.2025 р.\par\smallskip\smallskip\smallskip\smallskip
{\parbox {6.5 cm} {Зав. кафедри \dotfill Ю.В.~Крак}}\hspace {0.7cm} {\hbox to 6.5 cm { Екзаменатор \dotfill Т.О.~Карнаух}} \end{scriptsize}}
%end
\newpage
%begin
{{\begin {scriptsize}\par
{ \center  Київський національний університет імені Тараса Шевченка \par}
{\parbox {5 cm}{Освітня програма \hfill}{\parbox {7 cm}{Інформатика\hfill}}\parbox {6 cm}{\hfill Семестр 1}}\par
{\parbox {5 cm} {Навчальний предмет \hfill}\parbox {7 cm}{Програмування \hfill}\parbox {6cm}{\hfill}\par}\vspace{-15pt} \end{scriptsize}}{\center Екзаменаційний білет \No \textbf { 65 }\par}}
{%beginitem
1. Які з наведених типів не є цілими (integral), якщо такі є серед перелічених?

double, int, long*, unsigned char?
%enditem
\par
\vspace{2mm}
%beginitem
2. Скільки змінних, що є вказівниками, тут оголошено?
int *pn1, pn2, pn3, pn4;?
%enditem
\par
\vspace{2mm}
%beginitem
3. За наявних оголошень \code{void f(double); double x;} які з виразів \code{f(3.1), f(x+=2), f(\&x)} є синтаксично коректними?
%enditem
\par
\vspace{2mm}
%beginitem
4. Що таке просторова складність алгоритму?
%enditem
\par
\vspace{2mm}
%beginitem
5. Яку часову оцінку складності має алгоритм сортування вибором?
%enditem
\par
\vspace{2mm}
%beginitem
6. Написати програму, що запитує у користувача значення дійсних параметрів $x$ та $y$, обчислює для них значення виразу та повiдомляє введенi данi та результат обчислення.
$$\frac{\pi x - 53}{(\tg x - 0.3) (y - 14)}$$ 

%enditem
\par
%beginitem
7. Написати функцію, що для цілого числа знаходить простий дільник, який входить у його розклад у найбільшому степені. Якщо таких кілька, то повернути найменший. Наприклад, для числа $2^{9}{\cdot}3^5{\cdot}7^9{\cdot}11$ функція має повернути $2$.

По завданням 2 та 3 оцінюється правильність та якість коду, а також економність обчислень.

%enditem
}
{\smallskip\smallskip \begin {scriptsize} Затверджено на засіданні кафедри теоретичної кібернетики, протокол \No 3  від   08.11.2025 р.\par\smallskip\smallskip\smallskip\smallskip
{\parbox {6.5 cm} {Зав. кафедри \dotfill Ю.В.~Крак}}\hspace {0.7cm} {\hbox to 6.5 cm { Екзаменатор \dotfill Т.О.~Карнаух}} \end{scriptsize}}
%end
\vspace{0.9cm}
%begin
{{\begin {scriptsize}\par
{ \center  Київський національний університет імені Тараса Шевченка \par}
{\parbox {5 cm}{Освітня програма \hfill}{\parbox {7 cm}{Інформатика\hfill}}\parbox {6 cm}{\hfill Семестр 1}}\par
{\parbox {5 cm} {Навчальний предмет \hfill}\parbox {7 cm}{Програмування \hfill}\parbox {6cm}{\hfill}\par}\vspace{-15pt} \end{scriptsize}}{\center Екзаменаційний білет \No \textbf { 66 }\par}}
{%beginitem
1. Які з наведених типів не є цілими (integral), якщо такі є серед перелічених?

double, int, long, unsigned char?
%enditem
\par
\vspace{2mm}
%beginitem
2. Скільки змінних, що є вказівниками, тут оголошено?
long pointer1, pointer2, pointer3, pointer4;?
%enditem
\par
\vspace{2mm}
%beginitem
3. За наявних оголошень \code{void f(double\&); double x;} які з виразів \code{f(x*x), f(x+=2), f(\&x)} є синтаксично коректними?
%enditem
\par
\vspace{2mm}
%beginitem
4. Що таке масив?
%enditem
\par
\vspace{2mm}
%beginitem
5. Яку часову оцінку складності має алгоритм сортування злиттям?
%enditem
\par
\vspace{2mm}
%beginitem
6. Написати програму, що запитує у користувача значення дійсних параметрів $x$ та $y$, обчислює для них значення виразу та повiдомляє введенi данi та результат обчислення.
$$\frac{\sqrt x + 43}{(\cos x + 0.8) (y + 46)}$$ 

%enditem
\par
%beginitem
7. Написати функцію, що для інтервалу $[a, b)$ знаходить найближчу пару чисел з цього інтервалу, що кожне з них має парну кількість різних простих дільників. Якщо таких пар кілька, то цікавить пара з найбільшою сумою.

По завданням 2 та 3 оцінюється правильність та якість коду, а також економність обчислень.

%enditem
}
{\smallskip\smallskip \begin {scriptsize} Затверджено на засіданні кафедри теоретичної кібернетики, протокол \No 3  від   08.11.2025 р.\par\smallskip\smallskip\smallskip\smallskip
{\parbox {6.5 cm} {Зав. кафедри \dotfill Ю.В.~Крак}}\hspace {0.7cm} {\hbox to 6.5 cm { Екзаменатор \dotfill Т.О.~Карнаух}} \end{scriptsize}}
%end
\newpage
%begin
{{\begin {scriptsize}\par
{ \center  Київський національний університет імені Тараса Шевченка \par}
{\parbox {5 cm}{Освітня програма \hfill}{\parbox {7 cm}{Інформатика\hfill}}\parbox {6 cm}{\hfill Семестр 1}}\par
{\parbox {5 cm} {Навчальний предмет \hfill}\parbox {7 cm}{Програмування \hfill}\parbox {6cm}{\hfill}\par}\vspace{-15pt} \end{scriptsize}}{\center Екзаменаційний білет \No \textbf { 67 }\par}}
{%beginitem
1. Які з наведених типів не є цілими (integral), якщо такі є серед перелічених?

bool, char, int, unsigned?
%enditem
\par
\vspace{2mm}
%beginitem
2. Скільки змінних, що є вказівниками, тут оголошено?
double x, *px1, px2, ***px3;?
%enditem
\par
\vspace{2mm}
%beginitem
3. За наявних оголошень \code{void f(char); char c;} які з виразів \code{f(c*c), f(c+=2), f(\&c)} є синтаксично коректними?
%enditem
\par
\vspace{2mm}
%beginitem
4. Що таке масив?
%enditem
\par
\vspace{2mm}
%beginitem
5. Яку часову оцінку складності має алгоритм сортування купою (він же деревом, пірамідою)?
%enditem
\par
\vspace{2mm}
%beginitem
6. Написати програму, що запитує у користувача значення дійсних параметрів $x$ та $y$, обчислює для них значення виразу та повiдомляє введенi данi та результат обчислення.
$$\frac{\sqrt x - 66}{(\arctg x + 0.1) (y - 51)}$$ 

%enditem
\par
%beginitem
7. Написати функцію, що знаходить найменше число з послідовності цілих (задається масивом), у десятковому запису якого нема цифри 5. Повертає індекс останнього входження цього числа або $-1$ (за відсутності). Використовувати додаткові структури даних (у тому числі рядки) заборонено.

По завданням 2 та 3 оцінюється правильність та якість коду, а також економність обчислень.

%enditem
}
{\smallskip\smallskip \begin {scriptsize} Затверджено на засіданні кафедри теоретичної кібернетики, протокол \No 3  від   08.11.2025 р.\par\smallskip\smallskip\smallskip\smallskip
{\parbox {6.5 cm} {Зав. кафедри \dotfill Ю.В.~Крак}}\hspace {0.7cm} {\hbox to 6.5 cm { Екзаменатор \dotfill Т.О.~Карнаух}} \end{scriptsize}}
%end
\vspace{0.9cm}
%begin
{{\begin {scriptsize}\par
{ \center  Київський національний університет імені Тараса Шевченка \par}
{\parbox {5 cm}{Освітня програма \hfill}{\parbox {7 cm}{Інформатика\hfill}}\parbox {6 cm}{\hfill Семестр 1}}\par
{\parbox {5 cm} {Навчальний предмет \hfill}\parbox {7 cm}{Програмування \hfill}\parbox {6cm}{\hfill}\par}\vspace{-15pt} \end{scriptsize}}{\center Екзаменаційний білет \No \textbf { 68 }\par}}
{%beginitem
1. Які з наведених типів не є цілими (integral), якщо такі є серед перелічених?

int*, signed char, unsigned long, string?
%enditem
\par
\vspace{2mm}
%beginitem
2. Скільки змінних, що є вказівниками, тут оголошено?
long *n1, *n2, n3, pn4;?
%enditem
\par
\vspace{2mm}
%beginitem
3. За наявних оголошень \code{void f(int\&); int n;} які з виразів \code{f(13), f(++n), f(\&n)} є синтаксично коректними?
%enditem
\par
\vspace{2mm}
%beginitem
4. Що таке С-рядок?
%enditem
\par
\vspace{2mm}
%beginitem
5. Яку часову оцінку складності має алгоритм швидкого сортування?
%enditem
\par
\vspace{2mm}
%beginitem
6. Написати програму, що запитує у користувача значення дійсних параметрів $x$ та $y$, обчислює для них значення виразу та повiдомляє введенi данi та результат обчислення.
$$\frac{\log_{10} x + 21}{(\sin x - 0.5) (y + 81)}$$ 

%enditem
\par
%beginitem
7. Написати функцію, що в інтервалі $[a, b)$ знаходить ціле число, що ділиться на найбільший степінь числа 3. Якщо таких чисел декілька, повернути найменше.

По завданням 2 та 3 оцінюється правильність та якість коду, а також економність обчислень.

%enditem
}
{\smallskip\smallskip \begin {scriptsize} Затверджено на засіданні кафедри теоретичної кібернетики, протокол \No 3  від   08.11.2025 р.\par\smallskip\smallskip\smallskip\smallskip
{\parbox {6.5 cm} {Зав. кафедри \dotfill Ю.В.~Крак}}\hspace {0.7cm} {\hbox to 6.5 cm { Екзаменатор \dotfill Т.О.~Карнаух}} \end{scriptsize}}
%end
\newpage
%begin
{{\begin {scriptsize}\par
{ \center  Київський національний університет імені Тараса Шевченка \par}
{\parbox {5 cm}{Освітня програма \hfill}{\parbox {7 cm}{Інформатика\hfill}}\parbox {6 cm}{\hfill Семестр 1}}\par
{\parbox {5 cm} {Навчальний предмет \hfill}\parbox {7 cm}{Програмування \hfill}\parbox {6cm}{\hfill}\par}\vspace{-15pt} \end{scriptsize}}{\center Екзаменаційний білет \No \textbf { 69 }\par}}
{%beginitem
1. Які з наведених типів не є цілими (integral), якщо такі є серед перелічених?

char, double, int, float, unsigned*?
%enditem
\par
\vspace{2mm}
%beginitem
2. Скільки змінних, що є вказівниками, тут оголошено?
char *s1, pc, **s2;?
%enditem
\par
\vspace{2mm}
%beginitem
3. За наявних оголошень \code{void f(int); int n;} які з виразів \code{f(3), f(++n), f(\&n)} є синтаксично коректними?
%enditem
\par
\vspace{2mm}
%beginitem
4. Що таке просторова складність алгоритму?
%enditem
\par
\vspace{2mm}
%beginitem
5. Яку часову оцінку складності має алгоритм сортування бульбашкою?
%enditem
\par
\vspace{2mm}
%beginitem
6. Написати програму, що запитує у користувача значення дійсних параметрів $x$ та $y$, обчислює для них значення виразу та повiдомляє введенi данi та результат обчислення.
$$\frac{\pi x + 40}{(\arccos x + 0.9) (y + 94)}$$ 

%enditem
\par
%beginitem
7. Написати функцію, що для послідовності цілих (задається масивом) знаходить найдовший відрізок, в якому всі сусідні числа відрізняються якнайменш на 3. Якщо таких відрізків кілька, то повернути перший (як індекс початку та довжину). Використовувати додаткові структури даних (у тому числі рядки) заборонено.

По завданням 2 та 3 оцінюється правильність та якість коду, а також економність обчислень.

%enditem
}
{\smallskip\smallskip \begin {scriptsize} Затверджено на засіданні кафедри теоретичної кібернетики, протокол \No 3  від   08.11.2025 р.\par\smallskip\smallskip\smallskip\smallskip
{\parbox {6.5 cm} {Зав. кафедри \dotfill Ю.В.~Крак}}\hspace {0.7cm} {\hbox to 6.5 cm { Екзаменатор \dotfill Т.О.~Карнаух}} \end{scriptsize}}
%end
\vspace{0.9cm}
%begin
{{\begin {scriptsize}\par
{ \center  Київський національний університет імені Тараса Шевченка \par}
{\parbox {5 cm}{Освітня програма \hfill}{\parbox {7 cm}{Інформатика\hfill}}\parbox {6 cm}{\hfill Семестр 1}}\par
{\parbox {5 cm} {Навчальний предмет \hfill}\parbox {7 cm}{Програмування \hfill}\parbox {6cm}{\hfill}\par}\vspace{-15pt} \end{scriptsize}}{\center Екзаменаційний білет \No \textbf { 70 }\par}}
{%beginitem
1. Які з наведених типів не є цілими (integral), якщо такі є серед перелічених?

char, double, int, float, unsigned?
%enditem
\par
\vspace{2mm}
%beginitem
2. Скільки змінних, що є вказівниками, тут оголошено?
int n, m, *p1, p2;?
%enditem
\par
\vspace{2mm}
%beginitem
3. За наявних оголошень \code{void f(char*); char c;} які з виразів \code{f(c*c), f(c+=2), f(\&c)} є синтаксично коректними?
%enditem
\par
\vspace{2mm}
%beginitem
4. Що таке часова складність алгоритму?
%enditem
\par
\vspace{2mm}
%beginitem
5. Яку часову оцінку складності має алгоритм швидкого сортування?
%enditem
\par
\vspace{2mm}
%beginitem
6. Написати програму, що запитує у користувача значення дійсних параметрів $x$ та $y$, обчислює для них значення виразу та повiдомляє введенi данi та результат обчислення.
$$\frac{\log_{2} x - 85}{(\ctg x - 0.7) (y - 77)}$$ 

%enditem
\par
%beginitem
7. Написати функцію, що для послідовності цілих (задається масивом) знаходить найдовший відрізок, в якому всі сусідні числа різні. Якщо таких відрізків кілька, то повернути останній (як індекс початку та довжину). Використовувати додаткові структури даних (у тому числі рядки) заборонено.

По завданням 2 та 3 оцінюється правильність та якість коду, а також економність обчислень.

%enditem
}
{\smallskip\smallskip \begin {scriptsize} Затверджено на засіданні кафедри теоретичної кібернетики, протокол \No 3  від   08.11.2025 р.\par\smallskip\smallskip\smallskip\smallskip
{\parbox {6.5 cm} {Зав. кафедри \dotfill Ю.В.~Крак}}\hspace {0.7cm} {\hbox to 6.5 cm { Екзаменатор \dotfill Т.О.~Карнаух}} \end{scriptsize}}
%end
\newpage
%begin
{{\begin {scriptsize}\par
{ \center  Київський національний університет імені Тараса Шевченка \par}
{\parbox {5 cm}{Освітня програма \hfill}{\parbox {7 cm}{Інформатика\hfill}}\parbox {6 cm}{\hfill Семестр 1}}\par
{\parbox {5 cm} {Навчальний предмет \hfill}\parbox {7 cm}{Програмування \hfill}\parbox {6cm}{\hfill}\par}\vspace{-15pt} \end{scriptsize}}{\center Екзаменаційний білет \No \textbf { 71 }\par}}
{%beginitem
1. Які з наведених типів не є цілими (integral), якщо такі є серед перелічених?

char, double, int, float, unsigned*?
%enditem
\par
\vspace{2mm}
%beginitem
2. Скільки змінних, що є вказівниками, тут оголошено?
int *pn1, pn2, pn3, pn4;?
%enditem
\par
\vspace{2mm}
%beginitem
3. За наявних оголошень \code{void f(double*); double x;} які з виразів \code{f(x*x), f(x+=2), f(\&x)} є синтаксично коректними?
%enditem
\par
\vspace{2mm}
%beginitem
4. Що таке просторова складність задачі?
%enditem
\par
\vspace{2mm}
%beginitem
5. Яку часову оцінку складності має алгоритм сортування купою (він же деревом, пірамідою)?
%enditem
\par
\vspace{2mm}
%beginitem
6. Написати програму, що запитує у користувача значення дійсних параметрів $x$ та $y$, обчислює для них значення виразу та повiдомляє введенi данi та результат обчислення.
$$\frac{\ x - 16}{(\tg x - 0.4) (y + 83)}$$ 

%enditem
\par
%beginitem
7. Написати функцію, що для інтервалу $[a, b)$ знаходить найближчу пару чисел з цього інтервалу, що мають від трьох до п'яти (включно) простих дільників.

По завданням 2 та 3 оцінюється правильність та якість коду, а також економність обчислень.

%enditem
}
{\smallskip\smallskip \begin {scriptsize} Затверджено на засіданні кафедри теоретичної кібернетики, протокол \No 3  від   08.11.2025 р.\par\smallskip\smallskip\smallskip\smallskip
{\parbox {6.5 cm} {Зав. кафедри \dotfill Ю.В.~Крак}}\hspace {0.7cm} {\hbox to 6.5 cm { Екзаменатор \dotfill Т.О.~Карнаух}} \end{scriptsize}}
%end
\vspace{0.9cm}
%begin
{{\begin {scriptsize}\par
{ \center  Київський національний університет імені Тараса Шевченка \par}
{\parbox {5 cm}{Освітня програма \hfill}{\parbox {7 cm}{Інформатика\hfill}}\parbox {6 cm}{\hfill Семестр 1}}\par
{\parbox {5 cm} {Навчальний предмет \hfill}\parbox {7 cm}{Програмування \hfill}\parbox {6cm}{\hfill}\par}\vspace{-15pt} \end{scriptsize}}{\center Екзаменаційний білет \No \textbf { 72 }\par}}
{%beginitem
1. Які з наведених типів не є цілими (integral), якщо такі є серед перелічених?

bool, int, long, long long, unsigned char?
%enditem
\par
\vspace{2mm}
%beginitem
2. Скільки змінних, що є вказівниками, тут оголошено?
double x, *px1, px2, ***px3;?
%enditem
\par
\vspace{2mm}
%beginitem
3. За наявних оголошень \code{void f(int*); int n;} які з виразів \code{f(n++), f(++n), f(\&n)} є синтаксично коректними?
%enditem
\par
\vspace{2mm}
%beginitem
4. У чому відмінність літералів від констант?
%enditem
\par
\vspace{2mm}
%beginitem
5. Яку часову оцінку складності має алгоритм сортування бульбашкою?
%enditem
\par
\vspace{2mm}
%beginitem
6. Написати програму, що запитує у користувача значення дійсних параметрів $x$ та $y$, обчислює для них значення виразу та повiдомляє введенi данi та результат обчислення.
$$\frac{\ln x + 29}{(\ctg x + 0.6) (y - 31)}$$ 

%enditem
\par
%beginitem
7. Написати функцію, що для інтервалу $[a, b)$ знаходить найближчу пару чисел з цього інтервалу, що кожне з них має непарну кількість різних простих дільників. Якщо таких пар кілька, то цікавить пара з найбільшою сумою.

По завданням 2 та 3 оцінюється правильність та якість коду, а також економність обчислень.

%enditem
}
{\smallskip\smallskip \begin {scriptsize} Затверджено на засіданні кафедри теоретичної кібернетики, протокол \No 3  від   08.11.2025 р.\par\smallskip\smallskip\smallskip\smallskip
{\parbox {6.5 cm} {Зав. кафедри \dotfill Ю.В.~Крак}}\hspace {0.7cm} {\hbox to 6.5 cm { Екзаменатор \dotfill Т.О.~Карнаух}} \end{scriptsize}}
%end
\newpage
%begin
{{\begin {scriptsize}\par
{ \center  Київський національний університет імені Тараса Шевченка \par}
{\parbox {5 cm}{Освітня програма \hfill}{\parbox {7 cm}{Інформатика\hfill}}\parbox {6 cm}{\hfill Семестр 1}}\par
{\parbox {5 cm} {Навчальний предмет \hfill}\parbox {7 cm}{Програмування \hfill}\parbox {6cm}{\hfill}\par}\vspace{-15pt} \end{scriptsize}}{\center Екзаменаційний білет \No \textbf { 73 }\par}}
{%beginitem
1. Які з наведених типів не є цілими (integral), якщо такі є серед перелічених?

int*, signed char, unsigned long, string?
%enditem
\par
\vspace{2mm}
%beginitem
2. Скільки змінних, що є вказівниками, тут оголошено?
long pointer1, pointer2, pointer3, pointer4;?
%enditem
\par
\vspace{2mm}
%beginitem
3. За наявних оголошень \code{void f(double\&); double x;} які з виразів \code{f(x*x), f(x+=2), f(\&x)} є синтаксично коректними?
%enditem
\par
\vspace{2mm}
%beginitem
4. Що таке часова складність задачі?
%enditem
\par
\vspace{2mm}
%beginitem
5. Яку часову оцінку складності має алгоритм швидкого сортування?
%enditem
\par
\vspace{2mm}
%beginitem
6. Написати програму, що запитує у користувача значення дійсних параметрів $x$ та $y$, обчислює для них значення виразу та повiдомляє введенi данi та результат обчислення.
$$\frac{\log_{10} x + 92}{(\arctg x - 0.3) (y + 73)}$$ 

%enditem
\par
%beginitem
7. Написати функцію, що знаходить найбільше число з послідовності цілих (задається масивом), у десятковому запису якого нема цифри 3. Повертає індекс останнього входження цього числа або $-1$ (за відсутності). Використовувати додаткові структури даних (у тому числі рядки) заборонено.

По завданням 2 та 3 оцінюється правильність та якість коду, а також економність обчислень.

%enditem
}
{\smallskip\smallskip \begin {scriptsize} Затверджено на засіданні кафедри теоретичної кібернетики, протокол \No 3  від   08.11.2025 р.\par\smallskip\smallskip\smallskip\smallskip
{\parbox {6.5 cm} {Зав. кафедри \dotfill Ю.В.~Крак}}\hspace {0.7cm} {\hbox to 6.5 cm { Екзаменатор \dotfill Т.О.~Карнаух}} \end{scriptsize}}
%end
\vspace{0.9cm}
%begin
{{\begin {scriptsize}\par
{ \center  Київський національний університет імені Тараса Шевченка \par}
{\parbox {5 cm}{Освітня програма \hfill}{\parbox {7 cm}{Інформатика\hfill}}\parbox {6 cm}{\hfill Семестр 1}}\par
{\parbox {5 cm} {Навчальний предмет \hfill}\parbox {7 cm}{Програмування \hfill}\parbox {6cm}{\hfill}\par}\vspace{-15pt} \end{scriptsize}}{\center Екзаменаційний білет \No \textbf { 74 }\par}}
{%beginitem
1. Які з наведених типів не є цілими (integral), якщо такі є серед перелічених?

double, int, long, unsigned char?
%enditem
\par
\vspace{2mm}
%beginitem
2. Скільки змінних, що є вказівниками, тут оголошено?
char *s1, pc, **s2;?
%enditem
\par
\vspace{2mm}
%beginitem
3. За наявних оголошень \code{void f(char); char c;} які з виразів \code{f(c*c), f(c+=2), f(\&c)} є синтаксично коректними?
%enditem
\par
\vspace{2mm}
%beginitem
4. Чим префіксний оператор ++ відрізняється від постфіксного?
%enditem
\par
\vspace{2mm}
%beginitem
5. Яку часову оцінку складності має алгоритм сортування вставками?
%enditem
\par
\vspace{2mm}
%beginitem
6. Написати програму, що запитує у користувача значення дійсних параметрів $x$ та $y$, обчислює для них значення виразу та повiдомляє введенi данi та результат обчислення.
$$\frac{\sqrt x - 90}{(\cos x + 0.6) (y - 17)}$$ 

%enditem
\par
%beginitem
7. Написати функцію, що для цілого числа знаходить простий дільник, який входить у його розклад у найбільшому степені. Якщо таких кілька, то повернути найменший. Наприклад, для числа $2^{9}{\cdot}3^5{\cdot}7^9{\cdot}11$ функція має повернути $2$.

По завданням 2 та 3 оцінюється правильність та якість коду, а також економність обчислень.

%enditem
}
{\smallskip\smallskip \begin {scriptsize} Затверджено на засіданні кафедри теоретичної кібернетики, протокол \No 3  від   08.11.2025 р.\par\smallskip\smallskip\smallskip\smallskip
{\parbox {6.5 cm} {Зав. кафедри \dotfill Ю.В.~Крак}}\hspace {0.7cm} {\hbox to 6.5 cm { Екзаменатор \dotfill Т.О.~Карнаух}} \end{scriptsize}}
%end
\newpage
%begin
{{\begin {scriptsize}\par
{ \center  Київський національний університет імені Тараса Шевченка \par}
{\parbox {5 cm}{Освітня програма \hfill}{\parbox {7 cm}{Інформатика\hfill}}\parbox {6 cm}{\hfill Семестр 1}}\par
{\parbox {5 cm} {Навчальний предмет \hfill}\parbox {7 cm}{Програмування \hfill}\parbox {6cm}{\hfill}\par}\vspace{-15pt} \end{scriptsize}}{\center Екзаменаційний білет \No \textbf { 75 }\par}}
{%beginitem
1. Які з наведених типів не є цілими (integral), якщо такі є серед перелічених?

double, int, long*, unsigned char?
%enditem
\par
\vspace{2mm}
%beginitem
2. Скільки змінних, що є вказівниками, тут оголошено?
double *px, *py, **ppx;?
%enditem
\par
\vspace{2mm}
%beginitem
3. За наявних оголошень \code{void f(double); double x;} які з виразів \code{f(3.1), f(x+=2), f(\&x)} є синтаксично коректними?
%enditem
\par
\vspace{2mm}
%beginitem
4. Що означає термін ``часова складність у середньому''?
%enditem
\par
\vspace{2mm}
%beginitem
5. Яку часову оцінку складності має алгоритм швидкого сортування?
%enditem
\par
\vspace{2mm}
%beginitem
6. Написати програму, що запитує у користувача значення дійсних параметрів $x$ та $y$, обчислює для них значення виразу та повiдомляє введенi данi та результат обчислення.
$$\frac{\ln x - 33}{(\arcsin x - 0.4) (y + 91)}$$ 

%enditem
\par
%beginitem
7. Написати функцію, що знаходить найбільше число з послідовності цілих (задається масивом), у десятковому запису якого нема цифри 4. Повертає індекс першого входження цього числа або $-1$ (за відсутності). Використовувати додаткові структури даних (у тому числі рядки) заборонено.

По завданням 2 та 3 оцінюється правильність та якість коду, а також економність обчислень.

%enditem
}
{\smallskip\smallskip \begin {scriptsize} Затверджено на засіданні кафедри теоретичної кібернетики, протокол \No 3  від   08.11.2025 р.\par\smallskip\smallskip\smallskip\smallskip
{\parbox {6.5 cm} {Зав. кафедри \dotfill Ю.В.~Крак}}\hspace {0.7cm} {\hbox to 6.5 cm { Екзаменатор \dotfill Т.О.~Карнаух}} \end{scriptsize}}
%end
\vspace{0.9cm}
%begin
{{\begin {scriptsize}\par
{ \center  Київський національний університет імені Тараса Шевченка \par}
{\parbox {5 cm}{Освітня програма \hfill}{\parbox {7 cm}{Інформатика\hfill}}\parbox {6 cm}{\hfill Семестр 1}}\par
{\parbox {5 cm} {Навчальний предмет \hfill}\parbox {7 cm}{Програмування \hfill}\parbox {6cm}{\hfill}\par}\vspace{-15pt} \end{scriptsize}}{\center Екзаменаційний білет \No \textbf { 76 }\par}}
{%beginitem
1. Які з наведених типів не є цілими (integral), якщо такі є серед перелічених?

int, signed char, unsigned long, string?
%enditem
\par
\vspace{2mm}
%beginitem
2. Скільки змінних, що є вказівниками, тут оголошено?
long *n1, *n2, n3, pn4;?
%enditem
\par
\vspace{2mm}
%beginitem
3. За наявних оголошень \code{void f(char\&); char c;} які з виразів \code{f(c*c), f(c++), f(\&c)} є синтаксично коректними?
%enditem
\par
\vspace{2mm}
%beginitem
4. Що таке тип даних?
%enditem
\par
\vspace{2mm}
%beginitem
5. Яку часову оцінку складності має алгоритм сортування вибором?
%enditem
\par
\vspace{2mm}
%beginitem
6. Написати програму, що запитує у користувача значення дійсних параметрів $x$ та $y$, обчислює для них значення виразу та повiдомляє введенi данi та результат обчислення.
$$\frac{\ x + 65}{(\arccos x + 0.7) (y - 55)}$$ 

%enditem
\par
%beginitem
7. Написати функцію, що знаходить найменше число з послідовності цілих (задається масивом), у десятковому запису якого нема цифри 7. Повертає індекс першого входження цього числа або $-1$ (за відсутності). Використовувати додаткові структури даних (у тому числі рядки) заборонено.

По завданням 2 та 3 оцінюється правильність та якість коду, а також економність обчислень.

%enditem
}
{\smallskip\smallskip \begin {scriptsize} Затверджено на засіданні кафедри теоретичної кібернетики, протокол \No 3  від   08.11.2025 р.\par\smallskip\smallskip\smallskip\smallskip
{\parbox {6.5 cm} {Зав. кафедри \dotfill Ю.В.~Крак}}\hspace {0.7cm} {\hbox to 6.5 cm { Екзаменатор \dotfill Т.О.~Карнаух}} \end{scriptsize}}
%end
\newpage
%begin
{{\begin {scriptsize}\par
{ \center  Київський національний університет імені Тараса Шевченка \par}
{\parbox {5 cm}{Освітня програма \hfill}{\parbox {7 cm}{Інформатика\hfill}}\parbox {6 cm}{\hfill Семестр 1}}\par
{\parbox {5 cm} {Навчальний предмет \hfill}\parbox {7 cm}{Програмування \hfill}\parbox {6cm}{\hfill}\par}\vspace{-15pt} \end{scriptsize}}{\center Екзаменаційний білет \No \textbf { 77 }\par}}
{%beginitem
1. Які з наведених типів не є цілими (integral), якщо такі є серед перелічених?

bool, char, int, unsigned?
%enditem
\par
\vspace{2mm}
%beginitem
2. Скільки змінних, що є вказівниками, тут оголошено?
int n, m, *p1, p2;?
%enditem
\par
\vspace{2mm}
%beginitem
3. За наявних оголошень \code{void f(int); int n;} які з виразів \code{f(3), f(++n), f(\&n)} є синтаксично коректними?
%enditem
\par
\vspace{2mm}
%beginitem
4. Які параметри може приймати функція main? Що вони позначають та яких вони типів?
%enditem
\par
\vspace{2mm}
%beginitem
5. Яку часову оцінку складності має алгоритм сортування злиттям?
%enditem
\par
\vspace{2mm}
%beginitem
6. Написати програму, що запитує у користувача значення дійсних параметрів $x$ та $y$, обчислює для них значення виразу та повiдомляє введенi данi та результат обчислення.
$$\frac{\pi x + 18}{(\sin x - 0.8) (y + 98)}$$ 

%enditem
\par
%beginitem
7. Написати функцію, що для цілого числа знаходить кількість простих дільників, які входять у його розклад у найбільшому степені. Наприклад, для числа $3^5{\cdot}7^9{\cdot}11{\cdot}23^{9}$ таких дільників два – $23$ та $7$, тому функція має повернути $2$.

По завданням 2 та 3 оцінюється правильність та якість коду, а також економність обчислень.

%enditem
}
{\smallskip\smallskip \begin {scriptsize} Затверджено на засіданні кафедри теоретичної кібернетики, протокол \No 3  від   08.11.2025 р.\par\smallskip\smallskip\smallskip\smallskip
{\parbox {6.5 cm} {Зав. кафедри \dotfill Ю.В.~Крак}}\hspace {0.7cm} {\hbox to 6.5 cm { Екзаменатор \dotfill Т.О.~Карнаух}} \end{scriptsize}}
%end
\vspace{0.9cm}
%begin
{{\begin {scriptsize}\par
{ \center  Київський національний університет імені Тараса Шевченка \par}
{\parbox {5 cm}{Освітня програма \hfill}{\parbox {7 cm}{Інформатика\hfill}}\parbox {6 cm}{\hfill Семестр 1}}\par
{\parbox {5 cm} {Навчальний предмет \hfill}\parbox {7 cm}{Програмування \hfill}\parbox {6cm}{\hfill}\par}\vspace{-15pt} \end{scriptsize}}{\center Екзаменаційний білет \No \textbf { 78 }\par}}
{%beginitem
1. Які з наведених типів не є цілими (integral), якщо такі є серед перелічених?

char, double, int, float, unsigned?
%enditem
\par
\vspace{2mm}
%beginitem
2. Скільки змінних, що є вказівниками, тут оголошено?
int *pn1, pn2, pn3, pn4;?
%enditem
\par
\vspace{2mm}
%beginitem
3. За наявних оголошень \code{void f(int\&); int n;} які з виразів \code{f(13), f(++n), f(\&n)} є синтаксично коректними?
%enditem
\par
\vspace{2mm}
%beginitem
4. Що таке С-рядок?
%enditem
\par
\vspace{2mm}
%beginitem
5. Яку часову оцінку складності має алгоритм сортування купою (він же деревом, пірамідою)?
%enditem
\par
\vspace{2mm}
%beginitem
6. Написати програму, що запитує у користувача значення дійсних параметрів $x$ та $y$, обчислює для них значення виразу та повiдомляє введенi данi та результат обчислення.
$$\frac{\log_{2} x - 13}{(\arctg x + 0.2) (y - 45)}$$ 

%enditem
\par
%beginitem
7. Написати функцію, що для цілого числа знаходить кількість його простих дільників, що входять у його розклад рівно у другому степеню. Наприклад, для числа $3^2{\cdot}7^9{\cdot}11{\cdot}23^{2}$ таких дільників два – $3$ та $23$, тому функція має повернути $2$.

По завданням 2 та 3 оцінюється правильність та якість коду, а також економність обчислень.

%enditem
}
{\smallskip\smallskip \begin {scriptsize} Затверджено на засіданні кафедри теоретичної кібернетики, протокол \No 3  від   08.11.2025 р.\par\smallskip\smallskip\smallskip\smallskip
{\parbox {6.5 cm} {Зав. кафедри \dotfill Ю.В.~Крак}}\hspace {0.7cm} {\hbox to 6.5 cm { Екзаменатор \dotfill Т.О.~Карнаух}} \end{scriptsize}}
%end
\newpage
%begin
{{\begin {scriptsize}\par
{ \center  Київський національний університет імені Тараса Шевченка \par}
{\parbox {5 cm}{Освітня програма \hfill}{\parbox {7 cm}{Інформатика\hfill}}\parbox {6 cm}{\hfill Семестр 1}}\par
{\parbox {5 cm} {Навчальний предмет \hfill}\parbox {7 cm}{Програмування \hfill}\parbox {6cm}{\hfill}\par}\vspace{-15pt} \end{scriptsize}}{\center Екзаменаційний білет \No \textbf { 79 }\par}}
{%beginitem
1. Які з наведених типів не є цілими (integral), якщо такі є серед перелічених?

bool, int, long, long long, unsigned char?
%enditem
\par
\vspace{2mm}
%beginitem
2. Скільки змінних, що є вказівниками, тут оголошено?
long pointer1, pointer2, pointer3, pointer4;?
%enditem
\par
\vspace{2mm}
%beginitem
3. За наявних оголошень \code{void f(char*); char c;} які з виразів \code{f(c*c), f(c+=2), f(\&c)} є синтаксично коректними?
%enditem
\par
\vspace{2mm}
%beginitem
4. У чому відмінність літералів від констант?
%enditem
\par
\vspace{2mm}
%beginitem
5. Яку часову оцінку складності має алгоритм швидкого сортування?
%enditem
\par
\vspace{2mm}
%beginitem
6. Написати програму, що запитує у користувача значення дійсних параметрів $x$ та $y$, обчислює для них значення виразу та повiдомляє введенi данi та результат обчислення.
$$\frac{\pi x + 80}{(\tg x - 0.9) (y + 11)}$$ 

%enditem
\par
%beginitem
7. Написати функцію, що для послідовності цілих (задається масивом) знаходить найдовший відрізок, в якому всі сусідні числа різні. Якщо таких відрізків кілька, то повернути перший (як індекс початку та довжину). Використовувати додаткові структури даних (у тому числі рядки) заборонено.

По завданням 2 та 3 оцінюється правильність та якість коду, а також економність обчислень.

%enditem
}
{\smallskip\smallskip \begin {scriptsize} Затверджено на засіданні кафедри теоретичної кібернетики, протокол \No 3  від   08.11.2025 р.\par\smallskip\smallskip\smallskip\smallskip
{\parbox {6.5 cm} {Зав. кафедри \dotfill Ю.В.~Крак}}\hspace {0.7cm} {\hbox to 6.5 cm { Екзаменатор \dotfill Т.О.~Карнаух}} \end{scriptsize}}
%end
\vspace{0.9cm}
%begin
{{\begin {scriptsize}\par
{ \center  Київський національний університет імені Тараса Шевченка \par}
{\parbox {5 cm}{Освітня програма \hfill}{\parbox {7 cm}{Інформатика\hfill}}\parbox {6 cm}{\hfill Семестр 1}}\par
{\parbox {5 cm} {Навчальний предмет \hfill}\parbox {7 cm}{Програмування \hfill}\parbox {6cm}{\hfill}\par}\vspace{-15pt} \end{scriptsize}}{\center Екзаменаційний білет \No \textbf { 80 }\par}}
{%beginitem
1. Які з наведених типів не є цілими (integral), якщо такі є серед перелічених?

bool*, char, int, unsigned?
%enditem
\par
\vspace{2mm}
%beginitem
2. Скільки змінних, що є вказівниками, тут оголошено?
int n, m, *p1, p2;?
%enditem
\par
\vspace{2mm}
%beginitem
3. За наявних оголошень \code{void f(double*); double x;} які з виразів \code{f(x*x), f(x+=2), f(\&x)} є синтаксично коректними?
%enditem
\par
\vspace{2mm}
%beginitem
4. Що таке часова складність задачі?
%enditem
\par
\vspace{2mm}
%beginitem
5. Яку часову оцінку складності має алгоритм сортування вставками?
%enditem
\par
\vspace{2mm}
%beginitem
6. Написати програму, що запитує у користувача значення дійсних параметрів $x$ та $y$, обчислює для них значення виразу та повiдомляє введенi данi та результат обчислення.
$$\frac{\log_{10} x - 79}{(\arccos x + 0.1) (y - 97)}$$ 

%enditem
\par
%beginitem
7. Написати функцію, що в інтервалі $[a, b)$ знаходить ціле число, що ділиться на найбільший степінь числа 3. Якщо таких чисел декілька, повернути найбільше.

По завданням 2 та 3 оцінюється правильність та якість коду, а також економність обчислень.

%enditem
}
{\smallskip\smallskip \begin {scriptsize} Затверджено на засіданні кафедри теоретичної кібернетики, протокол \No 3  від   08.11.2025 р.\par\smallskip\smallskip\smallskip\smallskip
{\parbox {6.5 cm} {Зав. кафедри \dotfill Ю.В.~Крак}}\hspace {0.7cm} {\hbox to 6.5 cm { Екзаменатор \dotfill Т.О.~Карнаух}} \end{scriptsize}}
%end
\newpage
%begin
{{\begin {scriptsize}\par
{ \center  Київський національний університет імені Тараса Шевченка \par}
{\parbox {5 cm}{Освітня програма \hfill}{\parbox {7 cm}{Інформатика\hfill}}\parbox {6 cm}{\hfill Семестр 1}}\par
{\parbox {5 cm} {Навчальний предмет \hfill}\parbox {7 cm}{Програмування \hfill}\parbox {6cm}{\hfill}\par}\vspace{-15pt} \end{scriptsize}}{\center Екзаменаційний білет \No \textbf { 81 }\par}}
{%beginitem
1. Які з наведених типів не є цілими (integral), якщо такі є серед перелічених?

int, signed char, unsigned long, string?
%enditem
\par
\vspace{2mm}
%beginitem
2. Скільки змінних, що є вказівниками, тут оголошено?
long *n1, *n2, n3, pn4;?
%enditem
\par
\vspace{2mm}
%beginitem
3. За наявних оголошень \code{void f(int*); int n;} які з виразів \code{f(n++), f(++n), f(\&n)} є синтаксично коректними?
%enditem
\par
\vspace{2mm}
%beginitem
4. Що таке просторова складність алгоритму?
%enditem
\par
\vspace{2mm}
%beginitem
5. Яку часову оцінку складності має алгоритм сортування бульбашкою?
%enditem
\par
\vspace{2mm}
%beginitem
6. Написати програму, що запитує у користувача значення дійсних параметрів $x$ та $y$, обчислює для них значення виразу та повiдомляє введенi данi та результат обчислення.
$$\frac{\ln x - 20}{(\sin x + 0.5) (y + 54)}$$ 

%enditem
\par
%beginitem
7. Написати функцію, що для інтервалу $[a, b)$ знаходить найближчу пару чисел з цього інтервалу, що кожне з них має парну кількість різних простих дільників. Якщо таких пар кілька, то цікавить пара з найменшим добутком.

По завданням 2 та 3 оцінюється правильність та якість коду, а також економність обчислень.

%enditem
}
{\smallskip\smallskip \begin {scriptsize} Затверджено на засіданні кафедри теоретичної кібернетики, протокол \No 3  від   08.11.2025 р.\par\smallskip\smallskip\smallskip\smallskip
{\parbox {6.5 cm} {Зав. кафедри \dotfill Ю.В.~Крак}}\hspace {0.7cm} {\hbox to 6.5 cm { Екзаменатор \dotfill Т.О.~Карнаух}} \end{scriptsize}}
%end
\vspace{0.9cm}
%begin
{{\begin {scriptsize}\par
{ \center  Київський національний університет імені Тараса Шевченка \par}
{\parbox {5 cm}{Освітня програма \hfill}{\parbox {7 cm}{Інформатика\hfill}}\parbox {6 cm}{\hfill Семестр 1}}\par
{\parbox {5 cm} {Навчальний предмет \hfill}\parbox {7 cm}{Програмування \hfill}\parbox {6cm}{\hfill}\par}\vspace{-15pt} \end{scriptsize}}{\center Екзаменаційний білет \No \textbf { 82 }\par}}
{%beginitem
1. Які з наведених типів не є цілими (integral), якщо такі є серед перелічених?

bool, int, long, long long, unsigned char?
%enditem
\par
\vspace{2mm}
%beginitem
2. Скільки змінних, що є вказівниками, тут оголошено?
double x, *px1, px2, ***px3;?
%enditem
\par
\vspace{2mm}
%beginitem
3. За наявних оголошень \code{void f(double*); double x;} які з виразів \code{f(x*x), f(x+=2), f(\&x)} є синтаксично коректними?
%enditem
\par
\vspace{2mm}
%beginitem
4. Що таке просторова складність задачі?
%enditem
\par
\vspace{2mm}
%beginitem
5. Яку часову оцінку складності має алгоритм сортування вибором?
%enditem
\par
\vspace{2mm}
%beginitem
6. Написати програму, що запитує у користувача значення дійсних параметрів $x$ та $y$, обчислює для них значення виразу та повiдомляє введенi данi та результат обчислення.
$$\frac{\sqrt x + 25}{(\arcsin x - 0.1) (y - 71)}$$ 

%enditem
\par
%beginitem
7. Написати функцію, що повертає найменше число з послідовності цілих (задається масивом), що має не більше трьох різних простих дільників. За відсутності має повертати None. Використовувати додаткові структури даних (у тому числі рядки) заборонено.

По завданням 2 та 3 оцінюється правильність та якість коду, а також економність обчислень.

%enditem
}
{\smallskip\smallskip \begin {scriptsize} Затверджено на засіданні кафедри теоретичної кібернетики, протокол \No 3  від   08.11.2025 р.\par\smallskip\smallskip\smallskip\smallskip
{\parbox {6.5 cm} {Зав. кафедри \dotfill Ю.В.~Крак}}\hspace {0.7cm} {\hbox to 6.5 cm { Екзаменатор \dotfill Т.О.~Карнаух}} \end{scriptsize}}
%end
\newpage
%begin
{{\begin {scriptsize}\par
{ \center  Київський національний університет імені Тараса Шевченка \par}
{\parbox {5 cm}{Освітня програма \hfill}{\parbox {7 cm}{Інформатика\hfill}}\parbox {6 cm}{\hfill Семестр 1}}\par
{\parbox {5 cm} {Навчальний предмет \hfill}\parbox {7 cm}{Програмування \hfill}\parbox {6cm}{\hfill}\par}\vspace{-15pt} \end{scriptsize}}{\center Екзаменаційний білет \No \textbf { 83 }\par}}
{%beginitem
1. Які з наведених типів не є цілими (integral), якщо такі є серед перелічених?

bool, int, long, long long, unsigned char?
%enditem
\par
\vspace{2mm}
%beginitem
2. Скільки змінних, що є вказівниками, тут оголошено?
char *s1, pc, **s2;?
%enditem
\par
\vspace{2mm}
%beginitem
3. За наявних оголошень \code{void f(int\&); int n;} які з виразів \code{f(13), f(++n), f(\&n)} є синтаксично коректними?
%enditem
\par
\vspace{2mm}
%beginitem
4. Що означає термін ``часова складність у середньому''?
%enditem
\par
\vspace{2mm}
%beginitem
5. Яку часову оцінку складності має алгоритм швидкого сортування?
%enditem
\par
\vspace{2mm}
%beginitem
6. Написати програму, що запитує у користувача значення дійсних параметрів $x$ та $y$, обчислює для них значення виразу та повiдомляє введенi данi та результат обчислення.
$$\frac{\ x + 28}{(\ctg x - 0.8) (y + 75)}$$ 

%enditem
\par
%beginitem
7. Написати функцію, що повертає число з послідовності цілих (задається масивом), що має найбільшу кількість різних простих дільників. Якщо таких чисел кілька, то повернути найменше. Використовувати додаткові структури даних (у тому числі рядки) заборонено.

По завданням 2 та 3 оцінюється правильність та якість коду, а також економність обчислень.

%enditem
}
{\smallskip\smallskip \begin {scriptsize} Затверджено на засіданні кафедри теоретичної кібернетики, протокол \No 3  від   08.11.2025 р.\par\smallskip\smallskip\smallskip\smallskip
{\parbox {6.5 cm} {Зав. кафедри \dotfill Ю.В.~Крак}}\hspace {0.7cm} {\hbox to 6.5 cm { Екзаменатор \dotfill Т.О.~Карнаух}} \end{scriptsize}}
%end
\vspace{0.9cm}
%begin
{{\begin {scriptsize}\par
{ \center  Київський національний університет імені Тараса Шевченка \par}
{\parbox {5 cm}{Освітня програма \hfill}{\parbox {7 cm}{Інформатика\hfill}}\parbox {6 cm}{\hfill Семестр 1}}\par
{\parbox {5 cm} {Навчальний предмет \hfill}\parbox {7 cm}{Програмування \hfill}\parbox {6cm}{\hfill}\par}\vspace{-15pt} \end{scriptsize}}{\center Екзаменаційний білет \No \textbf { 84 }\par}}
{%beginitem
1. Які з наведених типів не є цілими (integral), якщо такі є серед перелічених?

double, int, long, unsigned char?
%enditem
\par
\vspace{2mm}
%beginitem
2. Скільки змінних, що є вказівниками, тут оголошено?
double *px, *py, **ppx;?
%enditem
\par
\vspace{2mm}
%beginitem
3. За наявних оголошень \code{void f(double\&); double x;} які з виразів \code{f(x*x), f(x+=2), f(\&x)} є синтаксично коректними?
%enditem
\par
\vspace{2mm}
%beginitem
4. Які параметри може приймати функція main? Що вони позначають та яких вони типів?
%enditem
\par
\vspace{2mm}
%beginitem
5. Яку часову оцінку складності має алгоритм сортування злиттям?
%enditem
\par
\vspace{2mm}
%beginitem
6. Написати програму, що запитує у користувача значення дійсних параметрів $x$ та $y$, обчислює для них значення виразу та повiдомляє введенi данi та результат обчислення.
$$\frac{\log_{2} x - 22}{(\cos x + 0.3) (y - 57)}$$ 

%enditem
\par
%beginitem
7. Написати функцію, що для цілого числа знаходить кількість його простих дільників, що входять у його розклад не більш ніж у 5му степеню. Наприклад, для числа $3^8{\cdot}7^9{\cdot}11{\cdot}23^{2}$ таких дільників два – $11$ та $23$, тому функція має повернути $2$.

По завданням 2 та 3 оцінюється правильність та якість коду, а також економність обчислень.

%enditem
}
{\smallskip\smallskip \begin {scriptsize} Затверджено на засіданні кафедри теоретичної кібернетики, протокол \No 3  від   08.11.2025 р.\par\smallskip\smallskip\smallskip\smallskip
{\parbox {6.5 cm} {Зав. кафедри \dotfill Ю.В.~Крак}}\hspace {0.7cm} {\hbox to 6.5 cm { Екзаменатор \dotfill Т.О.~Карнаух}} \end{scriptsize}}
%end
\newpage
%begin
{{\begin {scriptsize}\par
{ \center  Київський національний університет імені Тараса Шевченка \par}
{\parbox {5 cm}{Освітня програма \hfill}{\parbox {7 cm}{Інформатика\hfill}}\parbox {6 cm}{\hfill Семестр 1}}\par
{\parbox {5 cm} {Навчальний предмет \hfill}\parbox {7 cm}{Програмування \hfill}\parbox {6cm}{\hfill}\par}\vspace{-15pt} \end{scriptsize}}{\center Екзаменаційний білет \No \textbf { 85 }\par}}
{%beginitem
1. Які з наведених типів не є цілими (integral), якщо такі є серед перелічених?

char, double, int, float, unsigned?
%enditem
\par
\vspace{2mm}
%beginitem
2. Скільки змінних, що є вказівниками, тут оголошено?
long pointer1, pointer2, pointer3, pointer4;?
%enditem
\par
\vspace{2mm}
%beginitem
3. За наявних оголошень \code{void f(char*); char c;} які з виразів \code{f(c*c), f(c+=2), f(\&c)} є синтаксично коректними?
%enditem
\par
\vspace{2mm}
%beginitem
4. Що таке масив?
%enditem
\par
\vspace{2mm}
%beginitem
5. Яку часову оцінку складності має алгоритм швидкого сортування?
%enditem
\par
\vspace{2mm}
%beginitem
6. Написати програму, що запитує у користувача значення дійсних параметрів $x$ та $y$, обчислює для них значення виразу та повiдомляє введенi данi та результат обчислення.
$$\frac{\pi x + 30}{(\arccos x - 0.4) (y - 23)}$$ 

%enditem
\par
%beginitem
7. Написати функцію, що знаходить найменше число з послідовності цілих (задається масивом), у десятковому запису якого нема цифри 5. Повертає індекс останнього входження цього числа або $-1$ (за відсутності). Використовувати додаткові структури даних (у тому числі рядки) заборонено.

По завданням 2 та 3 оцінюється правильність та якість коду, а також економність обчислень.

%enditem
}
{\smallskip\smallskip \begin {scriptsize} Затверджено на засіданні кафедри теоретичної кібернетики, протокол \No 3  від   08.11.2025 р.\par\smallskip\smallskip\smallskip\smallskip
{\parbox {6.5 cm} {Зав. кафедри \dotfill Ю.В.~Крак}}\hspace {0.7cm} {\hbox to 6.5 cm { Екзаменатор \dotfill Т.О.~Карнаух}} \end{scriptsize}}
%end
\vspace{0.9cm}
%begin
{{\begin {scriptsize}\par
{ \center  Київський національний університет імені Тараса Шевченка \par}
{\parbox {5 cm}{Освітня програма \hfill}{\parbox {7 cm}{Інформатика\hfill}}\parbox {6 cm}{\hfill Семестр 1}}\par
{\parbox {5 cm} {Навчальний предмет \hfill}\parbox {7 cm}{Програмування \hfill}\parbox {6cm}{\hfill}\par}\vspace{-15pt} \end{scriptsize}}{\center Екзаменаційний білет \No \textbf { 86 }\par}}
{%beginitem
1. Які з наведених типів не є цілими (integral), якщо такі є серед перелічених?

double, int, long*, unsigned char?
%enditem
\par
\vspace{2mm}
%beginitem
2. Скільки змінних, що є вказівниками, тут оголошено?
int *pn1, pn2, pn3, pn4;?
%enditem
\par
\vspace{2mm}
%beginitem
3. За наявних оголошень \code{void f(int); int n;} які з виразів \code{f(3), f(++n), f(\&n)} є синтаксично коректними?
%enditem
\par
\vspace{2mm}
%beginitem
4. Що таке тип даних?
%enditem
\par
\vspace{2mm}
%beginitem
5. Яку часову оцінку складності має алгоритм сортування купою (він же деревом, пірамідою)?
%enditem
\par
\vspace{2mm}
%beginitem
6. Написати програму, що запитує у користувача значення дійсних параметрів $x$ та $y$, обчислює для них значення виразу та повiдомляє введенi данi та результат обчислення.
$$\frac{\log_{10} x - 63}{(\arcsin x + 0.7) (y + 34)}$$ 

%enditem
\par
%beginitem
7. Написати функцію, яка для цілого числа знаходить кількість його різних простих дільників, що мають у десятковому запису цифру 1.

По завданням 2 та 3 оцінюється правильність та якість коду, а також економність обчислень.

%enditem
}
{\smallskip\smallskip \begin {scriptsize} Затверджено на засіданні кафедри теоретичної кібернетики, протокол \No 3  від   08.11.2025 р.\par\smallskip\smallskip\smallskip\smallskip
{\parbox {6.5 cm} {Зав. кафедри \dotfill Ю.В.~Крак}}\hspace {0.7cm} {\hbox to 6.5 cm { Екзаменатор \dotfill Т.О.~Карнаух}} \end{scriptsize}}
%end
\newpage
%begin
{{\begin {scriptsize}\par
{ \center  Київський національний університет імені Тараса Шевченка \par}
{\parbox {5 cm}{Освітня програма \hfill}{\parbox {7 cm}{Інформатика\hfill}}\parbox {6 cm}{\hfill Семестр 1}}\par
{\parbox {5 cm} {Навчальний предмет \hfill}\parbox {7 cm}{Програмування \hfill}\parbox {6cm}{\hfill}\par}\vspace{-15pt} \end{scriptsize}}{\center Екзаменаційний білет \No \textbf { 87 }\par}}
{%beginitem
1. Які з наведених типів не є цілими (integral), якщо такі є серед перелічених?

char, double, int, float, unsigned*?
%enditem
\par
\vspace{2mm}
%beginitem
2. Скільки змінних, що є вказівниками, тут оголошено?
double *px, *py, **ppx;?
%enditem
\par
\vspace{2mm}
%beginitem
3. За наявних оголошень \code{void f(char\&); char c;} які з виразів \code{f(c*c), f(c++), f(\&c)} є синтаксично коректними?
%enditem
\par
\vspace{2mm}
%beginitem
4. Що таке часова складність алгоритму?
%enditem
\par
\vspace{2mm}
%beginitem
5. Яку часову оцінку складності має алгоритм швидкого сортування?
%enditem
\par
\vspace{2mm}
%beginitem
6. Написати програму, що запитує у користувача значення дійсних параметрів $x$ та $y$, обчислює для них значення виразу та повiдомляє введенi данi та результат обчислення.
$$\frac{\sqrt x - 69}{(\arctg x - 0.2) (y - 37)}$$ 

%enditem
\par
%beginitem
7. Написати функцію, що для інтервалу $[a, b)$ знаходить знаходить найбільший степінь числа 3, на який ділиться хоча б одне ціле число з цього інтервалу.

По завданням 2 та 3 оцінюється правильність та якість коду, а також економність обчислень.

%enditem
}
{\smallskip\smallskip \begin {scriptsize} Затверджено на засіданні кафедри теоретичної кібернетики, протокол \No 3  від   08.11.2025 р.\par\smallskip\smallskip\smallskip\smallskip
{\parbox {6.5 cm} {Зав. кафедри \dotfill Ю.В.~Крак}}\hspace {0.7cm} {\hbox to 6.5 cm { Екзаменатор \dotfill Т.О.~Карнаух}} \end{scriptsize}}
%end
\vspace{0.9cm}
%begin
{{\begin {scriptsize}\par
{ \center  Київський національний університет імені Тараса Шевченка \par}
{\parbox {5 cm}{Освітня програма \hfill}{\parbox {7 cm}{Інформатика\hfill}}\parbox {6 cm}{\hfill Семестр 1}}\par
{\parbox {5 cm} {Навчальний предмет \hfill}\parbox {7 cm}{Програмування \hfill}\parbox {6cm}{\hfill}\par}\vspace{-15pt} \end{scriptsize}}{\center Екзаменаційний білет \No \textbf { 88 }\par}}
{%beginitem
1. Які з наведених типів не є цілими (integral), якщо такі є серед перелічених?

bool*, char, int, unsigned?
%enditem
\par
\vspace{2mm}
%beginitem
2. Скільки змінних, що є вказівниками, тут оголошено?
long *n1, *n2, n3, pn4;?
%enditem
\par
\vspace{2mm}
%beginitem
3. За наявних оголошень \code{void f(char); char c;} які з виразів \code{f(c*c), f(c+=2), f(\&c)} є синтаксично коректними?
%enditem
\par
\vspace{2mm}
%beginitem
4. Чим префіксний оператор ++ відрізняється від постфіксного?
%enditem
\par
\vspace{2mm}
%beginitem
5. Яку часову оцінку складності має алгоритм сортування вибором?
%enditem
\par
\vspace{2mm}
%beginitem
6. Написати програму, що запитує у користувача значення дійсних параметрів $x$ та $y$, обчислює для них значення виразу та повiдомляє введенi данi та результат обчислення.
$$\frac{\ln x + 59}{(\ctg x + 0.5) (y + 35)}$$ 

%enditem
\par
%beginitem
7. Написати функцію, що для інтервалу $[a, b)$ знаходить найближчу пару чисел з цього інтервалу, що кожне з них має парну кількість різних простих дільників. Якщо таких пар кілька, то цікавить пара з найбільшою сумою.

По завданням 2 та 3 оцінюється правильність та якість коду, а також економність обчислень.

%enditem
}
{\smallskip\smallskip \begin {scriptsize} Затверджено на засіданні кафедри теоретичної кібернетики, протокол \No 3  від   08.11.2025 р.\par\smallskip\smallskip\smallskip\smallskip
{\parbox {6.5 cm} {Зав. кафедри \dotfill Ю.В.~Крак}}\hspace {0.7cm} {\hbox to 6.5 cm { Екзаменатор \dotfill Т.О.~Карнаух}} \end{scriptsize}}
%end
\newpage
%begin
{{\begin {scriptsize}\par
{ \center  Київський національний університет імені Тараса Шевченка \par}
{\parbox {5 cm}{Освітня програма \hfill}{\parbox {7 cm}{Інформатика\hfill}}\parbox {6 cm}{\hfill Семестр 1}}\par
{\parbox {5 cm} {Навчальний предмет \hfill}\parbox {7 cm}{Програмування \hfill}\parbox {6cm}{\hfill}\par}\vspace{-15pt} \end{scriptsize}}{\center Екзаменаційний білет \No \textbf { 89 }\par}}
{%beginitem
1. Які з наведених типів не є цілими (integral), якщо такі є серед перелічених?

int*, signed char, unsigned long, string?
%enditem
\par
\vspace{2mm}
%beginitem
2. Скільки змінних, що є вказівниками, тут оголошено?
int n, m, *p1, p2;?
%enditem
\par
\vspace{2mm}
%beginitem
3. За наявних оголошень \code{void f(int*); int n;} які з виразів \code{f(n++), f(++n), f(\&n)} є синтаксично коректними?
%enditem
\par
\vspace{2mm}
%beginitem
4. Що таке просторова складність задачі?
%enditem
\par
\vspace{2mm}
%beginitem
5. Яку часову оцінку складності має алгоритм сортування вставками?
%enditem
\par
\vspace{2mm}
%beginitem
6. Написати програму, що запитує у користувача значення дійсних параметрів $x$ та $y$, обчислює для них значення виразу та повiдомляє введенi данi та результат обчислення.
$$\frac{\ x + 72}{(\sin x - 0.6) (y - 67)}$$ 

%enditem
\par
%beginitem
7. Написати функцію, що для послідовності цілих (задається масивом) знаходить довжину найдовшого відрізку, в якому всі сусідні числа відрізняються якнайменш на 2. Використовувати додаткові структури даних (у тому числі рядки) заборонено.

По завданням 2 та 3 оцінюється правильність та якість коду, а також економність обчислень.

%enditem
}
{\smallskip\smallskip \begin {scriptsize} Затверджено на засіданні кафедри теоретичної кібернетики, протокол \No 3  від   08.11.2025 р.\par\smallskip\smallskip\smallskip\smallskip
{\parbox {6.5 cm} {Зав. кафедри \dotfill Ю.В.~Крак}}\hspace {0.7cm} {\hbox to 6.5 cm { Екзаменатор \dotfill Т.О.~Карнаух}} \end{scriptsize}}
%end
\vspace{0.9cm}
%begin
{{\begin {scriptsize}\par
{ \center  Київський національний університет імені Тараса Шевченка \par}
{\parbox {5 cm}{Освітня програма \hfill}{\parbox {7 cm}{Інформатика\hfill}}\parbox {6 cm}{\hfill Семестр 1}}\par
{\parbox {5 cm} {Навчальний предмет \hfill}\parbox {7 cm}{Програмування \hfill}\parbox {6cm}{\hfill}\par}\vspace{-15pt} \end{scriptsize}}{\center Екзаменаційний білет \No \textbf { 90 }\par}}
{%beginitem
1. Які з наведених типів не є цілими (integral), якщо такі є серед перелічених?

bool, char, int, unsigned?
%enditem
\par
\vspace{2mm}
%beginitem
2. Скільки змінних, що є вказівниками, тут оголошено?
char *s1, pc, **s2;?
%enditem
\par
\vspace{2mm}
%beginitem
3. За наявних оголошень \code{void f(double); double x;} які з виразів \code{f(3.1), f(x+=2), f(\&x)} є синтаксично коректними?
%enditem
\par
\vspace{2mm}
%beginitem
4. Що таке тип даних?
%enditem
\par
\vspace{2mm}
%beginitem
5. Яку часову оцінку складності має алгоритм сортування злиттям?
%enditem
\par
\vspace{2mm}
%beginitem
6. Написати програму, що запитує у користувача значення дійсних параметрів $x$ та $y$, обчислює для них значення виразу та повiдомляє введенi данi та результат обчислення.
$$\frac{\log_{2} x - 99}{(\tg x + 0.9) (y + 69)}$$ 

%enditem
\par
%beginitem
7. Написати функцію, що для цілого числа знаходить суму його простих дільників, що входять у його розклад не більш ніж у 4му степеню. Наприклад, для числа $3^5{\cdot}7^9{\cdot}11{\cdot}23^{2}$ таких дільників два – $11$ та $23$, тому функція має повернути $34$.

По завданням 2 та 3 оцінюється правильність та якість коду, а також економність обчислень.

%enditem
}
{\smallskip\smallskip \begin {scriptsize} Затверджено на засіданні кафедри теоретичної кібернетики, протокол \No 3  від   08.11.2025 р.\par\smallskip\smallskip\smallskip\smallskip
{\parbox {6.5 cm} {Зав. кафедри \dotfill Ю.В.~Крак}}\hspace {0.7cm} {\hbox to 6.5 cm { Екзаменатор \dotfill Т.О.~Карнаух}} \end{scriptsize}}
%end
\newpage
%begin
{{\begin {scriptsize}\par
{ \center  Київський національний університет імені Тараса Шевченка \par}
{\parbox {5 cm}{Освітня програма \hfill}{\parbox {7 cm}{Інформатика\hfill}}\parbox {6 cm}{\hfill Семестр 1}}\par
{\parbox {5 cm} {Навчальний предмет \hfill}\parbox {7 cm}{Програмування \hfill}\parbox {6cm}{\hfill}\par}\vspace{-15pt} \end{scriptsize}}{\center Екзаменаційний білет \No \textbf { 91 }\par}}
{%beginitem
1. Які з наведених типів не є цілими (integral), якщо такі є серед перелічених?

bool*, char, int, unsigned?
%enditem
\par
\vspace{2mm}
%beginitem
2. Скільки змінних, що є вказівниками, тут оголошено?
double x, *px1, px2, ***px3;?
%enditem
\par
\vspace{2mm}
%beginitem
3. За наявних оголошень \code{void f(int*); int n;} які з виразів \code{f(n++), f(++n), f(\&n)} є синтаксично коректними?
%enditem
\par
\vspace{2mm}
%beginitem
4. Що таке часова складність алгоритму?
%enditem
\par
\vspace{2mm}
%beginitem
5. Яку часову оцінку складності має алгоритм сортування бульбашкою?
%enditem
\par
\vspace{2mm}
%beginitem
6. Написати програму, що запитує у користувача значення дійсних параметрів $x$ та $y$, обчислює для них значення виразу та повiдомляє введенi данi та результат обчислення.
$$\frac{\log_{10} x + 27}{(\cos x + 0.3) (y - 38)}$$ 

%enditem
\par
%beginitem
7. Написати функцію, що для інтервалу $[a, b)$ знаходить найближчу пару чисел з цього інтервалу, що кожне з них має непарну кількість різних простих дільників. Якщо таких пар кілька, то цікавить пара з найменшим добутком.

По завданням 2 та 3 оцінюється правильність та якість коду, а також економність обчислень.

%enditem
}
{\smallskip\smallskip \begin {scriptsize} Затверджено на засіданні кафедри теоретичної кібернетики, протокол \No 3  від   08.11.2025 р.\par\smallskip\smallskip\smallskip\smallskip
{\parbox {6.5 cm} {Зав. кафедри \dotfill Ю.В.~Крак}}\hspace {0.7cm} {\hbox to 6.5 cm { Екзаменатор \dotfill Т.О.~Карнаух}} \end{scriptsize}}
%end
\vspace{0.9cm}
%begin
{{\begin {scriptsize}\par
{ \center  Київський національний університет імені Тараса Шевченка \par}
{\parbox {5 cm}{Освітня програма \hfill}{\parbox {7 cm}{Інформатика\hfill}}\parbox {6 cm}{\hfill Семестр 1}}\par
{\parbox {5 cm} {Навчальний предмет \hfill}\parbox {7 cm}{Програмування \hfill}\parbox {6cm}{\hfill}\par}\vspace{-15pt} \end{scriptsize}}{\center Екзаменаційний білет \No \textbf { 92 }\par}}
{%beginitem
1. Які з наведених типів не є цілими (integral), якщо такі є серед перелічених?

double, int, long, unsigned char?
%enditem
\par
\vspace{2mm}
%beginitem
2. Скільки змінних, що є вказівниками, тут оголошено?
double x, *px1, px2, ***px3;?
%enditem
\par
\vspace{2mm}
%beginitem
3. За наявних оголошень \code{void f(double*); double x;} які з виразів \code{f(x*x), f(x+=2), f(\&x)} є синтаксично коректними?
%enditem
\par
\vspace{2mm}
%beginitem
4. Чим префіксний оператор ++ відрізняється від постфіксного?
%enditem
\par
\vspace{2mm}
%beginitem
5. Яку часову оцінку складності має алгоритм сортування купою (він же деревом, пірамідою)?
%enditem
\par
\vspace{2mm}
%beginitem
6. Написати програму, що запитує у користувача значення дійсних параметрів $x$ та $y$, обчислює для них значення виразу та повiдомляє введенi данi та результат обчислення.
$$\frac{\ln x - 61}{(\sin x - 0.2) (y + 75)}$$ 

%enditem
\par
%beginitem
7. Написати функцію, що повертає найбільше число з послідовності цілих (задається масивом), що має рівно три різних простих дільники. За відсутності має повертати None. Використовувати додаткові структури даних (у тому числі рядки) заборонено.

По завданням 2 та 3 оцінюється правильність та якість коду, а також економність обчислень.

%enditem
}
{\smallskip\smallskip \begin {scriptsize} Затверджено на засіданні кафедри теоретичної кібернетики, протокол \No 3  від   08.11.2025 р.\par\smallskip\smallskip\smallskip\smallskip
{\parbox {6.5 cm} {Зав. кафедри \dotfill Ю.В.~Крак}}\hspace {0.7cm} {\hbox to 6.5 cm { Екзаменатор \dotfill Т.О.~Карнаух}} \end{scriptsize}}
%end
\newpage
%begin
{{\begin {scriptsize}\par
{ \center  Київський національний університет імені Тараса Шевченка \par}
{\parbox {5 cm}{Освітня програма \hfill}{\parbox {7 cm}{Інформатика\hfill}}\parbox {6 cm}{\hfill Семестр 1}}\par
{\parbox {5 cm} {Навчальний предмет \hfill}\parbox {7 cm}{Програмування \hfill}\parbox {6cm}{\hfill}\par}\vspace{-15pt} \end{scriptsize}}{\center Екзаменаційний білет \No \textbf { 93 }\par}}
{%beginitem
1. Які з наведених типів не є цілими (integral), якщо такі є серед перелічених?

double, int, long*, unsigned char?
%enditem
\par
\vspace{2mm}
%beginitem
2. Скільки змінних, що є вказівниками, тут оголошено?
double *px, *py, **ppx;?
%enditem
\par
\vspace{2mm}
%beginitem
3. За наявних оголошень \code{void f(double); double x;} які з виразів \code{f(3.1), f(x+=2), f(\&x)} є синтаксично коректними?
%enditem
\par
\vspace{2mm}
%beginitem
4. У чому відмінність літералів від констант?
%enditem
\par
\vspace{2mm}
%beginitem
5. Яку часову оцінку складності має алгоритм сортування вставками?
%enditem
\par
\vspace{2mm}
%beginitem
6. Написати програму, що запитує у користувача значення дійсних параметрів $x$ та $y$, обчислює для них значення виразу та повiдомляє введенi данi та результат обчислення.
$$\frac{\log_{2} x - 84}{(\arcsin x + 0.1) (y - 65)}$$ 

%enditem
\par
%beginitem
7. Написати функцію, що для цілого числа знаходить простий дільник, який входить у його розклад ц найбільшому степені. Якщо таких кілька, то повернути найбільший. Наприклад, для числа $2^{9}{\cdot}3^5{\cdot}7^9{\cdot}11$ функція має повернути $7$.

По завданням 2 та 3 оцінюється правильність та якість коду, а також економність обчислень.

%enditem
}
{\smallskip\smallskip \begin {scriptsize} Затверджено на засіданні кафедри теоретичної кібернетики, протокол \No 3  від   08.11.2025 р.\par\smallskip\smallskip\smallskip\smallskip
{\parbox {6.5 cm} {Зав. кафедри \dotfill Ю.В.~Крак}}\hspace {0.7cm} {\hbox to 6.5 cm { Екзаменатор \dotfill Т.О.~Карнаух}} \end{scriptsize}}
%end
\vspace{0.9cm}
%begin
{{\begin {scriptsize}\par
{ \center  Київський національний університет імені Тараса Шевченка \par}
{\parbox {5 cm}{Освітня програма \hfill}{\parbox {7 cm}{Інформатика\hfill}}\parbox {6 cm}{\hfill Семестр 1}}\par
{\parbox {5 cm} {Навчальний предмет \hfill}\parbox {7 cm}{Програмування \hfill}\parbox {6cm}{\hfill}\par}\vspace{-15pt} \end{scriptsize}}{\center Екзаменаційний білет \No \textbf { 94 }\par}}
{%beginitem
1. Які з наведених типів не є цілими (integral), якщо такі є серед перелічених?

bool, char, int, unsigned?
%enditem
\par
\vspace{2mm}
%beginitem
2. Скільки змінних, що є вказівниками, тут оголошено?
long *n1, *n2, n3, pn4;?
%enditem
\par
\vspace{2mm}
%beginitem
3. За наявних оголошень \code{void f(int\&); int n;} які з виразів \code{f(13), f(++n), f(\&n)} є синтаксично коректними?
%enditem
\par
\vspace{2mm}
%beginitem
4. Що таке часова складність задачі?
%enditem
\par
\vspace{2mm}
%beginitem
5. Яку часову оцінку складності має алгоритм сортування вибором?
%enditem
\par
\vspace{2mm}
%beginitem
6. Написати програму, що запитує у користувача значення дійсних параметрів $x$ та $y$, обчислює для них значення виразу та повiдомляє введенi данi та результат обчислення.
$$\frac{\sqrt x + 89}{(\arctg x - 0.8) (y + 18)}$$ 

%enditem
\par
%beginitem
7. Написати функцію, що повертає найбільше число з послідовності цілих (задається масивом), у сімковому запису якого нема цифри 2. За відсутності має повертати None. Використовувати додаткові структури даних (у тому числі рядки) заборонено.

По завданням 2 та 3 оцінюється правильність та якість коду, а також економність обчислень.

%enditem
}
{\smallskip\smallskip \begin {scriptsize} Затверджено на засіданні кафедри теоретичної кібернетики, протокол \No 3  від   08.11.2025 р.\par\smallskip\smallskip\smallskip\smallskip
{\parbox {6.5 cm} {Зав. кафедри \dotfill Ю.В.~Крак}}\hspace {0.7cm} {\hbox to 6.5 cm { Екзаменатор \dotfill Т.О.~Карнаух}} \end{scriptsize}}
%end
\newpage
%begin
{{\begin {scriptsize}\par
{ \center  Київський національний університет імені Тараса Шевченка \par}
{\parbox {5 cm}{Освітня програма \hfill}{\parbox {7 cm}{Інформатика\hfill}}\parbox {6 cm}{\hfill Семестр 1}}\par
{\parbox {5 cm} {Навчальний предмет \hfill}\parbox {7 cm}{Програмування \hfill}\parbox {6cm}{\hfill}\par}\vspace{-15pt} \end{scriptsize}}{\center Екзаменаційний білет \No \textbf { 95 }\par}}
{%beginitem
1. Які з наведених типів не є цілими (integral), якщо такі є серед перелічених?

bool, int, long, long long, unsigned char?
%enditem
\par
\vspace{2mm}
%beginitem
2. Скільки змінних, що є вказівниками, тут оголошено?
int n, m, *p1, p2;?
%enditem
\par
\vspace{2mm}
%beginitem
3. За наявних оголошень \code{void f(int); int n;} які з виразів \code{f(3), f(++n), f(\&n)} є синтаксично коректними?
%enditem
\par
\vspace{2mm}
%beginitem
4. Що означає термін ``часова складність у середньому''?
%enditem
\par
\vspace{2mm}
%beginitem
5. Яку часову оцінку складності має алгоритм швидкого сортування?
%enditem
\par
\vspace{2mm}
%beginitem
6. Написати програму, що запитує у користувача значення дійсних параметрів $x$ та $y$, обчислює для них значення виразу та повiдомляє введенi данi та результат обчислення.
$$\frac{\pi x - 11}{(\tg x + 0.7) (y + 78)}$$ 

%enditem
\par
%beginitem
7. Написати функцію, що повертає найбільше число з послідовності цілих (задається масивом), що має не більше трьох різних простих дільників. За відсутності має повертати None. Використовувати додаткові структури даних (у тому числі рядки) заборонено.

По завданням 2 та 3 оцінюється правильність та якість коду, а також економність обчислень.

%enditem
}
{\smallskip\smallskip \begin {scriptsize} Затверджено на засіданні кафедри теоретичної кібернетики, протокол \No 3  від   08.11.2025 р.\par\smallskip\smallskip\smallskip\smallskip
{\parbox {6.5 cm} {Зав. кафедри \dotfill Ю.В.~Крак}}\hspace {0.7cm} {\hbox to 6.5 cm { Екзаменатор \dotfill Т.О.~Карнаух}} \end{scriptsize}}
%end
\vspace{0.9cm}
%begin
{{\begin {scriptsize}\par
{ \center  Київський національний університет імені Тараса Шевченка \par}
{\parbox {5 cm}{Освітня програма \hfill}{\parbox {7 cm}{Інформатика\hfill}}\parbox {6 cm}{\hfill Семестр 1}}\par
{\parbox {5 cm} {Навчальний предмет \hfill}\parbox {7 cm}{Програмування \hfill}\parbox {6cm}{\hfill}\par}\vspace{-15pt} \end{scriptsize}}{\center Екзаменаційний білет \No \textbf { 96 }\par}}
{%beginitem
1. Які з наведених типів не є цілими (integral), якщо такі є серед перелічених?

bool, int, long, long long, unsigned char?
%enditem
\par
\vspace{2mm}
%beginitem
2. Скільки змінних, що є вказівниками, тут оголошено?
char *s1, pc, **s2;?
%enditem
\par
\vspace{2mm}
%beginitem
3. За наявних оголошень \code{void f(char\&); char c;} які з виразів \code{f(c*c), f(c++), f(\&c)} є синтаксично коректними?
%enditem
\par
\vspace{2mm}
%beginitem
4. Що таке масив?
%enditem
\par
\vspace{2mm}
%beginitem
5. Яку часову оцінку складності має алгоритм сортування бульбашкою?
%enditem
\par
\vspace{2mm}
%beginitem
6. Написати програму, що запитує у користувача значення дійсних параметрів $x$ та $y$, обчислює для них значення виразу та повiдомляє введенi данi та результат обчислення.
$$\frac{\ x + 28}{(\ctg x - 0.9) (y - 19)}$$ 

%enditem
\par
%beginitem
7. Написати функцію, що для послідовності цілих (задається масивом) знаходить довжину найдовшого відрізку, в якому всі сусідні числа різні. Використовувати додаткові структури даних (у тому числі рядки) заборонено.

По завданням 2 та 3 оцінюється правильність та якість коду, а також економність обчислень.

%enditem
}
{\smallskip\smallskip \begin {scriptsize} Затверджено на засіданні кафедри теоретичної кібернетики, протокол \No 3  від   08.11.2025 р.\par\smallskip\smallskip\smallskip\smallskip
{\parbox {6.5 cm} {Зав. кафедри \dotfill Ю.В.~Крак}}\hspace {0.7cm} {\hbox to 6.5 cm { Екзаменатор \dotfill Т.О.~Карнаух}} \end{scriptsize}}
%end
\newpage
%begin
{{\begin {scriptsize}\par
{ \center  Київський національний університет імені Тараса Шевченка \par}
{\parbox {5 cm}{Освітня програма \hfill}{\parbox {7 cm}{Інформатика\hfill}}\parbox {6 cm}{\hfill Семестр 1}}\par
{\parbox {5 cm} {Навчальний предмет \hfill}\parbox {7 cm}{Програмування \hfill}\parbox {6cm}{\hfill}\par}\vspace{-15pt} \end{scriptsize}}{\center Екзаменаційний білет \No \textbf { 97 }\par}}
{%beginitem
1. Які з наведених типів не є цілими (integral), якщо такі є серед перелічених?

char, double, int, float, unsigned*?
%enditem
\par
\vspace{2mm}
%beginitem
2. Скільки змінних, що є вказівниками, тут оголошено?
long pointer1, pointer2, pointer3, pointer4;?
%enditem
\par
\vspace{2mm}
%beginitem
3. За наявних оголошень \code{void f(double\&); double x;} які з виразів \code{f(x*x), f(x+=2), f(\&x)} є синтаксично коректними?
%enditem
\par
\vspace{2mm}
%beginitem
4. Що таке С-рядок?
%enditem
\par
\vspace{2mm}
%beginitem
5. Яку часову оцінку складності має алгоритм швидкого сортування?
%enditem
\par
\vspace{2mm}
%beginitem
6. Написати програму, що запитує у користувача значення дійсних параметрів $x$ та $y$, обчислює для них значення виразу та повiдомляє введенi данi та результат обчислення.
$$\frac{\ x + 43}{(\cos x + 0.5) (y - 32)}$$ 

%enditem
\par
%beginitem
7. Написати функцію, що для інтервалу $[a, b)$ знаходить кількість чисел, квадрати яких діляться на суму квадратів їх цифр. (Розглядаємо десятковий запис.)

По завданням 2 та 3 оцінюється правильність та якість коду, а також економність обчислень.

%enditem
}
{\smallskip\smallskip \begin {scriptsize} Затверджено на засіданні кафедри теоретичної кібернетики, протокол \No 3  від   08.11.2025 р.\par\smallskip\smallskip\smallskip\smallskip
{\parbox {6.5 cm} {Зав. кафедри \dotfill Ю.В.~Крак}}\hspace {0.7cm} {\hbox to 6.5 cm { Екзаменатор \dotfill Т.О.~Карнаух}} \end{scriptsize}}
%end
\vspace{0.9cm}
%begin
{{\begin {scriptsize}\par
{ \center  Київський національний університет імені Тараса Шевченка \par}
{\parbox {5 cm}{Освітня програма \hfill}{\parbox {7 cm}{Інформатика\hfill}}\parbox {6 cm}{\hfill Семестр 1}}\par
{\parbox {5 cm} {Навчальний предмет \hfill}\parbox {7 cm}{Програмування \hfill}\parbox {6cm}{\hfill}\par}\vspace{-15pt} \end{scriptsize}}{\center Екзаменаційний білет \No \textbf { 98 }\par}}
{%beginitem
1. Які з наведених типів не є цілими (integral), якщо такі є серед перелічених?

int, signed char, unsigned long, string?
%enditem
\par
\vspace{2mm}
%beginitem
2. Скільки змінних, що є вказівниками, тут оголошено?
int *pn1, pn2, pn3, pn4;?
%enditem
\par
\vspace{2mm}
%beginitem
3. За наявних оголошень \code{void f(char*); char c;} які з виразів \code{f(c*c), f(c+=2), f(\&c)} є синтаксично коректними?
%enditem
\par
\vspace{2mm}
%beginitem
4. Що таке просторова складність алгоритму?
%enditem
\par
\vspace{2mm}
%beginitem
5. Яку часову оцінку складності має алгоритм сортування злиттям?
%enditem
\par
\vspace{2mm}
%beginitem
6. Написати програму, що запитує у користувача значення дійсних параметрів $x$ та $y$, обчислює для них значення виразу та повiдомляє введенi данi та результат обчислення.
$$\frac{\log_{2} x - 76}{(\arccos x - 0.4) (y + 37)}$$ 

%enditem
\par
%beginitem
7. Написати функцію, що в інтервалі $[a, b)$ знаходить ціле число, що ділиться на найбільший степінь числа 3. Якщо таких чисел декілька, повернути найменше.

По завданням 2 та 3 оцінюється правильність та якість коду, а також економність обчислень.

%enditem
}
{\smallskip\smallskip \begin {scriptsize} Затверджено на засіданні кафедри теоретичної кібернетики, протокол \No 3  від   08.11.2025 р.\par\smallskip\smallskip\smallskip\smallskip
{\parbox {6.5 cm} {Зав. кафедри \dotfill Ю.В.~Крак}}\hspace {0.7cm} {\hbox to 6.5 cm { Екзаменатор \dotfill Т.О.~Карнаух}} \end{scriptsize}}
%end
\newpage
%begin
{{\begin {scriptsize}\par
{ \center  Київський національний університет імені Тараса Шевченка \par}
{\parbox {5 cm}{Освітня програма \hfill}{\parbox {7 cm}{Інформатика\hfill}}\parbox {6 cm}{\hfill Семестр 1}}\par
{\parbox {5 cm} {Навчальний предмет \hfill}\parbox {7 cm}{Програмування \hfill}\parbox {6cm}{\hfill}\par}\vspace{-15pt} \end{scriptsize}}{\center Екзаменаційний білет \No \textbf { 99 }\par}}
{%beginitem
1. Які з наведених типів не є цілими (integral), якщо такі є серед перелічених?

int*, signed char, unsigned long, string?
%enditem
\par
\vspace{2mm}
%beginitem
2. Скільки змінних, що є вказівниками, тут оголошено?
long pointer1, pointer2, pointer3, pointer4;?
%enditem
\par
\vspace{2mm}
%beginitem
3. За наявних оголошень \code{void f(char); char c;} які з виразів \code{f(c*c), f(c+=2), f(\&c)} є синтаксично коректними?
%enditem
\par
\vspace{2mm}
%beginitem
4. Які параметри може приймати функція main? Що вони позначають та яких вони типів?
%enditem
\par
\vspace{2mm}
%beginitem
5. Яку часову оцінку складності має алгоритм сортування бульбашкою?
%enditem
\par
\vspace{2mm}
%beginitem
6. Написати програму, що запитує у користувача значення дійсних параметрів $x$ та $y$, обчислює для них значення виразу та повiдомляє введенi данi та результат обчислення.
$$\frac{\ln x + 82}{(\ctg x + 0.6) (y + 25)}$$ 

%enditem
\par
%beginitem
7. Написати функцію, що повертає число з послідовності цілих (задається масивом), що має парну кількість різних простих дільників. Якщо таких чисел кілька, то повернути найбільше. Використовувати додаткові структури даних (у тому числі рядки) заборонено.

По завданням 2 та 3 оцінюється правильність та якість коду, а також економність обчислень.

%enditem
}
{\smallskip\smallskip \begin {scriptsize} Затверджено на засіданні кафедри теоретичної кібернетики, протокол \No 3  від   08.11.2025 р.\par\smallskip\smallskip\smallskip\smallskip
{\parbox {6.5 cm} {Зав. кафедри \dotfill Ю.В.~Крак}}\hspace {0.7cm} {\hbox to 6.5 cm { Екзаменатор \dotfill Т.О.~Карнаух}} \end{scriptsize}}
%end
\vspace{0.9cm}
%begin
{{\begin {scriptsize}\par
{ \center  Київський національний університет імені Тараса Шевченка \par}
{\parbox {5 cm}{Освітня програма \hfill}{\parbox {7 cm}{Інформатика\hfill}}\parbox {6 cm}{\hfill Семестр 1}}\par
{\parbox {5 cm} {Навчальний предмет \hfill}\parbox {7 cm}{Програмування \hfill}\parbox {6cm}{\hfill}\par}\vspace{-15pt} \end{scriptsize}}{\center Екзаменаційний білет \No \textbf { 100 }\par}}
{%beginitem
1. Які з наведених типів не є цілими (integral), якщо такі є серед перелічених?

char, double, int, float, unsigned?
%enditem
\par
\vspace{2mm}
%beginitem
2. Скільки змінних, що є вказівниками, тут оголошено?
int n, m, *p1, p2;?
%enditem
\par
\vspace{2mm}
%beginitem
3. За наявних оголошень \code{void f(int\&); int n;} які з виразів \code{f(13), f(++n), f(\&n)} є синтаксично коректними?
%enditem
\par
\vspace{2mm}
%beginitem
4. Що означає термін ``часова складність у середньому''?
%enditem
\par
\vspace{2mm}
%beginitem
5. Яку часову оцінку складності має алгоритм сортування вставками?
%enditem
\par
\vspace{2mm}
%beginitem
6. Написати програму, що запитує у користувача значення дійсних параметрів $x$ та $y$, обчислює для них значення виразу та повiдомляє введенi данi та результат обчислення.
$$\frac{\pi x - 71}{(\tg x - 0.9) (y - 62)}$$ 

%enditem
\par
%beginitem
7. Написати функцію, що повертає число з послідовності цілих (задається масивом), що має найменшу кількість різних простих дільників. Якщо таких чисел кілька, то повернути найбільше. Використовувати додаткові структури даних (у тому числі рядки) заборонено.

По завданням 2 та 3 оцінюється правильність та якість коду, а також економність обчислень.

%enditem
}
{\smallskip\smallskip \begin {scriptsize} Затверджено на засіданні кафедри теоретичної кібернетики, протокол \No 3  від   08.11.2025 р.\par\smallskip\smallskip\smallskip\smallskip
{\parbox {6.5 cm} {Зав. кафедри \dotfill Ю.В.~Крак}}\hspace {0.7cm} {\hbox to 6.5 cm { Екзаменатор \dotfill Т.О.~Карнаух}} \end{scriptsize}}
%end
\newpage
%begin
{{\begin {scriptsize}\par
{ \center  Київський національний університет імені Тараса Шевченка \par}
{\parbox {5 cm}{Освітня програма \hfill}{\parbox {7 cm}{Інформатика\hfill}}\parbox {6 cm}{\hfill Семестр 1}}\par
{\parbox {5 cm} {Навчальний предмет \hfill}\parbox {7 cm}{Програмування \hfill}\parbox {6cm}{\hfill}\par}\vspace{-15pt} \end{scriptsize}}{\center Екзаменаційний білет \No \textbf { 101 }\par}}
{%beginitem
1. Які з наведених типів не є цілими (integral), якщо такі є серед перелічених?

bool, char, int, unsigned?
%enditem
\par
\vspace{2mm}
%beginitem
2. Скільки змінних, що є вказівниками, тут оголошено?
char *s1, pc, **s2;?
%enditem
\par
\vspace{2mm}
%beginitem
3. За наявних оголошень \code{void f(char); char c;} які з виразів \code{f(c*c), f(c+=2), f(\&c)} є синтаксично коректними?
%enditem
\par
\vspace{2mm}
%beginitem
4. Які параметри може приймати функція main? Що вони позначають та яких вони типів?
%enditem
\par
\vspace{2mm}
%beginitem
5. Яку часову оцінку складності має алгоритм сортування вибором?
%enditem
\par
\vspace{2mm}
%beginitem
6. Написати програму, що запитує у користувача значення дійсних параметрів $x$ та $y$, обчислює для них значення виразу та повiдомляє введенi данi та результат обчислення.
$$\frac{\sqrt x - 54}{(\arctg x - 0.2) (y - 98)}$$ 

%enditem
\par
%beginitem
7. Написати функцію, що для послідовності цілих (задається масивом) знаходить довжину найдовшого відрізку, що складається з чисел, у сімковому запису якого нема цифри 2. За відсутності має повертати $0$. Використовувати додаткові структури даних (у тому числі рядки) заборонено.

По завданням 2 та 3 оцінюється правильність та якість коду, а також економність обчислень.

%enditem
}
{\smallskip\smallskip \begin {scriptsize} Затверджено на засіданні кафедри теоретичної кібернетики, протокол \No 3  від   08.11.2025 р.\par\smallskip\smallskip\smallskip\smallskip
{\parbox {6.5 cm} {Зав. кафедри \dotfill Ю.В.~Крак}}\hspace {0.7cm} {\hbox to 6.5 cm { Екзаменатор \dotfill Т.О.~Карнаух}} \end{scriptsize}}
%end
\vspace{0.9cm}
%begin
{{\begin {scriptsize}\par
{ \center  Київський національний університет імені Тараса Шевченка \par}
{\parbox {5 cm}{Освітня програма \hfill}{\parbox {7 cm}{Інформатика\hfill}}\parbox {6 cm}{\hfill Семестр 1}}\par
{\parbox {5 cm} {Навчальний предмет \hfill}\parbox {7 cm}{Програмування \hfill}\parbox {6cm}{\hfill}\par}\vspace{-15pt} \end{scriptsize}}{\center Екзаменаційний білет \No \textbf { 102 }\par}}
{%beginitem
1. Які з наведених типів не є цілими (integral), якщо такі є серед перелічених?

bool, int, long, long long, unsigned char?
%enditem
\par
\vspace{2mm}
%beginitem
2. Скільки змінних, що є вказівниками, тут оголошено?
double *px, *py, **ppx;?
%enditem
\par
\vspace{2mm}
%beginitem
3. За наявних оголошень \code{void f(int); int n;} які з виразів \code{f(3), f(++n), f(\&n)} є синтаксично коректними?
%enditem
\par
\vspace{2mm}
%beginitem
4. Чим префіксний оператор ++ відрізняється від постфіксного?
%enditem
\par
\vspace{2mm}
%beginitem
5. Яку часову оцінку складності має алгоритм швидкого сортування?
%enditem
\par
\vspace{2mm}
%beginitem
6. Написати програму, що запитує у користувача значення дійсних параметрів $x$ та $y$, обчислює для них значення виразу та повiдомляє введенi данi та результат обчислення.
$$\frac{\log_{10} x + 80}{(\arcsin x + 0.3) (y + 23)}$$ 

%enditem
\par
%beginitem
7. Написати функцію, що для інтервалу $[a, b)$ знаходить найближчу пару чисел з цього інтервалу, що мають не менше трьох простих дільників та мають цифру 3 у своєму десятковому запису.

По завданням 2 та 3 оцінюється правильність та якість коду, а також економність обчислень.

%enditem
}
{\smallskip\smallskip \begin {scriptsize} Затверджено на засіданні кафедри теоретичної кібернетики, протокол \No 3  від   08.11.2025 р.\par\smallskip\smallskip\smallskip\smallskip
{\parbox {6.5 cm} {Зав. кафедри \dotfill Ю.В.~Крак}}\hspace {0.7cm} {\hbox to 6.5 cm { Екзаменатор \dotfill Т.О.~Карнаух}} \end{scriptsize}}
%end
\newpage
%begin
{{\begin {scriptsize}\par
{ \center  Київський національний університет імені Тараса Шевченка \par}
{\parbox {5 cm}{Освітня програма \hfill}{\parbox {7 cm}{Інформатика\hfill}}\parbox {6 cm}{\hfill Семестр 1}}\par
{\parbox {5 cm} {Навчальний предмет \hfill}\parbox {7 cm}{Програмування \hfill}\parbox {6cm}{\hfill}\par}\vspace{-15pt} \end{scriptsize}}{\center Екзаменаційний білет \No \textbf { 103 }\par}}
{%beginitem
1. Які з наведених типів не є цілими (integral), якщо такі є серед перелічених?

bool*, char, int, unsigned?
%enditem
\par
\vspace{2mm}
%beginitem
2. Скільки змінних, що є вказівниками, тут оголошено?
long *n1, *n2, n3, pn4;?
%enditem
\par
\vspace{2mm}
%beginitem
3. За наявних оголошень \code{void f(int*); int n;} які з виразів \code{f(n++), f(++n), f(\&n)} є синтаксично коректними?
%enditem
\par
\vspace{2mm}
%beginitem
4. Що таке часова складність алгоритму?
%enditem
\par
\vspace{2mm}
%beginitem
5. Яку часову оцінку складності має алгоритм сортування купою (він же деревом, пірамідою)?
%enditem
\par
\vspace{2mm}
%beginitem
6. Написати програму, що запитує у користувача значення дійсних параметрів $x$ та $y$, обчислює для них значення виразу та повiдомляє введенi данi та результат обчислення.
$$\frac{\ x - 44}{(\arccos x - 0.1) (y - 21)}$$ 

%enditem
\par
%beginitem
7. Написати функцію, що для послідовності цілих (задається масивом) знаходить найдовший відрізок, в якому всі сусідні числа різні. Якщо таких відрізків кілька, то повернути перший (як індекс початку та довжину). Використовувати додаткові структури даних (у тому числі рядки) заборонено.

По завданням 2 та 3 оцінюється правильність та якість коду, а також економність обчислень.

%enditem
}
{\smallskip\smallskip \begin {scriptsize} Затверджено на засіданні кафедри теоретичної кібернетики, протокол \No 3  від   08.11.2025 р.\par\smallskip\smallskip\smallskip\smallskip
{\parbox {6.5 cm} {Зав. кафедри \dotfill Ю.В.~Крак}}\hspace {0.7cm} {\hbox to 6.5 cm { Екзаменатор \dotfill Т.О.~Карнаух}} \end{scriptsize}}
%end
\vspace{0.9cm}
%begin
{{\begin {scriptsize}\par
{ \center  Київський національний університет імені Тараса Шевченка \par}
{\parbox {5 cm}{Освітня програма \hfill}{\parbox {7 cm}{Інформатика\hfill}}\parbox {6 cm}{\hfill Семестр 1}}\par
{\parbox {5 cm} {Навчальний предмет \hfill}\parbox {7 cm}{Програмування \hfill}\parbox {6cm}{\hfill}\par}\vspace{-15pt} \end{scriptsize}}{\center Екзаменаційний білет \No \textbf { 104 }\par}}
{%beginitem
1. Які з наведених типів не є цілими (integral), якщо такі є серед перелічених?

double, int, long*, unsigned char?
%enditem
\par
\vspace{2mm}
%beginitem
2. Скільки змінних, що є вказівниками, тут оголошено?
int *pn1, pn2, pn3, pn4;?
%enditem
\par
\vspace{2mm}
%beginitem
3. За наявних оголошень \code{void f(double); double x;} які з виразів \code{f(3.1), f(x+=2), f(\&x)} є синтаксично коректними?
%enditem
\par
\vspace{2mm}
%beginitem
4. Що таке масив?
%enditem
\par
\vspace{2mm}
%beginitem
5. Яку часову оцінку складності має алгоритм сортування злиттям?
%enditem
\par
\vspace{2mm}
%beginitem
6. Написати програму, що запитує у користувача значення дійсних параметрів $x$ та $y$, обчислює для них значення виразу та повiдомляє введенi данi та результат обчислення.
$$\frac{\log_{10} x + 22}{(\cos x + 0.6) (y + 87)}$$ 

%enditem
\par
%beginitem
7. Написати функцію, що для інтервалу $[a, b)$ знаходить найближчу пару чисел з цього інтервалу, що кожне з них має непарну кількість різних простих дільників. Якщо таких пар кілька, то цікавить пара з найменшим добутком.

По завданням 2 та 3 оцінюється правильність та якість коду, а також економність обчислень.

%enditem
}
{\smallskip\smallskip \begin {scriptsize} Затверджено на засіданні кафедри теоретичної кібернетики, протокол \No 3  від   08.11.2025 р.\par\smallskip\smallskip\smallskip\smallskip
{\parbox {6.5 cm} {Зав. кафедри \dotfill Ю.В.~Крак}}\hspace {0.7cm} {\hbox to 6.5 cm { Екзаменатор \dotfill Т.О.~Карнаух}} \end{scriptsize}}
%end
\newpage
%begin
{{\begin {scriptsize}\par
{ \center  Київський національний університет імені Тараса Шевченка \par}
{\parbox {5 cm}{Освітня програма \hfill}{\parbox {7 cm}{Інформатика\hfill}}\parbox {6 cm}{\hfill Семестр 1}}\par
{\parbox {5 cm} {Навчальний предмет \hfill}\parbox {7 cm}{Програмування \hfill}\parbox {6cm}{\hfill}\par}\vspace{-15pt} \end{scriptsize}}{\center Екзаменаційний білет \No \textbf { 105 }\par}}
{%beginitem
1. Які з наведених типів не є цілими (integral), якщо такі є серед перелічених?

int, signed char, unsigned long, string?
%enditem
\par
\vspace{2mm}
%beginitem
2. Скільки змінних, що є вказівниками, тут оголошено?
double x, *px1, px2, ***px3;?
%enditem
\par
\vspace{2mm}
%beginitem
3. За наявних оголошень \code{void f(double*); double x;} які з виразів \code{f(x*x), f(x+=2), f(\&x)} є синтаксично коректними?
%enditem
\par
\vspace{2mm}
%beginitem
4. Що таке тип даних?
%enditem
\par
\vspace{2mm}
%beginitem
5. Яку часову оцінку складності має алгоритм швидкого сортування?
%enditem
\par
\vspace{2mm}
%beginitem
6. Написати програму, що запитує у користувача значення дійсних параметрів $x$ та $y$, обчислює для них значення виразу та повiдомляє введенi данi та результат обчислення.
$$\frac{\ln x - 60}{(\sin x - 0.8) (y - 56)}$$ 

%enditem
\par
%beginitem
7. Написати функцію, що повертає найбільше число з послідовності цілих (задається масивом), що має не більше трьох різних простих дільників. За відсутності має повертати None. Використовувати додаткові структури даних (у тому числі рядки) заборонено.

По завданням 2 та 3 оцінюється правильність та якість коду, а також економність обчислень.

%enditem
}
{\smallskip\smallskip \begin {scriptsize} Затверджено на засіданні кафедри теоретичної кібернетики, протокол \No 3  від   08.11.2025 р.\par\smallskip\smallskip\smallskip\smallskip
{\parbox {6.5 cm} {Зав. кафедри \dotfill Ю.В.~Крак}}\hspace {0.7cm} {\hbox to 6.5 cm { Екзаменатор \dotfill Т.О.~Карнаух}} \end{scriptsize}}
%end
\vspace{0.9cm}
%begin
{{\begin {scriptsize}\par
{ \center  Київський національний університет імені Тараса Шевченка \par}
{\parbox {5 cm}{Освітня програма \hfill}{\parbox {7 cm}{Інформатика\hfill}}\parbox {6 cm}{\hfill Семестр 1}}\par
{\parbox {5 cm} {Навчальний предмет \hfill}\parbox {7 cm}{Програмування \hfill}\parbox {6cm}{\hfill}\par}\vspace{-15pt} \end{scriptsize}}{\center Екзаменаційний білет \No \textbf { 106 }\par}}
{%beginitem
1. Які з наведених типів не є цілими (integral), якщо такі є серед перелічених?

int*, signed char, unsigned long, string?
%enditem
\par
\vspace{2mm}
%beginitem
2. Скільки змінних, що є вказівниками, тут оголошено?
double x, *px1, px2, ***px3;?
%enditem
\par
\vspace{2mm}
%beginitem
3. За наявних оголошень \code{void f(char\&); char c;} які з виразів \code{f(c*c), f(c++), f(\&c)} є синтаксично коректними?
%enditem
\par
\vspace{2mm}
%beginitem
4. Що таке просторова складність задачі?
%enditem
\par
\vspace{2mm}
%beginitem
5. Яку часову оцінку складності має алгоритм сортування вибором?
%enditem
\par
\vspace{2mm}
%beginitem
6. Написати програму, що запитує у користувача значення дійсних параметрів $x$ та $y$, обчислює для них значення виразу та повiдомляє введенi данi та результат обчислення.
$$\frac{\log_{2} x + 39}{(\arcsin x + 0.7) (y + 90)}$$ 

%enditem
\par
%beginitem
7. Написати функцію, що для послідовності цілих (задається масивом) знаходить довжину найдовшого відрізку, в якому всі сусідні числа відрізняються якнайменш на 2. Використовувати додаткові структури даних (у тому числі рядки) заборонено.

По завданням 2 та 3 оцінюється правильність та якість коду, а також економність обчислень.

%enditem
}
{\smallskip\smallskip \begin {scriptsize} Затверджено на засіданні кафедри теоретичної кібернетики, протокол \No 3  від   08.11.2025 р.\par\smallskip\smallskip\smallskip\smallskip
{\parbox {6.5 cm} {Зав. кафедри \dotfill Ю.В.~Крак}}\hspace {0.7cm} {\hbox to 6.5 cm { Екзаменатор \dotfill Т.О.~Карнаух}} \end{scriptsize}}
%end
\newpage
%begin
{{\begin {scriptsize}\par
{ \center  Київський національний університет імені Тараса Шевченка \par}
{\parbox {5 cm}{Освітня програма \hfill}{\parbox {7 cm}{Інформатика\hfill}}\parbox {6 cm}{\hfill Семестр 1}}\par
{\parbox {5 cm} {Навчальний предмет \hfill}\parbox {7 cm}{Програмування \hfill}\parbox {6cm}{\hfill}\par}\vspace{-15pt} \end{scriptsize}}{\center Екзаменаційний білет \No \textbf { 107 }\par}}
{%beginitem
1. Які з наведених типів не є цілими (integral), якщо такі є серед перелічених?

char, double, int, float, unsigned*?
%enditem
\par
\vspace{2mm}
%beginitem
2. Скільки змінних, що є вказівниками, тут оголошено?
long pointer1, pointer2, pointer3, pointer4;?
%enditem
\par
\vspace{2mm}
%beginitem
3. За наявних оголошень \code{void f(double\&); double x;} які з виразів \code{f(x*x), f(x+=2), f(\&x)} є синтаксично коректними?
%enditem
\par
\vspace{2mm}
%beginitem
4. Що таке С-рядок?
%enditem
\par
\vspace{2mm}
%beginitem
5. Яку часову оцінку складності має алгоритм швидкого сортування?
%enditem
\par
\vspace{2mm}
%beginitem
6. Написати програму, що запитує у користувача значення дійсних параметрів $x$ та $y$, обчислює для них значення виразу та повiдомляє введенi данi та результат обчислення.
$$\frac{\sqrt x - 15}{(\arctg x - 0.4) (y + 48)}$$ 

%enditem
\par
%beginitem
7. Написати функцію, яка для цілого числа знаходить кількість його різних простих дільників, що мають у десятковому запису цифру 1.

По завданням 2 та 3 оцінюється правильність та якість коду, а також економність обчислень.

%enditem
}
{\smallskip\smallskip \begin {scriptsize} Затверджено на засіданні кафедри теоретичної кібернетики, протокол \No 3  від   08.11.2025 р.\par\smallskip\smallskip\smallskip\smallskip
{\parbox {6.5 cm} {Зав. кафедри \dotfill Ю.В.~Крак}}\hspace {0.7cm} {\hbox to 6.5 cm { Екзаменатор \dotfill Т.О.~Карнаух}} \end{scriptsize}}
%end
\vspace{0.9cm}
%begin
{{\begin {scriptsize}\par
{ \center  Київський національний університет імені Тараса Шевченка \par}
{\parbox {5 cm}{Освітня програма \hfill}{\parbox {7 cm}{Інформатика\hfill}}\parbox {6 cm}{\hfill Семестр 1}}\par
{\parbox {5 cm} {Навчальний предмет \hfill}\parbox {7 cm}{Програмування \hfill}\parbox {6cm}{\hfill}\par}\vspace{-15pt} \end{scriptsize}}{\center Екзаменаційний білет \No \textbf { 108 }\par}}
{%beginitem
1. Які з наведених типів не є цілими (integral), якщо такі є серед перелічених?

char, double, int, float, unsigned?
%enditem
\par
\vspace{2mm}
%beginitem
2. Скільки змінних, що є вказівниками, тут оголошено?
char *s1, pc, **s2;?
%enditem
\par
\vspace{2mm}
%beginitem
3. За наявних оголошень \code{void f(char*); char c;} які з виразів \code{f(c*c), f(c+=2), f(\&c)} є синтаксично коректними?
%enditem
\par
\vspace{2mm}
%beginitem
4. У чому відмінність літералів від констант?
%enditem
\par
\vspace{2mm}
%beginitem
5. Яку часову оцінку складності має алгоритм сортування купою (він же деревом, пірамідою)?
%enditem
\par
\vspace{2mm}
%beginitem
6. Написати програму, що запитує у користувача значення дійсних параметрів $x$ та $y$, обчислює для них значення виразу та повiдомляє введенi данi та результат обчислення.
$$\frac{\pi x + 58}{(\sin x + 0.5) (y - 14)}$$ 

%enditem
\par
%beginitem
7. Написати функцію, що для інтервалу $[a, b)$ знаходить найближчу пару чисел з цього інтервалу, що кожне з них має парну кількість різних простих дільників. Якщо таких пар кілька, то цікавить пара з найбільшою сумою.

По завданням 2 та 3 оцінюється правильність та якість коду, а також економність обчислень.

%enditem
}
{\smallskip\smallskip \begin {scriptsize} Затверджено на засіданні кафедри теоретичної кібернетики, протокол \No 3  від   08.11.2025 р.\par\smallskip\smallskip\smallskip\smallskip
{\parbox {6.5 cm} {Зав. кафедри \dotfill Ю.В.~Крак}}\hspace {0.7cm} {\hbox to 6.5 cm { Екзаменатор \dotfill Т.О.~Карнаух}} \end{scriptsize}}
%end
\newpage
%begin
{{\begin {scriptsize}\par
{ \center  Київський національний університет імені Тараса Шевченка \par}
{\parbox {5 cm}{Освітня програма \hfill}{\parbox {7 cm}{Інформатика\hfill}}\parbox {6 cm}{\hfill Семестр 1}}\par
{\parbox {5 cm} {Навчальний предмет \hfill}\parbox {7 cm}{Програмування \hfill}\parbox {6cm}{\hfill}\par}\vspace{-15pt} \end{scriptsize}}{\center Екзаменаційний білет \No \textbf { 109 }\par}}
{%beginitem
1. Які з наведених типів не є цілими (integral), якщо такі є серед перелічених?

bool, int, long, long long, unsigned char?
%enditem
\par
\vspace{2mm}
%beginitem
2. Скільки змінних, що є вказівниками, тут оголошено?
int *pn1, pn2, pn3, pn4;?
%enditem
\par
\vspace{2mm}
%beginitem
3. За наявних оголошень \code{void f(int\&); int n;} які з виразів \code{f(13), f(++n), f(\&n)} є синтаксично коректними?
%enditem
\par
\vspace{2mm}
%beginitem
4. Що таке просторова складність алгоритму?
%enditem
\par
\vspace{2mm}
%beginitem
5. Яку часову оцінку складності має алгоритм сортування бульбашкою?
%enditem
\par
\vspace{2mm}
%beginitem
6. Написати програму, що запитує у користувача значення дійсних параметрів $x$ та $y$, обчислює для них значення виразу та повiдомляє введенi данi та результат обчислення.
$$\frac{\pi x - 26}{(\arccos x + 0.4) (y - 83)}$$ 

%enditem
\par
%beginitem
7. Написати функцію, що для цілого числа знаходить простий дільник, який входить у його розклад у найбільшому степені. Якщо таких кілька, то повернути найменший. Наприклад, для числа $2^{9}{\cdot}3^5{\cdot}7^9{\cdot}11$ функція має повернути $2$.

По завданням 2 та 3 оцінюється правильність та якість коду, а також економність обчислень.

%enditem
}
{\smallskip\smallskip \begin {scriptsize} Затверджено на засіданні кафедри теоретичної кібернетики, протокол \No 3  від   08.11.2025 р.\par\smallskip\smallskip\smallskip\smallskip
{\parbox {6.5 cm} {Зав. кафедри \dotfill Ю.В.~Крак}}\hspace {0.7cm} {\hbox to 6.5 cm { Екзаменатор \dotfill Т.О.~Карнаух}} \end{scriptsize}}
%end
\vspace{0.9cm}
%begin
{{\begin {scriptsize}\par
{ \center  Київський національний університет імені Тараса Шевченка \par}
{\parbox {5 cm}{Освітня програма \hfill}{\parbox {7 cm}{Інформатика\hfill}}\parbox {6 cm}{\hfill Семестр 1}}\par
{\parbox {5 cm} {Навчальний предмет \hfill}\parbox {7 cm}{Програмування \hfill}\parbox {6cm}{\hfill}\par}\vspace{-15pt} \end{scriptsize}}{\center Екзаменаційний білет \No \textbf { 110 }\par}}
{%beginitem
1. Які з наведених типів не є цілими (integral), якщо такі є серед перелічених?

double, int, long, unsigned char?
%enditem
\par
\vspace{2mm}
%beginitem
2. Скільки змінних, що є вказівниками, тут оголошено?
long *n1, *n2, n3, pn4;?
%enditem
\par
\vspace{2mm}
%beginitem
3. За наявних оголошень \code{void f(double*); double x;} які з виразів \code{f(x*x), f(x+=2), f(\&x)} є синтаксично коректними?
%enditem
\par
\vspace{2mm}
%beginitem
4. Що таке часова складність задачі?
%enditem
\par
\vspace{2mm}
%beginitem
5. Яку часову оцінку складності має алгоритм швидкого сортування?
%enditem
\par
\vspace{2mm}
%beginitem
6. Написати програму, що запитує у користувача значення дійсних параметрів $x$ та $y$, обчислює для них значення виразу та повiдомляє введенi данi та результат обчислення.
$$\frac{\sqrt x + 35}{(\cos x - 0.8) (y + 53)}$$ 

%enditem
\par
%beginitem
7. Написати функцію, що для інтервалу $[a, b)$ знаходить кількість чисел, квадрати яких діляться на суму квадратів їх цифр. (Розглядаємо десятковий запис.)

По завданням 2 та 3 оцінюється правильність та якість коду, а також економність обчислень.

%enditem
}
{\smallskip\smallskip \begin {scriptsize} Затверджено на засіданні кафедри теоретичної кібернетики, протокол \No 3  від   08.11.2025 р.\par\smallskip\smallskip\smallskip\smallskip
{\parbox {6.5 cm} {Зав. кафедри \dotfill Ю.В.~Крак}}\hspace {0.7cm} {\hbox to 6.5 cm { Екзаменатор \dotfill Т.О.~Карнаух}} \end{scriptsize}}
%end
\newpage
%begin
{{\begin {scriptsize}\par
{ \center  Київський національний університет імені Тараса Шевченка \par}
{\parbox {5 cm}{Освітня програма \hfill}{\parbox {7 cm}{Інформатика\hfill}}\parbox {6 cm}{\hfill Семестр 1}}\par
{\parbox {5 cm} {Навчальний предмет \hfill}\parbox {7 cm}{Програмування \hfill}\parbox {6cm}{\hfill}\par}\vspace{-15pt} \end{scriptsize}}{\center Екзаменаційний білет \No \textbf { 111 }\par}}
{%beginitem
1. Які з наведених типів не є цілими (integral), якщо такі є серед перелічених?

char, double, int, float, unsigned?
%enditem
\par
\vspace{2mm}
%beginitem
2. Скільки змінних, що є вказівниками, тут оголошено?
int n, m, *p1, p2;?
%enditem
\par
\vspace{2mm}
%beginitem
3. За наявних оголошень \code{void f(double); double x;} які з виразів \code{f(3.1), f(x+=2), f(\&x)} є синтаксично коректними?
%enditem
\par
\vspace{2mm}
%beginitem
4. Що таке часова складність алгоритму?
%enditem
\par
\vspace{2mm}
%beginitem
5. Яку часову оцінку складності має алгоритм сортування вставками?
%enditem
\par
\vspace{2mm}
%beginitem
6. Написати програму, що запитує у користувача значення дійсних параметрів $x$ та $y$, обчислює для них значення виразу та повiдомляє введенi данi та результат обчислення.
$$\frac{\ x - 41}{(\tg x + 0.5) (y - 12)}$$ 

%enditem
\par
%beginitem
7. Написати функцію, що для цілого числа знаходить кількість його простих дільників, що входять у його розклад рівно у другому степеню. Наприклад, для числа $3^2{\cdot}7^9{\cdot}11{\cdot}23^{2}$ таких дільників два – $3$ та $23$, тому функція має повернути $2$.

По завданням 2 та 3 оцінюється правильність та якість коду, а також економність обчислень.

%enditem
}
{\smallskip\smallskip \begin {scriptsize} Затверджено на засіданні кафедри теоретичної кібернетики, протокол \No 3  від   08.11.2025 р.\par\smallskip\smallskip\smallskip\smallskip
{\parbox {6.5 cm} {Зав. кафедри \dotfill Ю.В.~Крак}}\hspace {0.7cm} {\hbox to 6.5 cm { Екзаменатор \dotfill Т.О.~Карнаух}} \end{scriptsize}}
%end
\vspace{0.9cm}
%begin
{{\begin {scriptsize}\par
{ \center  Київський національний університет імені Тараса Шевченка \par}
{\parbox {5 cm}{Освітня програма \hfill}{\parbox {7 cm}{Інформатика\hfill}}\parbox {6 cm}{\hfill Семестр 1}}\par
{\parbox {5 cm} {Навчальний предмет \hfill}\parbox {7 cm}{Програмування \hfill}\parbox {6cm}{\hfill}\par}\vspace{-15pt} \end{scriptsize}}{\center Екзаменаційний білет \No \textbf { 112 }\par}}
{%beginitem
1. Які з наведених типів не є цілими (integral), якщо такі є серед перелічених?

bool, int, long, long long, unsigned char?
%enditem
\par
\vspace{2mm}
%beginitem
2. Скільки змінних, що є вказівниками, тут оголошено?
double *px, *py, **ppx;?
%enditem
\par
\vspace{2mm}
%beginitem
3. За наявних оголошень \code{void f(int*); int n;} які з виразів \code{f(n++), f(++n), f(\&n)} є синтаксично коректними?
%enditem
\par
\vspace{2mm}
%beginitem
4. Що таке часова складність задачі?
%enditem
\par
\vspace{2mm}
%beginitem
5. Яку часову оцінку складності має алгоритм сортування злиттям?
%enditem
\par
\vspace{2mm}
%beginitem
6. Написати програму, що запитує у користувача значення дійсних параметрів $x$ та $y$, обчислює для них значення виразу та повiдомляє введенi данi та результат обчислення.
$$\frac{\ln x + 94}{(\ctg x - 0.2) (y + 72)}$$ 

%enditem
\par
%beginitem
7. Написати функцію, що повертає найбільше число з послідовності цілих (задається масивом), у сімковому запису якого нема цифри 2. За відсутності має повертати None. Використовувати додаткові структури даних (у тому числі рядки) заборонено.

По завданням 2 та 3 оцінюється правильність та якість коду, а також економність обчислень.

%enditem
}
{\smallskip\smallskip \begin {scriptsize} Затверджено на засіданні кафедри теоретичної кібернетики, протокол \No 3  від   08.11.2025 р.\par\smallskip\smallskip\smallskip\smallskip
{\parbox {6.5 cm} {Зав. кафедри \dotfill Ю.В.~Крак}}\hspace {0.7cm} {\hbox to 6.5 cm { Екзаменатор \dotfill Т.О.~Карнаух}} \end{scriptsize}}
%end
\newpage
%begin
{{\begin {scriptsize}\par
{ \center  Київський національний університет імені Тараса Шевченка \par}
{\parbox {5 cm}{Освітня програма \hfill}{\parbox {7 cm}{Інформатика\hfill}}\parbox {6 cm}{\hfill Семестр 1}}\par
{\parbox {5 cm} {Навчальний предмет \hfill}\parbox {7 cm}{Програмування \hfill}\parbox {6cm}{\hfill}\par}\vspace{-15pt} \end{scriptsize}}{\center Екзаменаційний білет \No \textbf { 113 }\par}}
{%beginitem
1. Які з наведених типів не є цілими (integral), якщо такі є серед перелічених?

bool*, char, int, unsigned?
%enditem
\par
\vspace{2mm}
%beginitem
2. Скільки змінних, що є вказівниками, тут оголошено?
long pointer1, pointer2, pointer3, pointer4;?
%enditem
\par
\vspace{2mm}
%beginitem
3. За наявних оголошень \code{void f(int); int n;} які з виразів \code{f(3), f(++n), f(\&n)} є синтаксично коректними?
%enditem
\par
\vspace{2mm}
%beginitem
4. Чим префіксний оператор ++ відрізняється від постфіксного?
%enditem
\par
\vspace{2mm}
%beginitem
5. Яку часову оцінку складності має алгоритм сортування вибором?
%enditem
\par
\vspace{2mm}
%beginitem
6. Написати програму, що запитує у користувача значення дійсних параметрів $x$ та $y$, обчислює для них значення виразу та повiдомляє введенi данi та результат обчислення.
$$\frac{\log_{10} x + 85}{(\arccos x - 0.1) (y + 33)}$$ 

%enditem
\par
%beginitem
7. Написати функцію, що для інтервалу $[a, b)$ знаходить найближчу пару чисел з цього інтервалу, що мають не менше трьох простих дільників та мають цифру 3 у своєму десятковому запису.

По завданням 2 та 3 оцінюється правильність та якість коду, а також економність обчислень.

%enditem
}
{\smallskip\smallskip \begin {scriptsize} Затверджено на засіданні кафедри теоретичної кібернетики, протокол \No 3  від   08.11.2025 р.\par\smallskip\smallskip\smallskip\smallskip
{\parbox {6.5 cm} {Зав. кафедри \dotfill Ю.В.~Крак}}\hspace {0.7cm} {\hbox to 6.5 cm { Екзаменатор \dotfill Т.О.~Карнаух}} \end{scriptsize}}
%end
\vspace{0.9cm}
%begin
{{\begin {scriptsize}\par
{ \center  Київський національний університет імені Тараса Шевченка \par}
{\parbox {5 cm}{Освітня програма \hfill}{\parbox {7 cm}{Інформатика\hfill}}\parbox {6 cm}{\hfill Семестр 1}}\par
{\parbox {5 cm} {Навчальний предмет \hfill}\parbox {7 cm}{Програмування \hfill}\parbox {6cm}{\hfill}\par}\vspace{-15pt} \end{scriptsize}}{\center Екзаменаційний білет \No \textbf { 114 }\par}}
{%beginitem
1. Які з наведених типів не є цілими (integral), якщо такі є серед перелічених?

double, int, long*, unsigned char?
%enditem
\par
\vspace{2mm}
%beginitem
2. Скільки змінних, що є вказівниками, тут оголошено?
int *pn1, pn2, pn3, pn4;?
%enditem
\par
\vspace{2mm}
%beginitem
3. За наявних оголошень \code{void f(double\&); double x;} які з виразів \code{f(x*x), f(x+=2), f(\&x)} є синтаксично коректними?
%enditem
\par
\vspace{2mm}
%beginitem
4. Що таке просторова складність задачі?
%enditem
\par
\vspace{2mm}
%beginitem
5. Яку часову оцінку складності має алгоритм сортування купою (він же деревом, пірамідою)?
%enditem
\par
\vspace{2mm}
%beginitem
6. Написати програму, що запитує у користувача значення дійсних параметрів $x$ та $y$, обчислює для них значення виразу та повiдомляє введенi данi та результат обчислення.
$$\frac{\log_{2} x - 88}{(\cos x + 0.7) (y - 17)}$$ 

%enditem
\par
%beginitem
7. Написати функцію, що в інтервалі $[a, b)$ знаходить ціле число, що ділиться на найбільший степінь числа 3. Якщо таких чисел декілька, повернути найменше.

По завданням 2 та 3 оцінюється правильність та якість коду, а також економність обчислень.

%enditem
}
{\smallskip\smallskip \begin {scriptsize} Затверджено на засіданні кафедри теоретичної кібернетики, протокол \No 3  від   08.11.2025 р.\par\smallskip\smallskip\smallskip\smallskip
{\parbox {6.5 cm} {Зав. кафедри \dotfill Ю.В.~Крак}}\hspace {0.7cm} {\hbox to 6.5 cm { Екзаменатор \dotfill Т.О.~Карнаух}} \end{scriptsize}}
%end
\newpage
%begin
{{\begin {scriptsize}\par
{ \center  Київський національний університет імені Тараса Шевченка \par}
{\parbox {5 cm}{Освітня програма \hfill}{\parbox {7 cm}{Інформатика\hfill}}\parbox {6 cm}{\hfill Семестр 1}}\par
{\parbox {5 cm} {Навчальний предмет \hfill}\parbox {7 cm}{Програмування \hfill}\parbox {6cm}{\hfill}\par}\vspace{-15pt} \end{scriptsize}}{\center Екзаменаційний білет \No \textbf { 115 }\par}}
{%beginitem
1. Які з наведених типів не є цілими (integral), якщо такі є серед перелічених?

char, double, int, float, unsigned*?
%enditem
\par
\vspace{2mm}
%beginitem
2. Скільки змінних, що є вказівниками, тут оголошено?
int n, m, *p1, p2;?
%enditem
\par
\vspace{2mm}
%beginitem
3. За наявних оголошень \code{void f(char); char c;} які з виразів \code{f(c*c), f(c+=2), f(\&c)} є синтаксично коректними?
%enditem
\par
\vspace{2mm}
%beginitem
4. Що таке масив?
%enditem
\par
\vspace{2mm}
%beginitem
5. Яку часову оцінку складності має алгоритм сортування злиттям?
%enditem
\par
\vspace{2mm}
%beginitem
6. Написати програму, що запитує у користувача значення дійсних параметрів $x$ та $y$, обчислює для них значення виразу та повiдомляє введенi данi та результат обчислення.
$$\frac{\pi x - 50}{(\ctg x - 0.9) (y - 45)}$$ 

%enditem
\par
%beginitem
7. Написати функцію, що для послідовності цілих (задається масивом) знаходить найдовший відрізок, в якому всі сусідні числа відрізняються якнайменш на 3. Якщо таких відрізків кілька, то повернути перший (як індекс початку та довжину). Використовувати додаткові структури даних (у тому числі рядки) заборонено.

По завданням 2 та 3 оцінюється правильність та якість коду, а також економність обчислень.

%enditem
}
{\smallskip\smallskip \begin {scriptsize} Затверджено на засіданні кафедри теоретичної кібернетики, протокол \No 3  від   08.11.2025 р.\par\smallskip\smallskip\smallskip\smallskip
{\parbox {6.5 cm} {Зав. кафедри \dotfill Ю.В.~Крак}}\hspace {0.7cm} {\hbox to 6.5 cm { Екзаменатор \dotfill Т.О.~Карнаух}} \end{scriptsize}}
%end
\vspace{0.9cm}
%begin
{{\begin {scriptsize}\par
{ \center  Київський національний університет імені Тараса Шевченка \par}
{\parbox {5 cm}{Освітня програма \hfill}{\parbox {7 cm}{Інформатика\hfill}}\parbox {6 cm}{\hfill Семестр 1}}\par
{\parbox {5 cm} {Навчальний предмет \hfill}\parbox {7 cm}{Програмування \hfill}\parbox {6cm}{\hfill}\par}\vspace{-15pt} \end{scriptsize}}{\center Екзаменаційний білет \No \textbf { 116 }\par}}
{%beginitem
1. Які з наведених типів не є цілими (integral), якщо такі є серед перелічених?

bool, char, int, unsigned?
%enditem
\par
\vspace{2mm}
%beginitem
2. Скільки змінних, що є вказівниками, тут оголошено?
long *n1, *n2, n3, pn4;?
%enditem
\par
\vspace{2mm}
%beginitem
3. За наявних оголошень \code{void f(char\&); char c;} які з виразів \code{f(c*c), f(c++), f(\&c)} є синтаксично коректними?
%enditem
\par
\vspace{2mm}
%beginitem
4. Що означає термін ``часова складність у середньому''?
%enditem
\par
\vspace{2mm}
%beginitem
5. Яку часову оцінку складності має алгоритм швидкого сортування?
%enditem
\par
\vspace{2mm}
%beginitem
6. Написати програму, що запитує у користувача значення дійсних параметрів $x$ та $y$, обчислює для них значення виразу та повiдомляє введенi данi та результат обчислення.
$$\frac{\log_{2} x + 67}{(\tg x + 0.3) (y + 73)}$$ 

%enditem
\par
%beginitem
7. Написати функцію, що для цілого числа знаходить кількість його простих дільників, що входять у його розклад не більш ніж у 5му степеню. Наприклад, для числа $3^8{\cdot}7^9{\cdot}11{\cdot}23^{2}$ таких дільників два – $11$ та $23$, тому функція має повернути $2$.

По завданням 2 та 3 оцінюється правильність та якість коду, а також економність обчислень.

%enditem
}
{\smallskip\smallskip \begin {scriptsize} Затверджено на засіданні кафедри теоретичної кібернетики, протокол \No 3  від   08.11.2025 р.\par\smallskip\smallskip\smallskip\smallskip
{\parbox {6.5 cm} {Зав. кафедри \dotfill Ю.В.~Крак}}\hspace {0.7cm} {\hbox to 6.5 cm { Екзаменатор \dotfill Т.О.~Карнаух}} \end{scriptsize}}
%end
\newpage
%begin
{{\begin {scriptsize}\par
{ \center  Київський національний університет імені Тараса Шевченка \par}
{\parbox {5 cm}{Освітня програма \hfill}{\parbox {7 cm}{Інформатика\hfill}}\parbox {6 cm}{\hfill Семестр 1}}\par
{\parbox {5 cm} {Навчальний предмет \hfill}\parbox {7 cm}{Програмування \hfill}\parbox {6cm}{\hfill}\par}\vspace{-15pt} \end{scriptsize}}{\center Екзаменаційний білет \No \textbf { 117 }\par}}
{%beginitem
1. Які з наведених типів не є цілими (integral), якщо такі є серед перелічених?

double, int, long, unsigned char?
%enditem
\par
\vspace{2mm}
%beginitem
2. Скільки змінних, що є вказівниками, тут оголошено?
double *px, *py, **ppx;?
%enditem
\par
\vspace{2mm}
%beginitem
3. За наявних оголошень \code{void f(char*); char c;} які з виразів \code{f(c*c), f(c+=2), f(\&c)} є синтаксично коректними?
%enditem
\par
\vspace{2mm}
%beginitem
4. Що таке просторова складність алгоритму?
%enditem
\par
\vspace{2mm}
%beginitem
5. Яку часову оцінку складності має алгоритм сортування бульбашкою?
%enditem
\par
\vspace{2mm}
%beginitem
6. Написати програму, що запитує у користувача значення дійсних параметрів $x$ та $y$, обчислює для них значення виразу та повiдомляє введенi данi та результат обчислення.
$$\frac{\log_{10} x + 91}{(\arcsin x + 0.6) (y + 63)}$$ 

%enditem
\par
%beginitem
7. Написати функцію, що в інтервалі $[a, b)$ знаходить ціле число, що ділиться на найбільший степінь числа 3. Якщо таких чисел декілька, повернути найбільше.

По завданням 2 та 3 оцінюється правильність та якість коду, а також економність обчислень.

%enditem
}
{\smallskip\smallskip \begin {scriptsize} Затверджено на засіданні кафедри теоретичної кібернетики, протокол \No 3  від   08.11.2025 р.\par\smallskip\smallskip\smallskip\smallskip
{\parbox {6.5 cm} {Зав. кафедри \dotfill Ю.В.~Крак}}\hspace {0.7cm} {\hbox to 6.5 cm { Екзаменатор \dotfill Т.О.~Карнаух}} \end{scriptsize}}
%end
\vspace{0.9cm}
%begin
{{\begin {scriptsize}\par
{ \center  Київський національний університет імені Тараса Шевченка \par}
{\parbox {5 cm}{Освітня програма \hfill}{\parbox {7 cm}{Інформатика\hfill}}\parbox {6 cm}{\hfill Семестр 1}}\par
{\parbox {5 cm} {Навчальний предмет \hfill}\parbox {7 cm}{Програмування \hfill}\parbox {6cm}{\hfill}\par}\vspace{-15pt} \end{scriptsize}}{\center Екзаменаційний білет \No \textbf { 118 }\par}}
{%beginitem
1. Які з наведених типів не є цілими (integral), якщо такі є серед перелічених?

bool, int, long, long long, unsigned char?
%enditem
\par
\vspace{2mm}
%beginitem
2. Скільки змінних, що є вказівниками, тут оголошено?
char *s1, pc, **s2;?
%enditem
\par
\vspace{2mm}
%beginitem
3. За наявних оголошень \code{void f(int\&); int n;} які з виразів \code{f(13), f(++n), f(\&n)} є синтаксично коректними?
%enditem
\par
\vspace{2mm}
%beginitem
4. У чому відмінність літералів від констант?
%enditem
\par
\vspace{2mm}
%beginitem
5. Яку часову оцінку складності має алгоритм швидкого сортування?
%enditem
\par
\vspace{2mm}
%beginitem
6. Написати програму, що запитує у користувача значення дійсних параметрів $x$ та $y$, обчислює для них значення виразу та повiдомляє введенi данi та результат обчислення.
$$\frac{\ x - 95}{(\arctg x - 0.1) (y - 97)}$$ 

%enditem
\par
%beginitem
7. Написати функцію, що знаходить найбільше число з послідовності цілих (задається масивом), у десятковому запису якого нема цифри 4. Повертає індекс першого входження цього числа або $-1$ (за відсутності). Використовувати додаткові структури даних (у тому числі рядки) заборонено.

По завданням 2 та 3 оцінюється правильність та якість коду, а також економність обчислень.

%enditem
}
{\smallskip\smallskip \begin {scriptsize} Затверджено на засіданні кафедри теоретичної кібернетики, протокол \No 3  від   08.11.2025 р.\par\smallskip\smallskip\smallskip\smallskip
{\parbox {6.5 cm} {Зав. кафедри \dotfill Ю.В.~Крак}}\hspace {0.7cm} {\hbox to 6.5 cm { Екзаменатор \dotfill Т.О.~Карнаух}} \end{scriptsize}}
%end
\newpage
%begin
{{\begin {scriptsize}\par
{ \center  Київський національний університет імені Тараса Шевченка \par}
{\parbox {5 cm}{Освітня програма \hfill}{\parbox {7 cm}{Інформатика\hfill}}\parbox {6 cm}{\hfill Семестр 1}}\par
{\parbox {5 cm} {Навчальний предмет \hfill}\parbox {7 cm}{Програмування \hfill}\parbox {6cm}{\hfill}\par}\vspace{-15pt} \end{scriptsize}}{\center Екзаменаційний білет \No \textbf { 119 }\par}}
{%beginitem
1. Які з наведених типів не є цілими (integral), якщо такі є серед перелічених?

int*, signed char, unsigned long, string?
%enditem
\par
\vspace{2mm}
%beginitem
2. Скільки змінних, що є вказівниками, тут оголошено?
double x, *px1, px2, ***px3;?
%enditem
\par
\vspace{2mm}
%beginitem
3. За наявних оголошень \code{void f(double*); double x;} які з виразів \code{f(x*x), f(x+=2), f(\&x)} є синтаксично коректними?
%enditem
\par
\vspace{2mm}
%beginitem
4. Які параметри може приймати функція main? Що вони позначають та яких вони типів?
%enditem
\par
\vspace{2mm}
%beginitem
5. Яку часову оцінку складності має алгоритм сортування вставками?
%enditem
\par
\vspace{2mm}
%beginitem
6. Написати програму, що запитує у користувача значення дійсних параметрів $x$ та $y$, обчислює для них значення виразу та повiдомляє введенi данi та результат обчислення.
$$\frac{\ln x + 13}{(\sin x + 0.2) (y - 52)}$$ 

%enditem
\par
%beginitem
7. Написати функцію, що повертає число з послідовності цілих (задається масивом), що має найменшу кількість різних простих дільників. Якщо таких чисел кілька, то повернути найбільше. Використовувати додаткові структури даних (у тому числі рядки) заборонено.

По завданням 2 та 3 оцінюється правильність та якість коду, а також економність обчислень.

%enditem
}
{\smallskip\smallskip \begin {scriptsize} Затверджено на засіданні кафедри теоретичної кібернетики, протокол \No 3  від   08.11.2025 р.\par\smallskip\smallskip\smallskip\smallskip
{\parbox {6.5 cm} {Зав. кафедри \dotfill Ю.В.~Крак}}\hspace {0.7cm} {\hbox to 6.5 cm { Екзаменатор \dotfill Т.О.~Карнаух}} \end{scriptsize}}
%end
\vspace{0.9cm}
%begin
{{\begin {scriptsize}\par
{ \center  Київський національний університет імені Тараса Шевченка \par}
{\parbox {5 cm}{Освітня програма \hfill}{\parbox {7 cm}{Інформатика\hfill}}\parbox {6 cm}{\hfill Семестр 1}}\par
{\parbox {5 cm} {Навчальний предмет \hfill}\parbox {7 cm}{Програмування \hfill}\parbox {6cm}{\hfill}\par}\vspace{-15pt} \end{scriptsize}}{\center Екзаменаційний білет \No \textbf { 120 }\par}}
{%beginitem
1. Які з наведених типів не є цілими (integral), якщо такі є серед перелічених?

int, signed char, unsigned long, string?
%enditem
\par
\vspace{2mm}
%beginitem
2. Скільки змінних, що є вказівниками, тут оголошено?
double x, *px1, px2, ***px3;?
%enditem
\par
\vspace{2mm}
%beginitem
3. За наявних оголошень \code{void f(char\&); char c;} які з виразів \code{f(c*c), f(c++), f(\&c)} є синтаксично коректними?
%enditem
\par
\vspace{2mm}
%beginitem
4. Що таке тип даних?
%enditem
\par
\vspace{2mm}
%beginitem
5. Яку часову оцінку складності має алгоритм швидкого сортування?
%enditem
\par
\vspace{2mm}
%beginitem
6. Написати програму, що запитує у користувача значення дійсних параметрів $x$ та $y$, обчислює для них значення виразу та повiдомляє введенi данi та результат обчислення.
$$\frac{\sqrt x - 36}{(\arcsin x - 0.5) (y + 93)}$$ 

%enditem
\par
%beginitem
7. Написати функцію, що повертає число з послідовності цілих (задається масивом), що має найбільшу кількість різних простих дільників. Якщо таких чисел кілька, то повернути найменше. Використовувати додаткові структури даних (у тому числі рядки) заборонено.

По завданням 2 та 3 оцінюється правильність та якість коду, а також економність обчислень.

%enditem
}
{\smallskip\smallskip \begin {scriptsize} Затверджено на засіданні кафедри теоретичної кібернетики, протокол \No 3  від   08.11.2025 р.\par\smallskip\smallskip\smallskip\smallskip
{\parbox {6.5 cm} {Зав. кафедри \dotfill Ю.В.~Крак}}\hspace {0.7cm} {\hbox to 6.5 cm { Екзаменатор \dotfill Т.О.~Карнаух}} \end{scriptsize}}
%end
\newpage
%begin
{{\begin {scriptsize}\par
{ \center  Київський національний університет імені Тараса Шевченка \par}
{\parbox {5 cm}{Освітня програма \hfill}{\parbox {7 cm}{Інформатика\hfill}}\parbox {6 cm}{\hfill Семестр 1}}\par
{\parbox {5 cm} {Навчальний предмет \hfill}\parbox {7 cm}{Програмування \hfill}\parbox {6cm}{\hfill}\par}\vspace{-15pt} \end{scriptsize}}{\center Екзаменаційний білет \No \textbf { 121 }\par}}
{%beginitem
1. Які з наведених типів не є цілими (integral), якщо такі є серед перелічених?

bool, int, long, long long, unsigned char?
%enditem
\par
\vspace{2mm}
%beginitem
2. Скільки змінних, що є вказівниками, тут оголошено?
int n, m, *p1, p2;?
%enditem
\par
\vspace{2mm}
%beginitem
3. За наявних оголошень \code{void f(double\&); double x;} які з виразів \code{f(x*x), f(x+=2), f(\&x)} є синтаксично коректними?
%enditem
\par
\vspace{2mm}
%beginitem
4. Що таке С-рядок?
%enditem
\par
\vspace{2mm}
%beginitem
5. Яку часову оцінку складності має алгоритм сортування бульбашкою?
%enditem
\par
\vspace{2mm}
%beginitem
6. Написати програму, що запитує у користувача значення дійсних параметрів $x$ та $y$, обчислює для них значення виразу та повiдомляє введенi данi та результат обчислення.
$$\frac{\log_{10} x - 42}{(\ctg x - 0.4) (y + 29)}$$ 

%enditem
\par
%beginitem
7. Написати функцію, що для цілого числа знаходить простий дільник, який входить у його розклад ц найбільшому степені. Якщо таких кілька, то повернути найбільший. Наприклад, для числа $2^{9}{\cdot}3^5{\cdot}7^9{\cdot}11$ функція має повернути $7$.

По завданням 2 та 3 оцінюється правильність та якість коду, а також економність обчислень.

%enditem
}
{\smallskip\smallskip \begin {scriptsize} Затверджено на засіданні кафедри теоретичної кібернетики, протокол \No 3  від   08.11.2025 р.\par\smallskip\smallskip\smallskip\smallskip
{\parbox {6.5 cm} {Зав. кафедри \dotfill Ю.В.~Крак}}\hspace {0.7cm} {\hbox to 6.5 cm { Екзаменатор \dotfill Т.О.~Карнаух}} \end{scriptsize}}
%end
\vspace{0.9cm}
%begin
{{\begin {scriptsize}\par
{ \center  Київський національний університет імені Тараса Шевченка \par}
{\parbox {5 cm}{Освітня програма \hfill}{\parbox {7 cm}{Інформатика\hfill}}\parbox {6 cm}{\hfill Семестр 1}}\par
{\parbox {5 cm} {Навчальний предмет \hfill}\parbox {7 cm}{Програмування \hfill}\parbox {6cm}{\hfill}\par}\vspace{-15pt} \end{scriptsize}}{\center Екзаменаційний білет \No \textbf { 122 }\par}}
{%beginitem
1. Які з наведених типів не є цілими (integral), якщо такі є серед перелічених?

double, int, long, unsigned char?
%enditem
\par
\vspace{2mm}
%beginitem
2. Скільки змінних, що є вказівниками, тут оголошено?
char *s1, pc, **s2;?
%enditem
\par
\vspace{2mm}
%beginitem
3. За наявних оголошень \code{void f(int); int n;} які з виразів \code{f(3), f(++n), f(\&n)} є синтаксично коректними?
%enditem
\par
\vspace{2mm}
%beginitem
4. Які параметри може приймати функція main? Що вони позначають та яких вони типів?
%enditem
\par
\vspace{2mm}
%beginitem
5. Яку часову оцінку складності має алгоритм сортування вставками?
%enditem
\par
\vspace{2mm}
%beginitem
6. Написати програму, що запитує у користувача значення дійсних параметрів $x$ та $y$, обчислює для них значення виразу та повiдомляє введенi данi та результат обчислення.
$$\frac{\pi x + 68}{(\tg x + 0.7) (y - 47)}$$ 

%enditem
\par
%beginitem
7. Написати функцію, що для послідовності цілих (задається масивом) знаходить довжину найдовшого відрізку, що складається з чисел, у сімковому запису якого нема цифри 2. За відсутності має повертати $0$. Використовувати додаткові структури даних (у тому числі рядки) заборонено.

По завданням 2 та 3 оцінюється правильність та якість коду, а також економність обчислень.

%enditem
}
{\smallskip\smallskip \begin {scriptsize} Затверджено на засіданні кафедри теоретичної кібернетики, протокол \No 3  від   08.11.2025 р.\par\smallskip\smallskip\smallskip\smallskip
{\parbox {6.5 cm} {Зав. кафедри \dotfill Ю.В.~Крак}}\hspace {0.7cm} {\hbox to 6.5 cm { Екзаменатор \dotfill Т.О.~Карнаух}} \end{scriptsize}}
%end
\newpage
\end{document}
